\chapter{The trivially valued analytification}
\index{analytification!trivially valued}\index{semivaluation}
\label{chap:trivial-analytification}

The trivial absolute value does not make analytification trivial.  It makes
the radial variable visible.  A point of the resulting Berkovich space can
have a nonzero kernel, a valuation with full support can be rescaled by every
positive real number, and the fixed points of this rescaling recover the
irreducible closed subsets of $X$.  These three features distinguish the
theory in \Crefrange{chap:trivial-analytification}{chap:bj-stability} from
the discretely valued theory developed in
\Crefrange{chap:bfj-model-forms}{chap:bfj-examples}.

This chapter fixes the underlying topological space and the conventions used
in \Crefrange{chap:bj-test-configurations}{chap:bj-envelope-history}.
We first describe analytic points as semivaluations, then discuss support,
center, compactness, and scaling.  Divisorial and quasimonomial points provide
the computable core.  Finally, Gauss extension compares the trivially valued
space with the usual non-Archimedean generic fibre of a degeneration.

\section{Orientation}
\label{sec:ch32-orientation}

\subsection{What the trivial absolute value trivializes, and what it does not}
\label{sec:ch32-what-is-trivial}

The Berkovich spectrum of the trivially valued field $k$ is a single point.
It is tempting to conclude from this that the analytification of a
$k$-variety is uninteresting, and the conclusion is exactly wrong.  Three
structures become inseparable from the global theory, and each is used
continually in \Cref{part:bj}.

\emph{(i) Semivaluations with nonzero kernel are unavoidable.}
Over a complete nontrivially valued field $K$ one also allows points of
$X^{\An}$ whose associated seminorm has a kernel; but if $X$ is a variety
over $K$, the generic point of $X$ is already an analytic point and the
kernel phenomenon is a boundary effect that can, in the discretely valued
BFJ theory of \Cref{part:bfj}, largely be handled by working with model
functions and their limits.  Over a trivially valued field the situation is
different in kind.  We shall prove in
\Cref{prop:ch32-proper-support-dense} that if every irreducible component of
$X$ has dimension at least two, then the points with \emph{proper support}
--- those whose kernel is a nonzero prime --- are dense in $X^{\An}$.  They
cannot be discarded, and a psh function may take the value $-\infty$ at
such a point.  The standing convention of \Cref{part:bj}, announced in
\Cref{sec:ch32-standing}, is that the value $-\infty$ at a proper-support
semivaluation is a real feature of the theory, not an edge case to be
suppressed.

\emph{(ii) There is no distinguished slice.}
Over a discretely valued $K$ with uniformizer $\varpi$, every point of
$X^{\An}$ satisfies $v(\varpi)=1$; this single normalization is what makes
the discretely valued theory a theory of \emph{functions} on a compact space.
Over a trivially valued field there is no element of $k^\times$ with nonzero
value, and the rescaling $v\mapsto tv$ is an automorphism of the whole
situation.  It is a genuine $\mathbf R_{>0}$-action on $X^{\An}$
(\Cref{prop:ch32-scaling-action}), continuous and order preserving, whose
orbits are compactified rays.  Consequently the natural objects of
\Cref{part:bj} come in two flavours: homogeneous functions on the cone, and
bounded functions on the compact space.  \Cref{thm:ch32-no-compact-link}
shows that no positive finite homogeneous gauge can normalize all scaling
orbits of $X^{\An}$, because the trivial locus consists of fixed points.

\emph{(iii) The Zariski space sits inside the analytic space.}
The fixed points of the scaling action are exactly the trivial
semivaluations $v_{Y,\triv}$, one for each irreducible closed subvariety
$Y\subset X$ (\Cref{thm:ch32-trivial-locus}).  The underlying set of the
scheme $X$ is therefore a subset $X^{\triv}\subset X^{\An}$, and it is
even a closed subset; but the analytic topology restricted to it is the
compact Hausdorff topology induced from $X^{\An}$, not the Zariski topology.
The bijection $X\to X^{\triv}$ is a section of both the support map and the
center map, and its inverse is continuous while it is not.

The three structures do not transfer as a single unchanged package to a
nontrivially valued base: kernels still occur there, but the scaling action
and its fixed locus are special to the trivial norm.  Thus \Cref{part:bj} is
not a formal special case of \Cref{part:bfj}.  What relates the two is a
canonical \emph{section}
$\sigma$ of the restriction map $X_K^{\An}\to X^{\An}$, for
$K=k(\!(\varpi)\!)$, constructed in \Cref{sec:ch30-gauss} and recalled here
in \Cref{sec:ch32-gauss}.

\subsection{Standing hypotheses and notation}
\label{sec:ch32-standing}

\emph{The ground field.}  Throughout the chapter $k$ is an algebraically
closed field, of arbitrary characteristic unless the contrary is said,
equipped with the trivial absolute value; thus $v_k(a)=0$ for every
$a\in k^\times$ and $v_k(0)=+\infty$, in the sign convention
\eqref{eq:appJ-valuation-sign}.  It is complete for this absolute value,
so the entire Berkovich formalism of \Cref{chap:berkovich} applies verbatim;
see \Cref{sec:berk-trivial}, where the trivially valued case is treated as
a genuine case rather than as a degeneration.

\emph{Departure from the global convention on varieties.}  The global
convention of the front matter reserves the word \emph{variety} for a
separated integral scheme of finite type.  In this chapter, and in
\Crefrange{chap:trivial-analytification}{chap:bj-stability} generally, a
\emph{variety over $k$} means a separated reduced $k$-scheme of finite
type, not necessarily irreducible nor equidimensional.  This is the
convention of \cite{BJ22}, and it is forced on us: the analytification of a
reducible variety, the central fibre of a test configuration, and the
normalization of a nodal curve all occur below, and none of them is integral.
Whenever irreducibility or integrality is needed it is stated.  We write
$(X_\alpha)_\alpha$ for the irreducible components of $X$.

\emph{The variety.}  Unless stated otherwise, $X$ is a projective variety
over $k$ in the above sense, of dimension $n$.  Properness is used for
the existence of centers and for compactness; projectivity is used whenever
an ample line bundle is invoked, which happens in
\Cref{lem:ch32-sections-generate-local-ring},
\Cref{prop:ch32-valuation-criterion} and
\Cref{prop:ch32-divisorial-linear-growth}.  When valuations of the function
field $k(X)$ are used, $X$ is assumed integral.

\emph{Points of $X^{\An}$ are real semivaluations.}  A point of the
trivially valued analytification is treated \emph{consistently} as a real
semivaluation $v$, never as a seminorm, except where a multiplicative
formula is genuinely clearer.  We write $|f|(v)=\rme^{-v(f)}$ when the
multiplicative notation is wanted, with $\rme^{-\infty}=0$.  We do not
distinguish notationally between a point of $X^{\An}$ and the
semivaluation it is; thus $v\in X^{\An}$.

\emph{The trivial valuation and the proper-support points are retained.}
The trivial valuation $v_{X,\triv}$ of an integral $X$ is a point of
$X^{\An}$ and is counted as a divisorial point
(\Cref{def:ch32-divisorial}); the semivaluations with proper support are
points of $X^{\An}$ and are never removed.  Functions on $X^{\An}$ are
allowed the value $-\infty$, and the locus where this occurs is analysed
in \Cref{sec:ch32-proper-support}.  The two safeguards recorded in the
front matter are in force: values at proper-support points are retained,
and compactness is formulated with nets (\Cref{app:nets}).

\emph{Ideals.}  An \emph{ideal} on $X$ means a coherent ideal sheaf
$\mathfrak b\subset\cO_X$; we write $\mathcal I=\mathcal I(X)$ for the
set of these, $Z(\mathfrak b)\subset X$ for the closed subscheme it
defines, and $\mathfrak b=0$ for the zero ideal.  We allow
$\mathfrak b=\cO_X$ and $\mathfrak b=0$.

\emph{Uniformizer.}  In \Cref{sec:ch32-gauss} we put $K=k(\!(\varpi)\!)$
and $R=k[\![\varpi]\!]$ with $v_K(\varpi)=1$; the letter $\varpi$ is
also the coordinate on $\mathbf A^1_k$ in the test-configuration picture of
\Cref{chap:bj-test-configurations}.  This agrees with
\eqref{eq:gauss-extension} up to the change of name of the uniformizer noted
in \Cref{sec:ch30-hypotheses}.

\subsection{What is imported}
\label{sec:ch32-imported}

The chapter is largely self-contained; the following statements are used
from elsewhere and are the only external inputs.

\begin{enumerate}
\item\label{itm:ch32-import:berk}
The construction of the analytification and its basic properties,
\Cref{sec:berk-analytification} and \Cref{prop:berk-pair-description}, and
the compactness of the spectrum of a Banach ring,
\Cref{prop:berk-compact-spectrum}.  The properness theorem
\Cref{thm:berk-proper-compact} is quoted once, in
\Cref{thm:triv-berkovich-compactness}; \Cref{thm:ch32-center-compact} gives
an independent proof in the trivially valued case.

\item\label{itm:ch32-import:centers}
Existence and uniqueness of centers on a proper separated scheme,
\Cref{prop:centers-proper-separated}.

\item\label{itm:ch32-import:abhyankar}
The Abhyankar inequality \Cref{thm:abhyankar-inequality} and its consequence
\Cref{cor:rank-plus-trdeg}, together with the invariants of monomial
valuations, \Cref{lem:monomial-invariants}.

\item\label{itm:ch32-import:skeleta}
The construction of quasimonomial valuations from an SNC pair by graded
reduction, \Cref{lem:monomial-valuation}, and the skeleton formalism of
\Cref{chap:skeleta}.  These are stated there over a discretely valued base;
\Cref{sec:ch32-quasimonomial} explains the two changes needed over a
trivially valued base and repeats the argument in the form used here.

\item\label{itm:ch32-import:gauss}
The Gauss extension and everything proved about it in
\Cref{sec:ch30-gauss}: well-definedness
(\Cref{thm:ch30-gauss-well-defined}), centers and supports
(\Cref{thm:ch30-gauss-centers-supports}), the relation between scaling and
renormalization (\Cref{prop:ch30-gauss-scaling}), the characterization of the
image as the $\mathbf G_m(k)$-invariant slice
(\Cref{thm:ch30-gauss-image}), and the model dictionary
(\Cref{thm:ch30-gauss-model-dictionary}).  \Cref{sec:ch32-gauss} builds on
these and does not reconstruct them.

\item\label{itm:ch32-import:tc}
From \Cref{chap:bj-test-configurations}: the definition of a test
configuration \Cref{def:bj-test-configuration}, integral closedness
\Cref{def:integrally-closed-tc}, and the correspondence between PL functions
and vertical divisors \Cref{thm:tc-pl-correspondence}.  These are used only
in the proof of \Cref{thm:triv-divisorial-density} and in
\Cref{sec:ch32-gauss}.

\item\label{itm:ch32-import:nets}
The net formalism of \Cref{app:nets}, in particular
\Cref{thm:appE-net-compactness} and \Cref{thm:appE-cluster-subnet}.
\end{enumerate}

Nothing in this chapter uses the envelope property, energy, or the
Monge--Amp\`ere operator.  The examples of \Cref{sec:trivial-valuation-cone},
computed there for the affine cone of valuations of a curve and of the affine
plane together with their homogeneous potentials, are the concrete
counterpart of the present chapter; the reader who wants numbers before
definitions should read that section first.

\section{Semivaluations trivial on the ground field}
\label{sec:ch32-semivaluations}

\subsection{Seminorms and semivaluations}
\label{sec:ch32-seminorms}

Let $U=\Spec A\subset X$ be affine.  A multiplicative seminorm on $A$
extending the trivial norm on $k$ is a map
\[
 |\cdot|_x:A\longrightarrow \mathbf R_{\geq0}
\]
such that $|0|_x=0$, $|1|_x=1$,
$|fg|_x=|f|_x|g|_x$, and
$|f+g|_x\leq\max\{|f|_x,|g|_x\}$.  We use the additive notation
\begin{equation}
\label{eq:triv-additive-semivaluation}
             v_x(f):=-\log|f|_x\in\mathbf R\cup\{+\infty\},
             \qquad -\log0:=+\infty .
\end{equation}
Thus $v_x(fg)=v_x(f)+v_x(g)$ and
$v_x(f+g)\geq\min\{v_x(f),v_x(g)\}$.  Since the ground norm is
trivial, $v_x(a)=0$ for $a\in k^\times$.

The translation between the two notations is
\Cref{lem:seminorm-semivaluation-dictionary} and will be used without
comment.  Two elementary consequences of the axioms are used so often that
we record them here.

\begin{lemma}
\label{lem:ch32-semivaluation-calculus}
Let $v$ be a semivaluation on a ring $A$, trivial on $k^\times$.
\begin{enumerate}
\item\label{itm:ch32-sc:units} $v(a)=0$ for every $a\in k^\times$, and
$v(af)=v(f)$ for $a\in k^\times$ and $f\in A$.
\item\label{itm:ch32-sc:strict} If $v(f)\neq v(g)$, then
$v(f+g)=\min\{v(f),v(g)\}$.
\item\label{itm:ch32-sc:unique} If in a finite sum $f=\sum_{i}f_i$ the
minimum $\min_iv(f_i)$ is attained exactly once, then
$v(f)=\min_iv(f_i)$.
\item\label{itm:ch32-sc:sum} $v\bigl(\sum_if_i\bigr)\geq\min_iv(f_i)$ for
every finite family.
\end{enumerate}
\end{lemma}

\begin{proof}
\cref{itm:ch32-sc:units} is triviality on $k^\times$ together with
multiplicativity.  \cref{itm:ch32-sc:sum} is the ultrametric inequality and
induction.  For \cref{itm:ch32-sc:strict}, say $v(f)<v(g)$; then
$v(f+g)\geq v(f)$, and $v(f)=v((f+g)-g)\geq\min\{v(f+g),v(g)\}$; since
$v(f)<v(g)$, the minimum must be $v(f+g)$, so $v(f)\geq v(f+g)$.
\cref{itm:ch32-sc:unique} follows from \cref{itm:ch32-sc:strict} by grouping
the terms other than the minimizing one, whose sum has value strictly
greater than the minimum by \cref{itm:ch32-sc:sum}.  This last statement is
the mechanism behind every ``no cancellation'' argument below, in
particular in \Cref{thm:ch30-gauss-well-defined} and in
\Cref{ex:ch32-blowup-surface}.
\end{proof}

\begin{definition}
\label{def:triv-support}
The \emph{support} of $x$ on $U$ is the prime ideal
\[
                 \mathfrak p_x:=\{f\in A\mid v_x(f)=+\infty\}.
\]
The semivaluation descends to $A/\mathfrak p_x$, hence extends to a
real valuation of its fraction field
$\kappa(\mathfrak p_x)$.  Globally, the closure of
$\{\mathfrak p_x\}$ is an irreducible subvariety, denoted
$s(x)$, and called the support of $x$.
\end{definition}

That $\mathfrak p_x$ is prime is the multiplicativity of $v_x$: if
$v_x(fg)=+\infty$, then $v_x(f)+v_x(g)=+\infty$, so one of the two
summands is $+\infty$.  That the induced map on $A/\mathfrak p_x$
extends to the fraction field is the usual construction
$v_x(f/g)=v_x(f)-v_x(g)$, legitimate because both values are finite there.
By a mild and universal abuse we write $s(x)$ both for the scheme point
$\mathfrak p_x$ and for the integral closed subscheme it generates; the
context always disambiguates, and the two usages are related by
$s(x)=\overline{\{\mathfrak p_x\}}$.

The word \emph{semivaluation} records precisely the possibility
$\mathfrak p_x\ne0$.  A point has \emph{full support} if $s(x)$ is
an irreducible component of $X$.  For integral $X$, the full-support
locus is
\begin{equation}
\label{eq:triv-full-support-locus}
 X^{\val}:=\{v:k(X)^\times\to\mathbf R\mid v|_{k^\times}=0,
                    \text{$v$ has a center on $X$}\}.
\end{equation}
Properness guarantees the center in \eqref{eq:triv-full-support-locus}; it is
the unique point $c(v)\in X$ whose local ring is dominated by the valuation
ring.  For a general analytic point, $c(v)$ is defined after passing to the
support.  One always has $c(v)\in s(v)$.

For a possibly reducible $X$ with irreducible components $X_\alpha$ we
set $X^{\val}:=\coprod_\alpha X_\alpha^{\val}$, the disjoint union being
justified by \Cref{prop:ch32-support-partition} below: a semivaluation of
full support determines, and is determined by, the component carrying it.

\subsection{The partition by support}
\label{sec:ch32-partition}

\begin{proposition}
\label{prop:ch32-support-partition}
\index{support!partition of the analytification}
Let $X$ be a variety over $k$.  The assignment $x\mapsto s(x)$ sets up
a bijection
\[
 X^{\An}\;\xrightarrow{\ \sim\ }\;\coprod_{Y\subset X}Y^{\val},
\]
where $Y$ ranges over the integral closed subschemes of $X$ and
$Y^{\val}$ denotes the set of real valuations of $k(Y)^\times$ trivial on
$k^\times$, the point $x$ corresponding to the induced valuation of
$k(s(x))$.  Under this bijection:
\begin{enumerate}
\item\label{itm:ch32-sp:full} $x$ has full support if and only if the
corresponding $Y$ is an irreducible component of $X$;
\item\label{itm:ch32-sp:triv} the trivial valuation of $k(Y)$ corresponds
to the point denoted $v_{Y,\triv}$, and for $Y=\{p\}$ a closed point this
is the unique point of $X^{\An}$ with support $p$;
\item\label{itm:ch32-sp:closed} for $Y\subset X$ closed, the inclusion
induces a homeomorphism of $Y^{\An}$ onto the closed subset
$\{x\in X^{\An}:s(x)\subset Y\}$ of $X^{\An}$.
\end{enumerate}
\end{proposition}

\begin{proof}
Let $x\in X^{\An}$.  By \Cref{prop:berk-pair-description} the point $x$
is a pair $(\xi,|\cdot|)$ with $\xi\in X$ a scheme point and
$|\cdot|$ a multiplicative absolute value on the residue field
$\kappa(\xi)$ extending that of $k$; in additive notation, a valuation of
$\kappa(\xi)^\times$ trivial on $k^\times$.  By
\Cref{def:triv-support} the point $\xi$ is exactly $\mathfrak p_x$,
i.e.\ the generic point of $s(x)$, and $\kappa(\xi)=k(s(x))$.  This is
the asserted bijection; its inverse sends a valuation $w$ of $k(Y)$ to
the semivaluation $f\mapsto w(f|_Y)$, with the convention that
$w(0)=+\infty$, which is a semivaluation on $\cO_X$ with support the
generic point of $Y$.

\cref{itm:ch32-sp:full} is the definition of full support:
$s(x)$ is an irreducible component of $X$ exactly when its generic point
is a generic point of $X$.

\cref{itm:ch32-sp:triv}  If $Y=\{p\}$ with $p$ a closed point, then
$k(Y)=\kappa(p)=k$ since $k$ is algebraically closed, and the only
valuation of $k^\times$ trivial on $k^\times$ is the trivial one.  In
general $v_{Y,\triv}$ is defined by $v_{Y,\triv}(f)=0$ if $f|_Y\neq0$
and $+\infty$ otherwise.

\cref{itm:ch32-sp:closed}  A closed immersion $Y\hookrightarrow X$ induces a
closed immersion of analytic spaces, hence a homeomorphism onto a closed
subset (\Cref{sec:berk-functoriality}); its image is the set of points whose
seminorm kills the ideal of $Y$, that is, whose support is contained in
$Y$.
\end{proof}

\begin{remark}
\label{rem:ch32-partition-not-open}
The pieces $Y^{\val}$ of \Cref{prop:ch32-support-partition} are neither
open nor closed in general, and the partition is not locally finite.  The
closed subsets of $X^{\An}$ that the partition does produce are the unions
$\{s(x)\subset Y\}=Y^{\An}$, and the open ones are the unions
$\{s(x)\cap U\neq\varnothing\}=U^{\An}$ for $U\subset X$ open.  A point
of $X^{\val}$ can be a limit of points of $X^{\val}$ that converges to a
point of proper support: this is the content of
\Cref{prop:ch32-proper-support-dense}, and it is the reason the partition
cannot be used to reduce statements about $X^{\An}$ to statements about
valuations.
\end{remark}

\begin{definition}
\label{def:trivial-analytification}
The set $X^{\An}$ is obtained by gluing the spaces of seminorms on affine
opens.  Its topology is the weakest topology for which $U^{\An}$ is open
for every Zariski open $U\subset X$, and
\[
                  x\longmapsto |f(x)|=\rme^{-v_x(f)}
\]
is continuous on $U^{\An}$ for every $f\in\cO_X(U)$.
\end{definition}

This description is intrinsic: an analytic point is equivalently a pair
$(\xi,|\cdot|)$, with $\xi\in X$ and a multiplicative absolute value on
$\kappa(\xi)$ extending that of $k$.  The point $\xi$ is the generic
point of $s(x)$, not in general the center $c(x)$.

\subsection{Evaluation on ideals}
\label{sec:ch32-ideals}

For a coherent ideal $\mathfrak a\subset\cO_X$, set
\begin{equation}
\label{eq:ideal-evaluation-triv}
 v(\mathfrak a):=
 \min\{v(f)\mid f\text{ belongs to a local generating set at }c(v)\}.
\end{equation}
The minimum is independent of the generating set.  Moreover,
\begin{equation}
\label{eq:ideal-tropical-identities}
\begin{split}
 v(\mathfrak a\mathfrak b)&=v(\mathfrak a)+v(\mathfrak b),\\
 v(\mathfrak a+\mathfrak b)&=\min\{v(\mathfrak a),v(\mathfrak b)\},\\
 v(\mathfrak a)=+\infty&\Longleftrightarrow
                         s(v)\subset V(\mathfrak a).
\end{split}
\end{equation}
The functions $v\mapsto v(\mathfrak a)$, with values in the compact
extended half-line $[0,+\infty]$, are continuous and generate the topology.
In this sense $X^{\An}$ is the tropical spectrum of the semiring of coherent
ideals, with sum and product interchanged with minimum and addition.

We now prove each of these assertions.  The continuity and topology
statements are \Cref{thm:ch32-ideals-generate} below, after the center map
has been studied; the algebraic ones are the following proposition, whose
proof is short but which is used constantly.

\begin{proposition}
\label{prop:ch32-tropical-identities}
\index{ideal!evaluation by a semivaluation}
Let $X$ be a proper $k$-variety and $v\in X^{\An}$, with center
$c(v)\in X$ and support $s(v)$.  Then:
\begin{enumerate}
\item\label{itm:ch32-ti:welldef}
the quantity \eqref{eq:ideal-evaluation-triv} is well defined, that is,
independent of the chosen finite generating set of the stalk
$\mathfrak a_{c(v)}$, and equals $\min\{v(f):f\in\mathfrak a_{c(v)}\}$;
\item\label{itm:ch32-ti:nonneg}
$v(\mathfrak a)\in[0,+\infty]$, with $v(\cO_X)=0$ and $v(0)=+\infty$;
\item\label{itm:ch32-ti:identities}
the three identities \eqref{eq:ideal-tropical-identities} hold;
\item\label{itm:ch32-ti:center}
$v(\mathfrak a)>0$ if and only if $c(v)\in Z(\mathfrak a)$;
\item\label{itm:ch32-ti:monotone}
if $\mathfrak a\subset\mathfrak b$, then
$v(\mathfrak a)\geq v(\mathfrak b)$; in particular $v(\cdot)$ is order
reversing, and $v(\mathfrak a)=v(\overline{\mathfrak a})$ for the integral
closure $\overline{\mathfrak a}$.
\end{enumerate}
\end{proposition}

\begin{proof}
\cref{itm:ch32-ti:welldef}  Let $f_1,\ldots,f_r$ generate
$\mathfrak a_{c(v)}$ as an ideal of $R:=\cO_{X,c(v)}$.  Every
$f\in\mathfrak a_{c(v)}$ is $f=\sum_ig_if_i$ with $g_i\in R$, and
$v\geq0$ on $R$ by definition of the center, so
\[
 v(f)\geq\min_i\bigl(v(g_i)+v(f_i)\bigr)\geq\min_iv(f_i)
\]
by \Cref{lem:ch32-semivaluation-calculus}\,\cref{itm:ch32-sc:sum}.  Since the
$f_i$ themselves lie in $\mathfrak a_{c(v)}$, the minimum over the
generating set equals the infimum over the whole stalk, and is attained.
Independence of the generating set follows because the answer has been
identified with $\min\{v(f):f\in\mathfrak a_{c(v)}\}$, which mentions no
generators.  This is the ideal evaluation of \eqref{eq:appA-ideal-value},
specialized to the present setting.

\cref{itm:ch32-ti:nonneg}  Nonnegativity is $v\geq0$ on $R$.  For
$\mathfrak a=\cO_X$ the stalk contains $1$ and $v(1)=0$.  For
$\mathfrak a=0$ the only element is $0$, of value $+\infty$.

\cref{itm:ch32-ti:identities}  If $\mathfrak a_{c(v)}=(f_i)_i$ and
$\mathfrak b_{c(v)}=(g_j)_j$, then $(\mathfrak a\mathfrak b)_{c(v)}$ is
generated by the products $f_ig_j$, and
\[
 \min_{i,j}\bigl(v(f_i)+v(g_j)\bigr)=\min_iv(f_i)+\min_jv(g_j),
\]
which is the first identity.  The stalk of $\mathfrak a+\mathfrak b$ is
generated by the $f_i$ together with the $g_j$, and the minimum over a
union of two finite sets is the minimum of the two minima; this is the second
identity.  For the third, $v(\mathfrak a)=+\infty$ means that every element
of $\mathfrak a_{c(v)}$ lies in the support prime of $v$; since the
support prime is the kernel of $v$ on the stalk and $\mathfrak a$ is
coherent, this says $\mathfrak a\subset\mathcal I_{s(v)}$ in a neighbourhood
of $c(v)$.  As $s(v)$ is irreducible with generic point in every
neighbourhood of $c(v)$ --- indeed $c(v)\in s(v)$ --- this is equivalent
to $\mathfrak a\subset\mathcal I_{s(v)}$, that is,
$s(v)\subset V(\mathfrak a)$.  These are the identities
\eqref{eq:appA-ideal-value} of \Cref{prop:appA-ideal-arithmetic} in the
present notation.

\cref{itm:ch32-ti:center}  By definition of the center, $v>0$ exactly on the
maximal ideal $\mathfrak m_{c(v)}\subset R$.  Hence
$\min_iv(f_i)>0$ if and only if every $f_i\in\mathfrak m_{c(v)}$, that
is, if and only if $\mathfrak a_{c(v)}\subset\mathfrak m_{c(v)}$, which is
exactly $c(v)\in Z(\mathfrak a)$.

\cref{itm:ch32-ti:monotone}  If $\mathfrak a\subset\mathfrak b$, then
$\mathfrak b=\mathfrak a+\mathfrak b$, so
$v(\mathfrak b)=\min\{v(\mathfrak a),v(\mathfrak b)\}\leq v(\mathfrak a)$.
The statement about integral closure is
\Cref{thm:appA-valuative-integral-closure}: an element integral over
$\mathfrak a$ satisfies a monic relation whose lower coefficients lie in
powers of $\mathfrak a$, and evaluating $v$ on that relation and using
\Cref{lem:ch32-semivaluation-calculus}\,\cref{itm:ch32-sc:sum} gives
$v(f)\geq v(\mathfrak a)$.
\end{proof}

\begin{theorem}
\label{thm:ch32-tropical-spectrum}
\index{tropical spectrum}
Let $A$ be a finitely generated $k$-algebra and let $\mathcal I(A)$ be
the set of ideals of $A$.  Write $\mathcal M_1(A)$ for the set of
semivaluations $v:A\to[0,+\infty]$ trivial on $k^\times$ and
nonnegative, with the topology of pointwise convergence on $A$, and write
$\Trop(A)$ for the set of maps $\chi:\mathcal I(A)\to[0,+\infty]$
satisfying
\[
\begin{gathered}
 \chi(A)=0,\qquad \chi(0)=+\infty,\\
 \chi(\mathfrak a\mathfrak b)=\chi(\mathfrak a)+\chi(\mathfrak b),\qquad
 \chi(\mathfrak a+\mathfrak b)=\min\{\chi(\mathfrak a),\chi(\mathfrak b)\},
\end{gathered}
\]
with the topology of pointwise convergence on $\mathcal I(A)$.  Then
\[
 \Phi:\mathcal M_1(A)\longrightarrow\Trop(A),
 \qquad
 \Phi(v)(\mathfrak a)=\min\{v(f):f\in\mathfrak a\},
\]
is a homeomorphism.
\end{theorem}

\begin{proof}
\emph{$\Phi$ is well defined.}  $A$ is noetherian, so every ideal is
finitely generated and the infimum defining $\Phi(v)(\mathfrak a)$ is a
minimum over a finite generating set, by the computation in
\Cref{prop:ch32-tropical-identities}\,\cref{itm:ch32-ti:welldef} with
$R=A$ --- which is legitimate here because $v\geq0$ on $A$ by
hypothesis.  The four displayed identities were verified in
\Cref{prop:ch32-tropical-identities}, with $\chi(A)=v(1)=0$.

\emph{$\Phi$ is injective.}  For $f\in A$ we have
$\Phi(v)\bigl((f)\bigr)=v(f)$, since $(f)$ is generated by $f$.  Hence
$\Phi(v)$ determines $v$.

\emph{$\Phi$ is surjective.}  Let $\chi\in\Trop(A)$ and define
$v(f):=\chi((f))$.  Then $v(1)=\chi(A)=0$ and $v(0)=\chi(0)=+\infty$.
Multiplicativity: $(fg)=(f)(g)$, so
$v(fg)=\chi((f)(g))=\chi((f))+\chi((g))=v(f)+v(g)$.  Ultrametricity: the
map $\chi$ is order reversing, because
$\mathfrak a\subset\mathfrak b$ gives
$\chi(\mathfrak b)=\chi(\mathfrak a+\mathfrak b)
 =\min\{\chi(\mathfrak a),\chi(\mathfrak b)\}\leq\chi(\mathfrak a)$; and
$(f+g)\subset(f)+(g)$, so
\[
 v(f+g)=\chi\bigl((f+g)\bigr)\geq\chi\bigl((f)+(g)\bigr)
 =\min\{v(f),v(g)\}.
\]
Nonnegativity is $\chi\geq0$, and triviality on $k^\times$ holds because
$(a)=A$ for $a\in k^\times$.  Thus $v\in\mathcal M_1(A)$.  Finally
$\Phi(v)=\chi$: for $\mathfrak a=(f_1,\ldots,f_r)=(f_1)+\cdots+(f_r)$,
\[
 \chi(\mathfrak a)=\min_i\chi\bigl((f_i)\bigr)=\min_iv(f_i)
 =\Phi(v)(\mathfrak a).
\]

\emph{$\Phi$ is a homeomorphism.}  Both topologies are topologies of
pointwise convergence, and by the two displayed formulas
$\Phi(v)(\mathfrak a)=\min_iv(f_i)$ and $v(f)=\Phi(v)((f))$, each
coordinate of one side is a continuous function of finitely many coordinates
of the other.  Hence $\Phi$ and $\Phi^{-1}$ are continuous.
\end{proof}

\begin{remark}
\label{rem:ch32-global-tropical}
For $X$ proper the same statement holds with $\mathcal I(A)$ replaced by
the set $\mathcal I(X)$ of coherent ideal sheaves and
$\mathcal M_1$ replaced by $X^{\An}$; this is the identification asserted
in \eqref{eq:ideal-tropical-identities} and stated in
\cite[\S1.1.3]{BJ22}, where it is attributed to the local picture of
\cite[\S1.2]{JM12}.  We shall not use the surjectivity half of the global
statement.  What we do use is the injectivity half, together with the
statement that the ideal evaluations generate the topology; both are proved
in \Cref{thm:ch32-ideals-generate}, and both are elementary once
\Cref{thm:ch32-center-compact} is available.  The reason for the asymmetry is
that a coherent ideal on a projective $X$ is not determined by its stalks
at a single point, so the reconstruction step of
\Cref{thm:ch32-tropical-spectrum} --- ``take $\mathfrak a=(f)$'' --- has
no global analogue without first choosing an ample line bundle; the device
that replaces it is \Cref{lem:ch32-sections-generate-local-ring}.
\end{remark}

\subsection{Evaluation on sections, and the trivial metric}
\label{sec:ch32-sections}

Let $L$ be a line bundle on $X$ and $s\in H^0(X,L)$.  For
$v\in X^{\An}$ choose a local trivialization of $L$ at $c(v)$; the
section $s$ becomes a germ in $\cO_{X,c(v)}$, and we set
\[
 v(s):=v\bigl(\text{that germ}\bigr)\in[0,+\infty],
 \qquad
 |s|(v):=\rme^{-v(s)}\in[0,1].
\]
The value does not depend on the trivialization, because two trivializations
differ by a unit of $\cO_{X,c(v)}$, which has value $0$.  Equivalently,
$v(s)=v(\mathfrak a_s)$ for the coherent ideal
$\mathfrak a_s:=\operatorname{im}\bigl(L^{-1}\xrightarrow{\;\cdot s\;}
\cO_X\bigr)$, whose zero locus is the divisor of $s$ when $s$ is a
regular section.

\begin{remark}
\label{rem:ch32-trivial-metric}
\index{metric!trivial}
The recipe $|s|(v)=\rme^{-v(s)}$ is the trivial metric on $L^{\An}$: the
unique continuous metric for which $|\tau|\equiv1$ on $c^{-1}(U)$ for
every trivializing section $\tau\in H^0(U,L)$ over an open $U\subset X$.
This is \cite[Remark~1.24]{BJ22}; in the language of \Cref{chap:metrics} it
is the model metric attached to the trivial test configuration
(\Cref{ex:metrics-trivial-bundle}).  All the metrics of \Cref{part:bj} are
measured against it, and all Fubini--Study functions of
\Cref{def:gf-fs-functions} are built from it.
\end{remark}

The following lemma is the substitute, over a projective base, for the
device ``take the principal ideal generated by $f$'' that made
\Cref{thm:ch32-tropical-spectrum} work.  It is used three times below and is
elementary.

\begin{lemma}
\label{lem:ch32-sections-generate-local-ring}
Let $X$ be a projective $k$-scheme and $L$ an ample line bundle on
$X$.  For every scheme point $\xi\in X$ and every
$f\in\cO_{X,\xi}$ there are $m\geq1$ and sections
$s,u\in H^0(X,L^{\otimes m})$ with $u(\xi)\neq0$ --- that is, $u$
generating $L^{\otimes m}$ at $\xi$ --- such that $f=s/u$ in
$\cO_{X,\xi}$.  Consequently, for every $v\in X^{\An}$ with
$c(v)=\xi$,
\[
 v(f)=v(s)-v(u)=v(s),
\]
the last equality because $v(u)=0$.
\end{lemma}

\begin{proof}
Replacing $L$ by a power we may assume $L$ very ample, so that
$X\subset\mathbf P^N$ is a closed subscheme with
$L=\cO_{\mathbf P^N}(1)|_X$.  Choose a linear form $t$ on
$\mathbf P^N$ with $t(\xi)\neq0$; this is possible since the linear forms
generate $\cO(1)$ at every point.  The open subscheme
$X_t:=X\cap D_+(t)$ is affine with coordinate ring generated over $k$ by
the functions $x_i/t$, $0\leq i\leq N$.  As $\xi\in X_t$, we may write
$f=g/h$ with $g,h\in\cO(X_t)$ and $h(\xi)\neq0$.  Each of $g,h$ is a
polynomial in the $x_i/t$; choosing $m$ at least as large as the two
degrees, we get homogeneous forms $P,Q$ of degree $m$ with
$g=P/t^m$, $h=Q/t^m$ on $X_t$.  Put $s:=P|_X$ and $u:=Q|_X$, two
sections of $L^{\otimes m}$.  Then $f=g/h=s/u$ on $X_t$, hence in
$\cO_{X,\xi}$, and $u(\xi)=h(\xi)t(\xi)^m\neq0$.

For the last assertion, $v(u)=0$ because $u$ generates
$L^{\otimes m}$ at $\xi=c(v)$, so in a trivialization it is a unit of
$\cO_{X,\xi}$.  Multiplicativity of $v$ applied to $s=fu$ gives
$v(s)=v(f)+v(u)=v(f)$.
\end{proof}

The next proposition is the valuative criterion promised in
\Cref{sec:ch32-standing}; it characterizes full support in terms of sections,
and is the first place where ampleness enters.  It is
\cite[Lemma~1.25]{BJ22}.  Recall that a section of a line bundle on $X$ is
\emph{regular} if it does not vanish identically along any irreducible
component of $X$.

\begin{proposition}
\label{prop:ch32-valuation-criterion}
\index{support!criterion by sections}
Let $X$ be a projective $k$-variety and $L$ an ample line bundle.
A point $v\in X^{\An}$ lies in $X^{\val}$, that is, has full support, if
and only if $v(s)<+\infty$ for every $m\geq1$ and every regular section
$s\in H^0(X,L^{\otimes m})$.
\end{proposition}

\begin{proof}
Suppose first that $v\notin X^{\val}$, so that $Y:=s(v)$ is not an
irreducible component of $X$.  Let $\mathcal I_Y\subset\cO_X$ be its
ideal sheaf; it is nonzero on every irreducible component of $X$, since
$Y$ is contained in some component and is not equal to it, and is not
contained in any other component's generic point.  By ampleness,
$\mathcal I_Y\otimes L^{\otimes m}$ is globally generated for $m\gg1$; in
particular there is $s\in H^0(X,\mathcal I_Y\otimes L^{\otimes m})
\subset H^0(X,L^{\otimes m})$ which does not vanish identically along any
irreducible component of $X$.  Such a section exists because global
generation implies that, for each generic point $\eta_\alpha$ of $X$, the
stalk $(\mathcal I_Y)_{\eta_\alpha}=\cO_{X,\eta_\alpha}$ is generated by
the images of the global sections, so the sections vanishing at
$\eta_\alpha$ form a proper subspace; a finite union of proper subspaces of
a vector space over an infinite field is not the whole space.  Then $s$ is
regular and vanishes along $Y=s(v)$, so $v(s)=+\infty$.

Conversely, if $v\in X^{\val}$, then $s(v)=X_\alpha$ is an irreducible
component, and a regular section $s$ does not vanish identically along
$X_\alpha$; hence its germ at $c(v)$ is not in the support prime of
$v$, and $v(s)<+\infty$.
\end{proof}

\begin{remark}
\label{rem:ch32-valuation-criterion-scope}
Ampleness in \Cref{prop:ch32-valuation-criterion} is used only to produce one
section vanishing along a prescribed proper closed subset and along no
component; any line bundle with that property would do.  The hypothesis that
$k$ is infinite is used in the ``finite union of proper subspaces''
step, and is automatic here since $k$ is algebraically closed.  The
criterion is the starting point of the theory of valuations of linear
growth, \Cref{def:bj-linear-growth}: there one asks not merely that
$v(s)$ be finite but that it grow at most linearly in $m$.
\end{remark}

\section{The topology, compactness, and countability}
\label{sec:ch32-topology}

\subsection{Compactness, twice}
\label{sec:ch32-compactness}

\begin{theorem}\label{thm:triv-berkovich-compactness}
If $X$ is proper over $k$, then $X^{\An}$ is compact Hausdorff.  If
$f:Y\to X$ is a morphism of proper $k$-schemes, composition of
seminorms defines a continuous map $f^{\An}:Y^{\An}\to X^{\An}$.
\end{theorem}

\begin{proof}
This is \Cref{thm:berk-proper-compact}, already proved for arbitrary
complete ground fields.  Applying it to the trivially valued field gives the
first assertion, and analytic functoriality gives the second.  The
semivaluation description used here is identified with that construction in
\cite[\S1.1]{BJ22}.
\end{proof}

The general theorem is not cheap: its proof in
\Cref{thm:berk-proper-compact} goes through Chow's lemma and the compactness
of the closed polydisc.  Over a trivially valued base there is a shorter and
more informative argument, which produces at the same time the anticontinuity
of the center map and the compactness of the sets $c^{-1}(U)$.  We give it
in full, because these compact sets are used repeatedly --- in the proof of
\Cref{thm:ch32-ideals-generate}, in the study of normalized links, and in
every argument below that localizes on $X$.

\begin{theorem}
\label{thm:ch32-center-compact}
\index{center map!compact fibres}
Let $X$ be a proper $k$-scheme of finite type.
\begin{enumerate}
\item\label{itm:ch32-cc:affine}
For every affine open $U=\Spec A\subset X$, the set
\[
 c^{-1}(U)=\{v\in U^{\An}:v\geq0\text{ on }A\}
\]
is a compact subset of $X^{\An}$; it is canonically the Berkovich spectrum
$\mathcal M(A)$ of $A$ equipped with the trivial norm.
\item\label{itm:ch32-cc:general}
For every open $U\subset X$, the set $c^{-1}(U)$ is compact, hence
closed; equivalently, $c:X^{\An}\to X$ is anticontinuous.
\item\label{itm:ch32-cc:compact}
$X^{\An}=c^{-1}(X)$ is compact, and it is Hausdorff.
\end{enumerate}
\end{theorem}

\begin{proof}
\cref{itm:ch32-cc:affine}  We first check the displayed description of
$c^{-1}(U)$.  If $v\in X^{\An}$ has $c(v)\in U$, then its support meets
$U$, so $v\in U^{\An}$, and $v\geq0$ on
$\cO_{X,c(v)}\supset A$.  Conversely let $v\in U^{\An}$ with $v\geq0$
on $A$.  Then $\mathfrak q:=\{f\in A:v(f)>0\}$ is a prime ideal of $A$:
it is closed under addition by ultrametricity, absorbs multiplication by
\Cref{lem:ch32-semivaluation-calculus}\,\cref{itm:ch32-sc:sum} and
nonnegativity, and is prime by multiplicativity.  The corresponding point
$\xi\in U$ satisfies $v\geq0$ on $A_{\mathfrak q}=\cO_{X,\xi}$ --- a
fraction $f/g$ with $g\notin\mathfrak q$ has $v(g)=0$ hence
$v(f/g)=v(f)\geq0$ --- and $v>0$ on $\mathfrak q A_{\mathfrak q}$.
Thus $v$ dominates $\cO_{X,\xi}$, so $\xi$ is a center of $v$, and by
uniqueness of centers on a separated scheme
(\Cref{prop:centers-proper-separated}) we get $c(v)=\xi\in U$.

Now equip $A$ with the trivial norm, making it a Banach ring.  By
definition, $\mathcal M(A)$ is the set of multiplicative seminorms
$|\cdot|$ on $A$ bounded by the norm, that is, with $|f|\leq1$ for all
$f\in A$; in additive notation, the set of semivaluations $v$ on $A$
with $v\geq0$.  Triviality on $k^\times$ is automatic: for
$a\in k^\times$ both $v(a)\geq0$ and $v(a^{-1})\geq0$, and their sum is
$v(1)=0$.  Thus $c^{-1}(U)=\mathcal M(A)$ as a set, and the topologies
agree, both being the topology of pointwise convergence on $A$.  By
\Cref{prop:berk-compact-spectrum} the space $\mathcal M(A)$ is compact.

\cref{itm:ch32-cc:general}  A scheme of finite type is noetherian, so an open
$U\subset X$ is a finite union $U=U_1\cup\cdots\cup U_r$ of affine opens.
Hence $c^{-1}(U)=\bigcup_ic^{-1}(U_i)$ is a finite union of compact sets,
therefore compact.  Since $X^{\An}$ is Hausdorff --- see below --- compact
subsets are closed, which is precisely the anticontinuity of $c$: the
preimage of an open set is closed.

\cref{itm:ch32-cc:compact}  Every $v\in X^{\An}$ has a center because $X$
is proper: the valuation ring of the induced valuation on $k(s(v))$
provides, by the valuative criterion of properness, a unique
$k$-morphism $\Spec\cO_v\to X$ extending $\Spec k(s(v))\to X$, and its
closed point is $c(v)$; this is \Cref{prop:centers-proper-separated}.
Therefore $X^{\An}=c^{-1}(X)$ is compact by \cref{itm:ch32-cc:general}.
Hausdorffness holds because $X$ is separated: it is
\Cref{thm:berk-proper-compact}\,(1), and can also be seen directly from
\Cref{thm:ch32-ideals-generate}\,\cref{itm:ch32-ig:separate} below, whose
proof does not use Hausdorffness.
\end{proof}

\begin{remark}
\label{rem:ch32-logical-order}
The only point at which \Cref{thm:ch32-center-compact} borrows from
\Cref{thm:berk-proper-compact} is Hausdorffness, and even that can be
avoided: \Cref{thm:ch32-ideals-generate}\,\cref{itm:ch32-ig:separate}
produces, for any two distinct points, a continuous
$[0,+\infty]$-valued function separating them, and a space admitting a
separating family of continuous maps to a Hausdorff space is Hausdorff.  We
have not arranged the exposition to exploit this, since
\Cref{thm:berk-proper-compact} is available; but the reader who wants a
self-contained treatment of the trivially valued case may take
\Cref{prop:berk-compact-spectrum} as the only input.
\end{remark}

\subsection{Ideals generate the topology}
\label{sec:ch32-ideals-topology}

\begin{theorem}
\label{thm:ch32-ideals-generate}
\index{ideal!evaluations generate the topology}
Let $X$ be a projective $k$-variety.
\begin{enumerate}
\item\label{itm:ch32-ig:continuous}
For every coherent ideal $\mathfrak a\subset\cO_X$, the function
$v\mapsto v(\mathfrak a)$ is continuous from $X^{\An}$ to
$[0,+\infty]$.
\item\label{itm:ch32-ig:separate}
The family $\{v\mapsto v(\mathfrak a)\}_{\mathfrak a\in\mathcal I(X)}$
separates the points of $X^{\An}$.
\item\label{itm:ch32-ig:embed}
The induced map
\[
 \iota:X^{\An}\longrightarrow[0,+\infty]^{\mathcal I(X)},
 \qquad
 \iota(v)=\bigl(v(\mathfrak a)\bigr)_{\mathfrak a},
\]
is a homeomorphism onto a compact subset.  In particular the ideal
evaluations generate the topology of $X^{\An}$.
\end{enumerate}
\end{theorem}

\begin{proof}
\cref{itm:ch32-ig:continuous}  Choose a finite affine open cover
$X=U_1\cup\cdots\cup U_N$.  By
\Cref{thm:ch32-center-compact}\,\cref{itm:ch32-cc:affine} the sets
$F_j:=c^{-1}(U_j)$ are compact, hence closed, and they cover $X^{\An}$
because every $v$ has a center, which lies in some $U_j$.  A function on
a topological space which is continuous on each member of a finite closed
cover is continuous.  It therefore suffices to prove continuity of
$v\mapsto v(\mathfrak a)$ on each $F_j$.

Fix $j$ and write $U_j=\Spec A_j$.  Since $\mathfrak a$ is coherent and
$U_j$ affine, the ideal $\mathfrak a(U_j)\subset A_j$ generates
$\mathfrak a_\xi$ at every $\xi\in U_j$; choose generators
$f_1,\ldots,f_r\in\mathfrak a(U_j)$.  For $v\in F_j$ we have
$c(v)\in U_j$ and the $f_i$ generate $\mathfrak a_{c(v)}$, so
\[
 v(\mathfrak a)=\min_{1\leq i\leq r}v(f_i)
 \qquad\text{on }F_j
\]
by \Cref{prop:ch32-tropical-identities}\,\cref{itm:ch32-ti:welldef}.  Each
function $v\mapsto v(f_i)$ is continuous on $U_j^{\An}\supset F_j$ by
\Cref{def:trivial-analytification}, and a minimum of finitely many continuous
$[0,+\infty]$-valued functions is continuous.

\cref{itm:ch32-ig:separate}  Let $v\neq w$ in $X^{\An}$.  Suppose first
that $s(v)\neq s(w)$.  Taking $\mathfrak a=\mathcal I_{s(v)}$ gives
$v(\mathfrak a)=+\infty$ by
\Cref{prop:ch32-tropical-identities}\,\cref{itm:ch32-ti:identities}; if also
$w(\mathfrak a)=+\infty$, then $s(w)\subset s(v)$, and interchanging the
roles gives $s(v)\subset s(w)$, hence $s(v)=s(w)$, a contradiction.  So
one of the two ideals $\mathcal I_{s(v)}$, $\mathcal I_{s(w)}$
separates.

Suppose now $Y:=s(v)=s(w)$, so that $v$ and $w$ are distinct
valuations of $k(Y)^\times$.  Pick $f\in k(Y)^\times$ with
$v(f)\neq w(f)$.  Let $L$ be an ample line bundle on $X$.  Applying
\Cref{lem:ch32-sections-generate-local-ring} to the projective integral
scheme $Y$, the ample bundle $L|_Y$ and the generic point of $Y$, we
obtain $m_0$ and nonzero
$s_0,u_0\in H^0(Y,L^{\otimes m_0}|_Y)$ with $f=s_0/u_0$.  By Serre
vanishing the restriction map
$H^0(X,L^{\otimes m})\to H^0(Y,L^{\otimes m}|_Y)$ is surjective for
$m\gg1$; fixing such an $m\geq m_0$ and a nonzero
$w_0\in H^0(Y,L^{\otimes(m-m_0)}|_Y)$, we have
$f=(s_0w_0)/(u_0w_0)$ with both terms nonzero sections of
$L^{\otimes m}|_Y$, and we may lift them to
$s,u\in H^0(X,L^{\otimes m})$.  Then
$v(s),v(u),w(s),w(u)$ are all finite, since $s|_Y\neq0\neq u|_Y$ and
$Y$ is the support of both points, and
\[
 v(s)-v(u)=v(f)\neq w(f)=w(s)-w(u).
\]
Hence $v(s)\neq w(s)$ or $v(u)\neq w(u)$.  Passing to the associated
ideals $\mathfrak a_s,\mathfrak a_u$ of
\Cref{sec:ch32-sections} and using $v(\mathfrak a_s)=v(s)$, we obtain an
ideal separating $v$ from $w$.

\cref{itm:ch32-ig:embed}  The map $\iota$ is continuous by
\cref{itm:ch32-ig:continuous} and injective by \cref{itm:ch32-ig:separate}.
Its source is compact (\Cref{thm:ch32-center-compact}) and its target is
Hausdorff, being a product of compact Hausdorff spaces
$[0,+\infty]\simeq[0,1]$.  A continuous injection from a compact space to
a Hausdorff space is a homeomorphism onto its image, and the image is
compact.  Since the product topology is generated by the coordinate
projections, the induced topology on $X^{\An}$ is generated by the
functions $v\mapsto v(\mathfrak a)$.
\end{proof}

\begin{remark}
\label{rem:ch32-subbasis}
By \Cref{thm:ch32-ideals-generate}\,\cref{itm:ch32-ig:embed} a subbasis for
the topology of $X^{\An}$ is given by the sets
\[
 \{v:v(\mathfrak a)<t\}\quad\text{and}\quad\{v:v(\mathfrak a)>t\},
 \qquad \mathfrak a\in\mathcal I(X),\ t\in\mathbf R .
\]
Both kinds occur.  Near a point of $X^{\triv}$ only the first kind is
informative: at a trivial semivaluation every ideal evaluation lies in
$\{0,+\infty\}$, so a condition $v(\mathfrak a)>t$ with $t<0$ is
vacuous and one with $t\geq0$ already forces $v(\mathfrak a)=+\infty$.
This observation is what makes the neighbourhood bases of
\Cref{lem:ch32-neighbourhood-basis-triv} and
\Cref{thm:ch32-not-first-countable} manageable.
\end{remark}

\begin{lemma}
\label{lem:ch32-neighbourhood-basis-triv}
Let $X$ be an integral projective $k$-variety and $v_{X,\triv}$ the
trivial valuation of $k(X)$.  The sets
\[
 W(\mathfrak a_1,\ldots,\mathfrak a_r;\varepsilon)
 :=\{v\in X^{\An}:v(\mathfrak a_i)<\varepsilon,\ 1\leq i\leq r\},
 \qquad \mathfrak a_i\neq0,\ \varepsilon>0,
\]
form a neighbourhood basis of $v_{X,\triv}$.
\end{lemma}

\begin{proof}
Each such set is open by
\Cref{thm:ch32-ideals-generate}\,\cref{itm:ch32-ig:continuous} and contains
$v_{X,\triv}$, because a nonzero ideal on an integral scheme has stalk
$\cO_{X,\eta}$ at the generic point $\eta=c(v_{X,\triv})$, so
$v_{X,\triv}(\mathfrak a_i)=0$.  Conversely, a basic open set containing
$v_{X,\triv}$ is a finite intersection of subbasic sets as in
\Cref{rem:ch32-subbasis}.  A subbasic set of the form
$\{v(\mathfrak a)>t\}$ containing $v_{X,\triv}$ forces
$0=v_{X,\triv}(\mathfrak a)>t$, so $t<0$ and the set is all of
$X^{\An}$ because $v\geq0$ always; such factors may be discarded.  A
subbasic set $\{v(\mathfrak a)<t\}$ containing $v_{X,\triv}$ has
$t>0$, and $\mathfrak a\neq0$ since otherwise
$v_{X,\triv}(\mathfrak a)=+\infty$.  Taking $\varepsilon$ to be the
smallest of the finitely many $t$'s gives the claim.
\end{proof}

\subsection{Countability: metrizable or not even first countable}
\label{sec:ch32-countability}

Compactness does not imply metrizability.  Over an uncountable field,
$X^{\An}$ is generally not first countable.  All compactness arguments in
this part therefore use nets.  A sequence may be used only after an explicit
countability or metrization reduction.

The two extreme cases are settled by the following two statements, and they
are the reason for the standing convention.  We first record an elementary
counting lemma.

\begin{lemma}
\label{lem:ch32-countably-many-ideals}
If $k$ is countable and $X$ is a $k$-scheme of finite type, then the
set $\mathcal I(X)$ of coherent ideal sheaves on $X$ is countable.
\end{lemma}

\begin{proof}
Fix a finite affine open cover $X=U_1\cup\cdots\cup U_N$, with
$U_j=\Spec A_j$ and $A_j$ a finitely generated $k$-algebra, hence a
countable set.  A coherent ideal is determined by the finitely generated
ideals $\mathfrak a(U_j)\subset A_j$, and there are only countably many
finite subsets of the countable set $A_j$, hence countably many finitely
generated ideals of $A_j$.  A coherent sheaf is determined by its
restrictions to the members of the cover, so $\mathcal I(X)$ injects into
the countable set $\prod_{j=1}^N\{\text{ideals of }A_j\}$.
\end{proof}

\begin{proposition}
\label{prop:ch32-metrizable}
\index{metrizability!of the trivially valued analytification}
If $k$ is countable and $X$ is a projective $k$-variety, then
$X^{\An}$ is a compact metrizable space.
\end{proposition}

\begin{proof}
By \Cref{thm:ch32-ideals-generate}\,\cref{itm:ch32-ig:embed}, $X^{\An}$ is
homeomorphic to a subspace of $[0,+\infty]^{\mathcal I(X)}$.  By
\Cref{lem:ch32-countably-many-ideals} the index set is countable, so the
product is a countable product of second-countable spaces and is therefore
second countable; hence so is the subspace $X^{\An}$.  A compact Hausdorff
second-countable space is metrizable by Urysohn's metrization theorem.
\end{proof}

\begin{lemma}
\label{lem:ch32-uncountable-field}
Let $k$ be an uncountable algebraically closed field and $X$ an integral
$k$-variety with $\dim X\geq1$.  If $(Z_i)_{i\in\mathbf N}$ is a
sequence of proper closed subsets of $X$, then
$X(k)\not\subset\bigcup_iZ_i$.
\end{lemma}

\begin{proof}
Fix a nonempty affine open $U=\Spec B\subset X$; since $X$ is of finite
type over $k$, the closed points of $U$ are closed in $X$, so it
suffices to find a closed point of $U$ outside all $Z_i\cap U$, and each
$Z_i\cap U$ is a proper closed subset of the integral affine variety
$U$ of dimension $d:=\dim X\geq1$.  By Noether normalization there is a
finite surjective morphism $\pi:U\to\mathbf A^d_k$.  A finite morphism is
closed and preserves dimension, so $W_i:=\pi(Z_i\cap U)$ is a closed
subset of $\mathbf A^d$ of dimension $<d$, in particular a proper closed
subset.  If we find $q\in\mathbf A^d(k)\setminus\bigcup_iW_i$, then any
preimage $p\in U(k)$ of $q$ --- nonempty by surjectivity of the finite
morphism $\pi$, and consisting of closed points since $k$ is
algebraically closed --- avoids every $Z_i$.  It therefore suffices to
show: \emph{$k^d$ is not the union of countably many proper closed
subsets}.

Write $Z_i\subset V(f_i)$ for some nonzero
$f_i\in k[x_1,\ldots,x_d]$.  We prove by induction on $d$ that
$k^d\neq\bigcup_iV(f_i)$ for any sequence of nonzero polynomials.  For
$d=1$ each $V(f_i)$ is finite, so the union is countable while $k$ is
uncountable.  For $d\geq2$, consider for $c\in k$ the hyperplane
$H_c=\{x_d=c\}\simeq k^{d-1}$.  The restriction $f_i|_{H_c}$, a
polynomial in $d-1$ variables, vanishes identically if and only if
$(x_d-c)$ divides $f_i$, which happens for at most $\deg f_i$ values
of $c$.  The set $B$ of $c\in k$ for which some $f_i|_{H_c}$ vanishes
identically is thus a countable union of finite sets, hence countable, and
we may choose $c\in k\setminus B$.  On $H_c\simeq k^{d-1}$ each
$f_i|_{H_c}$ is a nonzero polynomial, and the inductive hypothesis provides
a point of $H_c$ outside all $V(f_i|_{H_c})$.
\end{proof}

\begin{theorem}
\label{thm:ch32-not-first-countable}
\index{first countability!failure of}
Let $k$ be uncountable and algebraically closed, and let $X$ be an
integral projective $k$-variety with $\dim X\geq1$.  Then $X^{\An}$ is
not first countable at $v_{X,\triv}$.  In particular $X^{\An}$ is
compact Hausdorff and not metrizable.
\end{theorem}

\begin{proof}
Suppose $(N_\ell)_{\ell\in\mathbf N}$ were a countable neighbourhood basis
at $v:=v_{X,\triv}$.  By \Cref{lem:ch32-neighbourhood-basis-triv} we may
assume
\[
 N_\ell=W(\mathfrak a_{\ell,1},\ldots,\mathfrak a_{\ell,r_\ell};
 \varepsilon_\ell)
\]
with all $\mathfrak a_{\ell,j}\neq0$.  The zero loci
$Z(\mathfrak a_{\ell,j})$ form a countable family of proper closed subsets
of $X$; by \Cref{lem:ch32-uncountable-field} there is a closed point
$p\in X(k)$ lying in none of them.  Consider the trivial semivaluation
$v_{p,\triv}$ supported at $p$.  For each $\ell,j$ we have
$p\notin Z(\mathfrak a_{\ell,j})$, so the stalk of
$\mathfrak a_{\ell,j}$ at $p$ is the whole local ring and therefore
$v_{p,\triv}(\mathfrak a_{\ell,j})=0<\varepsilon_\ell$.  Hence
$v_{p,\triv}\in N_\ell$ for every $\ell$.

On the other hand let $\mathfrak m_p\subset\cO_X$ be the ideal sheaf of
$p$ and set $N:=\{w\in X^{\An}:w(\mathfrak m_p)<1\}$.  This is an open
neighbourhood of $v$, since $v(\mathfrak m_p)=0$, and it does not contain
$v_{p,\triv}$, since $v_{p,\triv}(\mathfrak m_p)=+\infty$ --- the support
of $v_{p,\triv}$ is $\{p\}=Z(\mathfrak m_p)$, so
\Cref{prop:ch32-tropical-identities}\,\cref{itm:ch32-ti:identities} applies.
Thus no $N_\ell$ is contained in $N$, contradicting the assumption that
the $N_\ell$ form a basis.

Since a metrizable space is first countable, $X^{\An}$ is not metrizable;
it is compact Hausdorff by \Cref{thm:ch32-center-compact}.
\end{proof}

\begin{warning}{Sequences are not enough}
\Cref{thm:ch32-not-first-countable} is the reason every compactness
argument in \Crefrange{chap:trivial-analytification}{chap:bj-stability} is
formulated with nets, in the sense of \Cref{app:nets}, and why the
regularized supremum $(\sup_i\varphi_i)^*$ of the front matter is taken
over families rather than sequences.  Concretely: a function on $X^{\An}$
that is continuous along every sequence need not be continuous; a family of
psh functions may have a supremum that no sequence from the family
approaches; and an increasing net of continuous functions may fail to admit a
cofinal increasing subsequence.  The three chapters that most depend on this
discipline are \Cref{chap:bj-envelope-statement},
\Cref{chap:bj-envelope-consequences} and \Cref{chap:bj-strong-topologies}.
There is a partial converse in the literature --- every point in the closure
of a subset of $X^{\An}$ is the limit of a \emph{sequence} from that
subset, so closed subsets are sequentially compact --- but it is a
nontrivial theorem about Berkovich spaces which we neither prove nor use;
see \cite[Remark~1.1]{BJ22} for the statement and its provenance.
\end{warning}

\begin{remark}
\label{rem:ch32-countability-summary}
Three statements hold for every algebraically closed $k$:
$X^{\An}$ is compact Hausdorff; the ideal evaluations generate its
topology; and $X^{\divis}$ is dense in it
(\Cref{thm:triv-divisorial-density}).  Density of a countable set is not
available in general, since $X^{\divis}$ has the cardinality of $k$ when
$\dim X\geq1$.  The reader should keep the two regimes separate: over
$\overline{\mathbf Q}$ the space $X^{\An}$ is a compact metric space and
sequential arguments are legitimate; over $\mathbf C$ they are not.
\end{remark}

\section{Centers, supports, and specialization}
\label{sec:ch32-centers}

\subsection{Two maps to the scheme \texorpdfstring{$X$}{X}}
\label{sec:ch32-two-maps}

There are two natural maps from $X^{\An}$ to the underlying topological
space of the scheme $X$, and they behave in opposite ways.

\begin{theorem}
\label{thm:ch32-support-center-maps}
\index{center map!anticontinuity}\index{support map!continuity}
Let $X$ be a proper $k$-scheme of finite type.
\begin{enumerate}
\item\label{itm:ch32-scm:s}
The support map $s:X^{\An}\to X$, $v\mapsto$ the generic point of
$s(v)$, is continuous for the Zariski topology on $X$; indeed
$s^{-1}(U)=U^{\An}$ for every open $U\subset X$.
\item\label{itm:ch32-scm:c}
The center map $c:X^{\An}\to X$ is anticontinuous: for every open
$U\subset X$, the set $c^{-1}(U)$ is compact and closed in
$X^{\An}$, and for every closed $Z\subset X$ the set $c^{-1}(Z)$ is
open.
\item\label{itm:ch32-scm:order}
For every $v\in X^{\An}$ one has $c(v)\in s(v)$; equivalently, writing
$Y=s(v)$ and $Z=\overline{\{c(v)\}}$, one has $Z\subset Y$.
\item\label{itm:ch32-scm:neither}
Neither $s$ nor $c$ is continuous for the other's convention: for
$\dim X\geq1$, $s$ is not anticontinuous and $c$ is not continuous.
\end{enumerate}
\end{theorem}

\begin{proof}
\cref{itm:ch32-scm:s}  By construction $U^{\An}$ is the set of points whose
support meets $U$, that is, whose support generic point lies in $U$; and
$U^{\An}$ is open in $X^{\An}$ by \Cref{def:trivial-analytification}.

\cref{itm:ch32-scm:c}  This is
\Cref{thm:ch32-center-compact}\,\cref{itm:ch32-cc:general}; the statement
about closed $Z$ is the complementary formulation, since
$c^{-1}(Z)=X^{\An}\setminus c^{-1}(X\setminus Z)$.

\cref{itm:ch32-scm:order}  The valuation $v$ is a valuation of
$k(Y)=\kappa(s(v))$, and its valuation ring dominates a local ring of
$Y$ --- the valuative criterion is applied to the integral proper scheme
$Y$, giving a center $c_Y(v)\in Y$.  On the other hand, applying the
criterion to $X$ gives $c(v)\in X$.  The two agree, because the
composite $\Spec\cO_v\to Y\hookrightarrow X$ is a morphism extending
$\Spec k(Y)\to X$, and the center is unique by
\Cref{prop:centers-proper-separated}.  Hence $c(v)=c_Y(v)\in Y$.

\cref{itm:ch32-scm:neither}  Take $X$ integral with $\dim X\geq1$, let
$p\in X(k)$ be a closed point, $\eta$ the generic point, and let $E$ be
a prime divisor over $X$ with center $p$, for instance an exceptional
divisor of the blowup of $p$.  For $s$: the set
$s^{-1}(\{p\})=\{v_{p,\triv}\}$ is a single point and is not open, since
$v_{p,\triv}=\lim_{t\to+\infty}t\ord_E$ by
\Cref{prop:ch32-scaling-limits}, so every neighbourhood of it contains
points $t\ord_E$ of support $X$.  As $\{p\}$ is closed in $X$, this
shows $s$ is not anticontinuous.  For $c$: the singleton $\{p\}$ is
closed in $X$, so $c^{-1}(\{p\})$ is open by \cref{itm:ch32-scm:c}; were
$c$ continuous it would also be closed.  It is not closed, because it
contains $t\ord_E$ for every $t>0$ while its closure contains
$v_{X,\triv}=\lim_{t\to0^+}t\ord_E$, whose center is $\eta\neq p$.
Hence $c$ is not continuous.
\end{proof}

\begin{remark}
\label{rem:ch32-anticontinuity}
Anticontinuity of the center map is the trivially valued avatar of the
anticontinuity of the reduction map in Berkovich geometry,
\Cref{prop:berk-reduction-anticontinuous}.  It is what makes the sets
$c^{-1}(U)$ --- compact and closed, and $k$-analytic domains --- the
correct local pieces of $X^{\An}$: they are not open, so they do not form
a topological covering in the naive sense, but they do form a finite closed
covering, which is exactly what one needs to check continuity of a function,
as in the proof of \Cref{thm:ch32-ideals-generate}.  In the discretely
valued theory the analogous sets are the preimages of open subsets of the
special fibre under the reduction map, and they play the same role.
\end{remark}

\subsection{The partial order}
\label{sec:ch32-order}

There is an order on $X^{\An}$ defined by
\begin{equation}
\label{eq:semivaluation-order}
 v\geq w \quad\Longleftrightarrow\quad
 v(\mathfrak a)\geq w(\mathfrak a)
 \quad\text{for every coherent ideal }\mathfrak a.
\end{equation}
If $Y=s(v)$ and $Z=\overline{\{c(v)\}}$, then
\begin{equation}
\label{eq:trivial-endpoints}
                 v_{Y,\triv}\leq v\leq v_{Z,\triv}.
\end{equation}

The relation \eqref{eq:semivaluation-order} is reflexive and transitive by
inspection, and antisymmetric by
\Cref{thm:ch32-ideals-generate}\,\cref{itm:ch32-ig:separate}.  Its
usefulness comes from the following concrete description, which is the
trivially valued form of \cite[Lemma~4.4]{JM12}.

\begin{theorem}
\label{thm:ch32-order-centers}
\index{order!on semivaluations}
Let $X$ be a projective $k$-variety and $v,w\in X^{\An}$.  Then
$v\geq w$ in the sense of \eqref{eq:semivaluation-order} if and only if
\begin{equation}
\label{eq:ch32-order-criterion}
 c(v)\in\overline{\{c(w)\}}
 \qquad\text{and}\qquad
 v(f)\geq w(f)\ \text{ for all }f\in\cO_{X,c(v)} .
\end{equation}
Moreover \eqref{eq:trivial-endpoints} holds, and the order is preserved by
the scaling action of \Cref{sec:ch32-scaling}.
\end{theorem}

\begin{proof}
Observe first that the second condition in
\eqref{eq:ch32-order-criterion} makes sense once the first holds: if
$c(v)$ lies in the closure of $\{c(w)\}$, then $c(w)$ is a
generization of $c(v)$ and $\cO_{X,c(w)}$ is a localization of
$\cO_{X,c(v)}$; in particular both $v$ and $w$ are defined and
nonnegative on $\cO_{X,c(v)}$.

\emph{\eqref{eq:ch32-order-criterion} implies $v\geq w$.}  Let
$\mathfrak a$ be a coherent ideal and let $f_1,\ldots,f_r$ generate
$\mathfrak a_{c(v)}$.  Since $\cO_{X,c(w)}$ is a localization of
$\cO_{X,c(v)}$, the same elements generate $\mathfrak a_{c(w)}$.  Hence
\[
 v(\mathfrak a)=\min_iv(f_i)\geq\min_iw(f_i)=w(\mathfrak a).
\]

\emph{$v\geq w$ implies \eqref{eq:ch32-order-criterion}.}  Let
$Z_w:=\overline{\{c(w)\}}$ and let $\mathfrak a$ be its ideal sheaf.
Every local generator of $\mathfrak a$ at $c(w)$ lies in the maximal
ideal $\mathfrak m_{c(w)}$, so $w(\mathfrak a)>0$ by definition of the
center; hence $v(\mathfrak a)\geq w(\mathfrak a)>0$, and
\Cref{prop:ch32-tropical-identities}\,\cref{itm:ch32-ti:center} gives
$c(v)\in Z(\mathfrak a)=Z_w$, which is the first condition.

For the second, fix $f\in\cO_{X,c(v)}$ and let $L$ be an ample line
bundle.  By \Cref{lem:ch32-sections-generate-local-ring} there are
$m\geq1$ and $s,u\in H^0(X,L^{\otimes m})$ with $u(c(v))\neq0$ and
$f=s/u$ in $\cO_{X,c(v)}$.  Since $c(w)$ is a generization of
$c(v)$, the section $u$ also generates $L^{\otimes m}$ at $c(w)$, so
$w(u)=0$; and $v(u)=0$.  Therefore, using the coherent ideals
$\mathfrak a_s,\mathfrak a_u$ of \Cref{sec:ch32-sections},
\[
 v(f)=v(s)-v(u)=v(\mathfrak a_s)
      \geq w(\mathfrak a_s)=w(s)-w(u)=w(f).
\]

\emph{Proof of \eqref{eq:trivial-endpoints}.}  For the left inequality, let
$\mathfrak a$ be any coherent ideal.  If $Y\subset Z(\mathfrak a)$, then
$v(\mathfrak a)=+\infty=v_{Y,\triv}(\mathfrak a)$ by
\Cref{prop:ch32-tropical-identities}\,\cref{itm:ch32-ti:identities}, since
$s(v)=Y=s(v_{Y,\triv})$.  If $Y\not\subset Z(\mathfrak a)$, then
$v_{Y,\triv}(\mathfrak a)=0\leq v(\mathfrak a)$.  Either way
$v\geq v_{Y,\triv}$.  For the right inequality, note first that
$c(v_{Z,\triv})=c(v)$: the trivial semivaluation with support $Z$ has
center the generic point of $Z$, which is $c(v)$.  If
$c(v)\in Z(\mathfrak a)$, then
$Z\subset Z(\mathfrak a)$ and
$v_{Z,\triv}(\mathfrak a)=+\infty\geq v(\mathfrak a)$.  If
$c(v)\notin Z(\mathfrak a)$, then some local generator of
$\mathfrak a$ at $c(v)$ is a unit, so $v(\mathfrak a)=0$, while
$v_{Z,\triv}(\mathfrak a)=0$ as well.  Either way
$v\leq v_{Z,\triv}$.

\emph{Scaling.}  If $v\geq w$, then $tv(\mathfrak a)=t\,v(\mathfrak a)
\geq t\,w(\mathfrak a)=tw(\mathfrak a)$ for every $t>0$, where
$t\cdot(+\infty)=+\infty$; so $tv\geq tw$.
\end{proof}

\begin{remark}
\label{rem:ch32-endpoints-reading}
The two inequalities in \eqref{eq:trivial-endpoints} say that a
semivaluation is squeezed between the trivial semivaluation of its support
and the trivial semivaluation of the closure of its center.  Both bounds are
attained: the left one exactly when $v=v_{Y,\triv}$, the right one exactly
when $v=v_{Z,\triv}$.  In the language of
\Cref{sec:ch32-scaling}, the two bounds are the two endpoints of the
compactified scaling orbit through $v$, reached as $t\to0^+$ and
$t\to+\infty$ respectively; see \Cref{prop:ch32-scaling-limits}.  This is
the trivially valued reason why a \emph{decreasing} function --- one that
reverses \eqref{eq:semivaluation-order} --- automatically satisfies a
maximum principle, \Cref{prop:ch32-maximum-principle}, and it is used in that
form throughout \Cref{part:bj}.
\end{remark}

\subsection{Trivial semivaluations and the Zariski space}
\label{sec:ch32-trivial-locus}

Every irreducible subvariety $Y\subset X$ defines the trivial
semivaluation $v_{Y,\triv}$: it has support $Y$ and gives value zero to
every nonzero rational function on $Y$.  These points are exactly the fixed
points of scaling below.  The assignment $Y\mapsto v_{Y,\triv}$ identifies
the underlying set of the Zariski space $X$ with a subset
$X^{\triv}\subset X^{\An}$; the induced topologies are not the same.

\begin{theorem}
\label{thm:ch32-trivial-locus}
\index{trivial semivaluation}\index{Zariski space!inside the analytification}
Let $X$ be a projective $k$-variety.  Set
$X^{\triv}:=\{v_{Y,\triv}:Y\subset X\text{ integral closed}\}$.
\begin{enumerate}
\item\label{itm:ch32-tl:char}
$X^{\triv}=\{v\in X^{\An}:v(\mathfrak a)\in\{0,+\infty\}\text{ for every
coherent }\mathfrak a\}$, and this set is closed in $X^{\An}$, hence
compact.
\item\label{itm:ch32-tl:fixed}
$X^{\triv}$ is exactly the fixed-point set of the scaling action
\eqref{eq:valuation-scaling}.
\item\label{itm:ch32-tl:bijection}
The map $\xi\mapsto v_{\overline{\{\xi\}},\triv}$ is a bijection from the
underlying set of $X$ onto $X^{\triv}$, and it is a section of both
$s$ and $c$.
\item\label{itm:ch32-tl:closure}
$X^{\triv}=\overline{X(k)}$, the closure in $X^{\An}$ of the set of
trivial semivaluations at closed points.
\item\label{itm:ch32-tl:not-homeo}
The inverse map $X^{\triv}\to X$ is continuous, but it is not a
homeomorphism as soon as $\dim X\geq1$.
\end{enumerate}
\end{theorem}

\begin{proof}
\cref{itm:ch32-tl:char}  If $v=v_{Y,\triv}$ and $\mathfrak a$ is coherent,
then $v(\mathfrak a)=+\infty$ if $Y\subset Z(\mathfrak a)$ and
$v(\mathfrak a)=0$ otherwise, since every element of $\cO_{X,c(v)}$ has
value $0$ or $+\infty$.  Conversely suppose
$v(\mathfrak a)\in\{0,+\infty\}$ for all $\mathfrak a$.  Let
$Y=s(v)$; we must show that $v$ is trivial on $k(Y)^\times$.  Let
$f\in\cO_{X,c(v)}$ with $f\notin$ the support prime, so
$v(f)<+\infty$; by \Cref{lem:ch32-sections-generate-local-ring} write
$f=s/u$ with $u(c(v))\neq0$, whence $v(f)=v(\mathfrak a_s)\in\{0,\infty\}
$ and so $v(f)=0$.  Thus $v$ takes only the values $0$ and $+\infty$
on $\cO_{X,c(v)}$, and therefore $v=v_{Y,\triv}$.  Closedness: the set
$\{0,+\infty\}\subset[0,+\infty]$ is closed, so
$X^{\triv}=\bigcap_{\mathfrak a}(v\mapsto v(\mathfrak a))^{-1}
(\{0,+\infty\})$ is an intersection of closed sets by
\Cref{thm:ch32-ideals-generate}\,\cref{itm:ch32-ig:continuous}.

\cref{itm:ch32-tl:fixed}  A point $v$ is fixed by all $t>0$ if and only if
$t\,v(\mathfrak a)=v(\mathfrak a)$ for all $t>0$ and all
$\mathfrak a$, that is, if and only if
$v(\mathfrak a)\in\{0,+\infty\}$ for all $\mathfrak a$; now apply
\cref{itm:ch32-tl:char}.

\cref{itm:ch32-tl:bijection}  Distinct integral closed subschemes have
distinct trivial semivaluations, because their supports differ, so the map is
injective; it is surjective onto $X^{\triv}$ by definition.  Its
composition with $s$ is the identity of $X$, since
$s(v_{Y,\triv})=$ generic point of $Y$; and its composition with $c$ is
also the identity, since $v_{Y,\triv}$ is nonnegative on
$\cO_{X,\eta_Y}$ and strictly positive --- indeed $+\infty$ --- on its
maximal ideal, which is precisely the support prime of $v_{Y,\triv}$
(\Cref{cor:ch32-orbit-degenerate}); so $c(v_{Y,\triv})=\eta_Y$.

\cref{itm:ch32-tl:closure}  The inclusion $\overline{X(k)}\subset X^{\triv}$
follows from \cref{itm:ch32-tl:char}, $X^{\triv}$ being closed.  For the
converse, let $Y\subset X$ be integral closed and let
$W=W(\mathfrak a_1,\ldots,\mathfrak a_r)$ be a basic neighbourhood of
$v_{Y,\triv}$, that is, a finite intersection of sets
$\{v(\mathfrak a_i)<t_i\}$ and $\{v(\mathfrak a_j)>t_j\}$.  Partition the
indices according to whether $v_{Y,\triv}(\mathfrak a_i)=0$, which means
$Y\not\subset Z(\mathfrak a_i)$, or $=+\infty$, which means
$Y\subset Z(\mathfrak a_i)$.  The open set
\[
 V:=Y\setminus\bigcup_{i\,:\,Y\not\subset Z(\mathfrak a_i)}
 \bigl(Y\cap Z(\mathfrak a_i)\bigr)
\]
is a nonempty open subset of the integral scheme $Y$, hence contains a
closed point $p\in X(k)$, because $k$ is algebraically closed and $Y$
is of finite type.  For this $p$: if $Y\subset Z(\mathfrak a_i)$, then
$p\in Z(\mathfrak a_i)$, so $v_{p,\triv}(\mathfrak a_i)=+\infty
=v_{Y,\triv}(\mathfrak a_i)$; if $Y\not\subset Z(\mathfrak a_i)$, then
$p\notin Z(\mathfrak a_i)$, so
$v_{p,\triv}(\mathfrak a_i)=0=v_{Y,\triv}(\mathfrak a_i)$.  Hence
$v_{p,\triv}$ satisfies exactly the same strict inequalities as
$v_{Y,\triv}$ and lies in $W$.  Thus every neighbourhood of
$v_{Y,\triv}$ meets $X(k)$.

\cref{itm:ch32-tl:not-homeo}  The inverse is $s|_{X^{\triv}}$, continuous by
\Cref{thm:ch32-support-center-maps}\,\cref{itm:ch32-scm:s}.  It is not a
homeomorphism when $\dim X\geq1$: the space $X^{\triv}$ is compact
Hausdorff, being a closed subspace of $X^{\An}$, while the Zariski space
$X$ is not Hausdorff: a positive-dimensional irreducible component has a
generic point and a closed point that cannot be separated by disjoint
Zariski neighbourhoods.
\end{proof}

\begin{example}
\label{ex:ch32-trivial-locus-curve}
Let $X$ be a smooth connected projective curve over $k$.  Then
$X^{\triv}=\{v_{X,\triv}\}\sqcup\{v_{p,\triv}:p\in X(k)\}$.  By
\Cref{thm:ch32-trivial-locus}\,\cref{itm:ch32-tl:closure} the set
$\{v_{p,\triv}\}$ is dense in $X^{\triv}$, and each $v_{p,\triv}$ is
isolated in $X^{\triv}$: the ideal $\mathfrak m_p$ satisfies
$v_{p,\triv}(\mathfrak m_p)=+\infty$ while
$v_{q,\triv}(\mathfrak m_p)=0$ for $q\neq p$ and
$v_{X,\triv}(\mathfrak m_p)=0$, so $\{v(\mathfrak m_p)>1\}\cap X^{\triv}
=\{v_{p,\triv}\}$.  A compact Hausdorff space with a dense set of isolated
points whose complement is a single point is the one-point compactification
of that discrete set.  Hence
\[
 X^{\triv}\simeq X(k)^+\!,
\]
the one-point compactification of the infinite discrete set $X(k)$,
with $v_{X,\triv}$ as the point at infinity.  Compare the Zariski topology
on $X$, in which the generic point is dense and no closed point is open:
the two topologies are, in a precise sense, opposite.
\end{example}

\begin{remark}
\label{rem:ch32-specialization}
The two maps $s$ and $c$ and the order \eqref{eq:semivaluation-order}
fit together as follows.  Write $\rightsquigarrow$ for specialization in
$X$, so $\xi\rightsquigarrow\xi'$ means
$\xi'\in\overline{\{\xi\}}$.  Then, for every $v\in X^{\An}$,
\[
 s(v)\rightsquigarrow c(v),
\]
by \Cref{thm:ch32-support-center-maps}\,\cref{itm:ch32-scm:order}; and by
\Cref{thm:ch32-order-centers}, $v\geq w$ forces
$c(w)\rightsquigarrow c(v)$.  It also forces $s(v)\subseteq s(w)$:
testing \eqref{eq:semivaluation-order} on $\mathfrak a=\mathcal I_{s(w)}$,
for which $w(\mathfrak a)=+\infty$, gives $v(\mathfrak a)=+\infty$ and
hence $s(v)\subseteq Z(\mathfrak a)=s(w)$.  Thus increasing in the
analytic order specializes the center and can only shrink the support; it
does not leave the support fixed.  Both movements are visible in
\eqref{eq:trivial-endpoints}: passing from $v$ up to
$v_{Z,\triv}$ replaces the support $s(v)=Y$ by the smaller
$Z=\overline{\{c(v)\}}$.  On $\mathbf P^1$, for instance,
$v_{p,\triv}\geq\ord_p$ while
$s(\ord_p)=\mathbf P^1\supsetneq\{p\}=s(v_{p,\triv})$.  This is the exact
sense in which the analytic order refines the specialization order of the
scheme, and it is the reason the two endpoints in
\eqref{eq:trivial-endpoints} are the trivial semivaluations at the two ends
of the specialization $s(v)\rightsquigarrow c(v)$.
\end{remark}

\section{The scaling action, cones, and normalized links}
\label{sec:ch32-scaling}

\subsection{The action}
\label{sec:ch32-action}

For $t>0$, define $tv$ by $(tv)(f)=t\,v(f)$, with
$t(+\infty)=+\infty$.  This gives a continuous action
\begin{equation}
\label{eq:valuation-scaling}
              \mathbf R_{>0}\times X^{\An}\longrightarrow X^{\An}.
\end{equation}
The support is unchanged by scaling.  As $t\to0^+$, $tv$ converges to
$v_{s(v),\triv}$.  As $t\to+\infty$, it converges to the trivial
semivaluation supported on the closure of $c(v)$.

\begin{proposition}
\label{prop:ch32-scaling-action}
\index{scaling action!continuity}
Let $X$ be a projective $k$-variety.
\begin{enumerate}
\item\label{itm:ch32-sa:point}
For $v\in X^{\An}$ and $t>0$, the function $tv$ is again a point of
$X^{\An}$, with $s(tv)=s(v)$ and $c(tv)=c(v)$, and
$(tv)(\mathfrak a)=t\,v(\mathfrak a)$ for every coherent ideal.
\item\label{itm:ch32-sa:cont}
The map \eqref{eq:valuation-scaling} is a continuous action of the
topological group $\mathbf R_{>0}$.
\item\label{itm:ch32-sa:order}
The action preserves the order \eqref{eq:semivaluation-order}.
\item\label{itm:ch32-sa:functions}
The induced action on functions $\varphi:X^{\An}\to\Rbar$, defined by
\begin{equation}
\label{eq:ch32-scaling-on-functions}
 (t\cdot\varphi)(v):=t\,\varphi(t^{-1}v),
\end{equation}
is an action of $\mathbf R_{>0}$ preserving $C^0(X^{\An})$; a function is
fixed by every $t$ if and only if it is homogeneous, that is,
$\varphi(tv)=t\varphi(v)$ for all $t>0$ and $v$.
\end{enumerate}
\end{proposition}

\begin{proof}
\cref{itm:ch32-sa:point}  The axioms of a semivaluation are preserved by
multiplication by $t>0$: $tv(fg)=t(v(f)+v(g))=tv(f)+tv(g)$ and
$tv(f+g)\geq t\min\{v(f),v(g)\}=\min\{tv(f),tv(g)\}$, while
$tv(1)=0$ and $tv(0)=+\infty$.  The kernel is unchanged since
$t\cdot(+\infty)=+\infty$ and $ts<+\infty$ for $s<+\infty$; hence
$s(tv)=s(v)$.  The set $\{f:tv(f)>0\}=\{f:v(f)>0\}$ is unchanged, hence
$c(tv)=c(v)$.  Finally
$(tv)(\mathfrak a)=\min_i tv(f_i)=t\min_iv(f_i)=t\,v(\mathfrak a)$, the
generators at $c(tv)=c(v)$ being the same.

\cref{itm:ch32-sa:cont}  The group axioms $1\cdot v=v$ and
$t(t'v)=(tt')v$ are immediate.  For continuity, use the embedding
$\iota$ of \Cref{thm:ch32-ideals-generate}\,\cref{itm:ch32-ig:embed}: it
suffices to check that $(t,v)\mapsto(tv)(\mathfrak a)=t\,v(\mathfrak a)$ is
continuous $\mathbf R_{>0}\times X^{\An}\to[0,+\infty]$ for each
$\mathfrak a$.  This is the composition of the continuous map
$(t,v)\mapsto(t,v(\mathfrak a))$ with the multiplication map
$\mathbf R_{>0}\times[0,+\infty]\to[0,+\infty]$, $(t,\lambda)\mapsto
t\lambda$.  The latter is continuous, including at $\lambda=+\infty$:
if $t_i\to t>0$ and $\lambda_i\to+\infty$, then
$t_i\lambda_i\geq(t/2)\lambda_i\to+\infty$ eventually.

\cref{itm:ch32-sa:order}  Proved in \Cref{thm:ch32-order-centers}.

\cref{itm:ch32-sa:functions}  That \eqref{eq:ch32-scaling-on-functions}
defines an action is a computation:
\[
 \bigl(t\cdot(t'\cdot\varphi)\bigr)(v)
 =t(t'\cdot\varphi)(t^{-1}v)
 =t\,t'\varphi\bigl(t'^{-1}t^{-1}v\bigr)
 =\bigl((tt')\cdot\varphi\bigr)(v).
\]
Continuity is preserved because $v\mapsto t^{-1}v$ is a homeomorphism.
Finally $t\cdot\varphi=\varphi$ for all $t$ says
$t\varphi(t^{-1}v)=\varphi(v)$ for all $t,v$, which upon substituting
$v\mapsto tv$ is $\varphi(tv)=t\varphi(v)$.
\end{proof}

\begin{remark}
\label{rem:ch32-scaling-factor}
Two actions on functions suggest themselves: $\varphi\mapsto
\varphi\circ t^{-1}$ and the one in \eqref{eq:ch32-scaling-on-functions}
with the extra factor $t$.  It is the second that is used throughout
\Cref{part:bj}, and \Cref{def:homogeneous-function} is stated for it.  The
reason is that the objects being acted on are not functions but metrics:
scaling a valuation by $t$ scales the additive weight of a metric by $t$
as well, so the pairing between potentials and valuations is
$\mathbf R_{>0}$-equivariant only with the factor.  The geometric
counterpart --- base change along $\varpi\mapsto\varpi^t$ on the
test-configuration side --- is \Cref{prop:ch30-gauss-scaling} together with
\eqref{eq:ch30-gauss-lambda}; the identity
\eqref{eq:function-scaling-action} in \Cref{chap:bj-homogeneity-bdivisors} is
the same formula in the notation used there.
\end{remark}

\subsection{Orbits and their two endpoints}
\label{sec:ch32-orbits}

\begin{proposition}
\label{prop:ch32-scaling-limits}
\index{scaling action!limits}
Let $X$ be a projective $k$-variety, $v\in X^{\An}$, $Y=s(v)$ and
$Z=\overline{\{c(v)\}}$.  Then, in $X^{\An}$,
\[
 \lim_{t\to0^+}tv=v_{Y,\triv},
 \qquad
 \lim_{t\to+\infty}tv=v_{Z,\triv}.
\]
Consequently the map
\[
\begin{gathered}
 \iota_v:[0,+\infty]\longrightarrow X^{\An},\\
 \iota_v(0)=v_{Y,\triv},\quad
 \iota_v(t)=tv\ (0<t<\infty),\quad
 \iota_v(+\infty)=v_{Z,\triv},
\end{gathered}
\]
is continuous; and if $v\notin X^{\triv}$ it is a homeomorphism onto its
image, a \emph{compactified ray} through $v$.
\end{proposition}

\begin{proof}
By \Cref{thm:ch32-ideals-generate}\,\cref{itm:ch32-ig:embed} it suffices to
compute the limit of $(tv)(\mathfrak a)=t\,v(\mathfrak a)$ for every
coherent ideal $\mathfrak a$.  Three cases occur.

If $v(\mathfrak a)=0$, then $t\,v(\mathfrak a)=0$ for all $t$.  By
\Cref{prop:ch32-tropical-identities}\,\cref{itm:ch32-ti:center} this happens
exactly when $c(v)\notin Z(\mathfrak a)$, and then also
$Y\not\subset Z(\mathfrak a)$ --- as $c(v)\in Y$ --- so
$v_{Y,\triv}(\mathfrak a)=0$ and, since $Z\not\subset Z(\mathfrak a)$,
$v_{Z,\triv}(\mathfrak a)=0$.  Both limits agree with the constant value.

If $0<v(\mathfrak a)<+\infty$, then $t\,v(\mathfrak a)\to0$ as
$t\to0^+$ and $\to+\infty$ as $t\to+\infty$.  Here
$c(v)\in Z(\mathfrak a)$ but $Y\not\subset Z(\mathfrak a)$, so
$v_{Y,\triv}(\mathfrak a)=0$ and $v_{Z,\triv}(\mathfrak a)=+\infty$,
matching the two limits.

If $v(\mathfrak a)=+\infty$, then $t\,v(\mathfrak a)=+\infty$ for all
$t$; here $Y\subset Z(\mathfrak a)$, so $Z\subset Y\subset
Z(\mathfrak a)$ and both trivial semivaluations take the value $+\infty$.

Thus $\iota_v$ is continuous at $0$ and $+\infty$; it is continuous on
$(0,+\infty)$ by \Cref{prop:ch32-scaling-action}\,\cref{itm:ch32-sa:cont}.

Suppose now $v\notin X^{\triv}$, so that by
\Cref{thm:ch32-trivial-locus}\,\cref{itm:ch32-tl:char} there is a coherent
ideal $\mathfrak a$ with $0<v(\mathfrak a)<+\infty$.  Then
$t\mapsto t\,v(\mathfrak a)$ is injective on $(0,+\infty)$ with values in
$(0,+\infty)$, so $\iota_v$ is injective on $(0,+\infty)$; it takes
the value $0$ at $t=0$ and $+\infty$ at $t=+\infty$, so
$\iota_v$ is injective on all of $[0,+\infty]$.  A continuous injection
from a compact space to a Hausdorff space is a homeomorphism onto its image.
\end{proof}

\begin{corollary}
\label{cor:ch32-orbit-degenerate}
For $v\in X^{\An}$ the following are equivalent: $v\in X^{\triv}$; the
orbit $\{tv:t>0\}$ is a single point; $s(v)=\overline{\{c(v)\}}$.
\end{corollary}

\begin{proof}
If $v\in X^{\triv}$ the orbit is a point by
\Cref{thm:ch32-trivial-locus}\,\cref{itm:ch32-tl:fixed}, and
$s(v)=\overline{\{c(v)\}}$ by
\Cref{thm:ch32-trivial-locus}\,\cref{itm:ch32-tl:bijection}.  Conversely
suppose $Y:=s(v)=\overline{\{c(v)\}}=:Z$.  Then $c(v)$ is the generic
point of $Y$, so $\cO_{X,c(v)}$ has residue field $k(Y)$ and maximal
ideal the support prime of $v$; since $v>0$ exactly on the maximal ideal
and $v=+\infty$ exactly on the support prime, $v$ takes only the values
$0,+\infty$, i.e.\ $v=v_{Y,\triv}\in X^{\triv}$.  Finally if the orbit
is a single point then $v=2v$, so $v(\mathfrak a)\in\{0,+\infty\}$ for
every $\mathfrak a$ and $v\in X^{\triv}$.
\end{proof}

\subsection{Cones and links}
\label{sec:ch32-links}

\begin{definition}
\label{def:normalized-link}
Let $\Omega\subset X^{\An}$ be a scaling-stable set and let
$A:\Omega\to\mathbf R_{>0}$ be continuous and homogeneous:
$A(tv)=tA(v)$.  The normalized link associated to $A$ is
\begin{equation}
\label{eq:normalized-link}
 \operatorname{Link}_A(\Omega):=\{v\in\Omega\mid A(v)=1\}.
\end{equation}
\end{definition}

Every scaling orbit in $\Omega$ meets the link exactly once, at
$A(v)^{-1}v$.  There is deliberately no claim that all of $X^{\An}$ has
a canonical normalization.  For example, a nonzero ideal $\mathfrak a$
gives $A(v)=v(\mathfrak a)$ only where
$0<v(\mathfrak a)<+\infty$.  Points centered outside
$V(\mathfrak a)$ have gauge zero, while points supported inside it have
gauge $+\infty$.  Ignoring either boundary would erase proper-support
points needed later.

\begin{theorem}
\label{thm:ch32-link-cone}
\index{link!normalized}\index{cone!structure of a scaling-stable set}
Let $\Omega\subset X^{\An}$ be scaling stable and let
$A:\Omega\to\mathbf R_{>0}$ be continuous and homogeneous.  Put
$\Lambda:=\operatorname{Link}_A(\Omega)$.  Then:
\begin{enumerate}
\item\label{itm:ch32-lc:homeo}
the map
\[
 \mathbf R_{>0}\times\Lambda\longrightarrow\Omega,
 \qquad (t,\lambda)\longmapsto t\lambda,
\]
is a homeomorphism, with inverse
$v\mapsto\bigl(A(v),A(v)^{-1}v\bigr)$;
\item\label{itm:ch32-lc:functions}
restriction to $\Lambda$ is a bijection from the set of homogeneous
functions $\Omega\to\Rbar$ onto the set of all functions
$\Lambda\to\Rbar$; it restricts to a bijection between homogeneous
continuous real functions on $\Omega$ and continuous real functions on
$\Lambda$;
\item\label{itm:ch32-lc:notriv}
$\Omega\cap X^{\triv}=\varnothing$; in particular no gauge exists on a
scaling-stable set meeting $X^{\triv}$.
\end{enumerate}
\end{theorem}

\begin{proof}
\cref{itm:ch32-lc:notriv}  If $v\in\Omega\cap X^{\triv}$, then $2v=v$, so
$A(v)=A(2v)=2A(v)$, forcing $A(v)=0$, contrary to
$A>0$.

\cref{itm:ch32-lc:homeo}  The map $(t,\lambda)\mapsto t\lambda$ is
continuous by \Cref{prop:ch32-scaling-action}\,\cref{itm:ch32-sa:cont} and
lands in $\Omega$ by scaling stability.  The proposed inverse is
continuous, $A$ being continuous and the action being continuous.  They are
mutually inverse: $A(t\lambda)=tA(\lambda)=t$ and
$A(t\lambda)^{-1}(t\lambda)=\lambda$; and conversely
$A(v)\cdot\bigl(A(v)^{-1}v\bigr)=v$.  Note that
$A(v)^{-1}v\in\Lambda$ since $A(A(v)^{-1}v)=A(v)^{-1}A(v)=1$.

\cref{itm:ch32-lc:functions}  A homogeneous $\varphi$ satisfies
$\varphi(v)=A(v)\varphi\bigl(A(v)^{-1}v\bigr)$, so it is determined by
$\varphi|_\Lambda$; conversely any $\psi:\Lambda\to\Rbar$ extends by this
formula, and the extension is homogeneous.  Here the product
$A(v)\psi(\lambda)$ is unambiguous because $A(v)\in(0,+\infty)$; no
$0\cdot\infty$ arises, in accordance with the extended-real conventions of
the front matter.  Continuity transfers in both directions through the
homeomorphism of \cref{itm:ch32-lc:homeo}.
\end{proof}

\begin{example}
\label{ex:ch32-three-gauges}
Let $X$ be an integral projective $k$-variety.
\begin{enumerate}
\item\emph{An ideal gauge.}  Let $\mathfrak a\neq0$ be a coherent ideal and
put $A_{\mathfrak a}(v):=v(\mathfrak a)$.  It is continuous and homogeneous
on all of $X^{\An}$, but it is strictly positive and finite only on
\[
 \Omega_{\mathfrak a}
 :=\{v:c(v)\in Z(\mathfrak a)\}\setminus\{v:s(v)\subset Z(\mathfrak a)\},
\]
by \Cref{prop:ch32-tropical-identities}\,\cref{itm:ch32-ti:center} and
\cref{itm:ch32-ti:identities}.  The set $\Omega_{\mathfrak a}$ is scaling
stable, and $\operatorname{Link}_{A_{\mathfrak a}}(\Omega_{\mathfrak a})
=\{v:v(\mathfrak a)=1\}$.  For $\mathfrak a=\mathfrak m_p$ the ideal of a
closed point, this is the classical normalization
$v(\mathfrak m_p)=1$ of local valuation theory.
\item\emph{The log-discrepancy gauge.}  If $X$ is normal,
$\mathbf Q$-Gorenstein, and klt, the log discrepancy $A_X$ of
\eqref{eq:appJ-div-discrepancy} is homogeneous by
\eqref{eq:appJ-discrepancy-homogeneous} and takes values in
$[0,+\infty]$.  It is finite and positive on the quasimonomial locus of any
fixed SNC model, where it is even affine on each simplex
(\Cref{prop:ch32-log-discrepancy-affine}); on all of $X^{\An}$ it is only
lower semicontinuous, so \Cref{thm:ch32-link-cone} does not apply globally.
The link $\{A_X=1\}$ restricted to a fixed model is the compact
polyhedral complex one draws in K-stability.  No assertion about the union
over all models is made here.
\item\emph{A section gauge.}  For $s$ a regular section of an ample
$L^{\otimes m}$, the function $v\mapsto v(s)$ is a gauge on the
scaling-stable set where it is finite and positive, namely the set of points
whose center lies on $\operatorname{div}(s)$ and whose support does not.
This is the gauge implicitly used when one normalizes by
``$\varphi(v)=-1$ on the divisor''.
\end{enumerate}
In each case the gauge exists only after removing both the locus where it
vanishes and the locus where it is infinite.  Both removals delete points
that matter later: the first deletes $v_{X,\triv}$, the second deletes
proper-support points.
\end{example}

\begin{theorem}
\label{thm:ch32-no-compact-link}
\index{link!failure of global normalization}
Let $X$ be a nonempty projective $k$-variety.  There is no continuous
homogeneous function
\[
                 A:X^{\An}\longrightarrow\mathbf R_{>0}.
\]
More generally, no positive finite homogeneous gauge exists on a
scaling-stable subset of $X^{\An}$ that meets $X^{\triv}$.  Thus no
normalized link can parametrize all scaling orbits of $X^{\An}$.
\end{theorem}

\begin{proof}
If $v\in X^{\triv}$, then $2v=v$ by
\Cref{thm:ch32-trivial-locus}.  Homogeneity would therefore give
\[
                       A(v)=A(2v)=2A(v),
\]
contradicting $A(v)>0$.  The same argument applies to any scaling-stable
domain meeting $X^{\triv}$.
\end{proof}

\begin{remark}
\label{rem:ch32-what-link-parametrizes}
\Cref{thm:ch32-link-cone} says exactly what a link is good for: it
parametrizes the scaling orbits in $\Omega$, and homogeneous functions on
$\Omega$ correspond to arbitrary functions on it.  It does \emph{not}
parametrize the points of $X^{\An}$: the trivial locus and the
proper-support points are excluded by
\Cref{thm:ch32-link-cone}\,\cref{itm:ch32-lc:notriv} and by the domain
restrictions of \Cref{ex:ch32-three-gauges}.  A link on a fixed model may be
compact, but compactness for a larger scaling-stable domain is a separate
question and is not asserted here.  The global obstruction is instead that
no gauge includes the fixed locus $X^{\triv}$.  This is one reason
\Cref{part:bj} is a theory of homogeneous
functions on a cone paired with bounded functions on a compact space, related
by the Legendre correspondence of
\Cref{thm:homogeneous-legendre-decomposition}, rather than a theory of
functions on a compact link as in \Cref{part:bfj}.  The curve case, where the
link is a discrete set indexed by the closed points, is computed in
\Cref{prop:ch08-trivial-curve-scaling}\,(3).
\end{remark}

\begin{remark}
\label{rem:ch32-two-rescalings}
The scaling action \eqref{eq:valuation-scaling} multiplies a point of
$X^{\An}$ by a positive real number.  Over a discretely valued field there
is a superficially similar operation: replacing the ground absolute value
$|\cdot|$ by $|\cdot|^c$ multiplies every valuation by $c$, see
\eqref{eq:appJ-uniformizer}.  The two are different.  The second is a
renormalization of a fixed geometric situation and changes no point of the
space; the first moves a point of $X^{\An}$ along its ray.  Only the first
has fixed points, namely $X^{\triv}$.  Under the Gauss extension the first
is converted into the second, and this is the content of
\Cref{prop:ch30-gauss-scaling}; see also the warning following it.
\end{remark}

\section{The line, computed completely}
\label{sec:ch32-line}

Before the general theory of divisorial and quasimonomial points we compute
the smallest interesting example in full: every point of
$(\mathbf A^1_k)^{\An}$ and of $(\mathbf P^1_k)^{\An}$ is named, the
topology is described by an explicit neighbourhood basis, and the failure of
local compactness of $X^{\val}$ is proved.  Everything in this section is
elementary and self-contained.  The reader who has read
\Cref{sec:ch08-trivial-curve} will recognize
\Cref{prop:trivial-p1-tree} in \Cref{thm:ch32-p1-classification} below; the
present treatment adds the classification of the points of the affine line,
the neighbourhood basis, and the non-local-compactness.

\subsection{All points of the affine line}
\label{sec:ch32-affine-line}

\begin{theorem}
\label{thm:ch32-a1-classification}
\index{affine line!trivially valued analytification}
Let $k$ be algebraically closed with the trivial absolute value and let
$A=k[T]$.  Every point of $(\mathbf A^1_k)^{\An}$ is exactly one of the
following:
\begin{enumerate}
\item\label{itm:ch32-a1:triv}
the trivial valuation $v_{\triv}$ of $k(T)$, with
$v_\triv(f)=0$ for all $f\neq0$;
\item\label{itm:ch32-a1:ray}
$v_{a,t}:=t\ord_a$ for a unique closed point $a\in k$ and a unique
$t\in(0,+\infty)$, where $\ord_a(f)$ is the order of vanishing of $f$
at $a$;
\item\label{itm:ch32-a1:endpoint}
the proper-support semivaluation $v_{a,\infty}$ for a unique $a\in k$,
defined by $v_{a,\infty}(f)=+\infty$ if $f(a)=0$ and $=0$ otherwise;
\item\label{itm:ch32-a1:infty}
$v_{\infty,t}:=-t\deg$ for a unique $t\in(0,+\infty)$, that is,
$v_{\infty,t}(f)=-t\deg f$.
\end{enumerate}
In multiplicative notation, the points of types
\cref{itm:ch32-a1:triv}, \cref{itm:ch32-a1:ray}, \cref{itm:ch32-a1:infty} are
the seminorms
\[
 \Bigl|\sum_ic_i(T-a)^i\Bigr|_{a,r}=\max_{c_i\neq0}r^i,
 \qquad r\in(0,+\infty),
\]
with $r=\rme^{-t}<1$ giving \cref{itm:ch32-a1:ray}, $r=1$ giving
\cref{itm:ch32-a1:triv} independently of $a$, and $r=\rme^{t}>1$ giving
\cref{itm:ch32-a1:infty}; and \cref{itm:ch32-a1:endpoint} is the limit
$r\to0^+$.  The exponent changes sign between the two unbounded branches
because $t$ runs over $(0,+\infty)$ in both
\cref{itm:ch32-a1:ray} and \cref{itm:ch32-a1:infty}, while
$v_{a,t}(T-a)=+t$ and $v_{\infty,t}(T)=-t$.
\end{theorem}

\begin{proof}
Let $v$ be a point, i.e.\ a semivaluation on $A$ trivial on $k^\times$.

\emph{Case 1: $v$ has nonzero support.}  The support is a nonzero prime of
$k[T]$, hence $(T-a)$ for a unique $a\in k$, $k$ being algebraically
closed.  The residue field is $k[T]/(T-a)=k$, and the only valuation of
$k^\times$ trivial on $k^\times$ is the trivial one.  Thus $v$ is
determined: $v(f)=+\infty$ if $(T-a)\mid f$ and $v(f)=0$ otherwise,
which is \cref{itm:ch32-a1:endpoint}.

\emph{Case 2: $v$ has zero support}, so $v$ is a valuation of
$k(T)^\times$ trivial on $k^\times$.  Two subcases.

If $v\geq0$ on $k[T]$, then either $v\equiv0$ on
$k[T]\setminus\{0\}$, giving \cref{itm:ch32-a1:triv}; or
$\mathfrak q:=\{f:v(f)>0\}$ is a nonzero prime, hence $(T-a)$, and we set
$t:=v(T-a)>0$.  Every $f\in k[T]$ factors as
$f=c\prod_j(T-a_j)^{m_j}$, and $v(T-a_j)=0$ for $a_j\neq a$ since
$T-a_j\notin\mathfrak q$.  Hence $v(f)=t\,m$ with $m=\ord_a(f)$, which
is \cref{itm:ch32-a1:ray}.

If $v(g)<0$ for some $g\in k[T]$, write $g=c\prod_j(T-a_j)^{m_j}$; as
$v(c)=0$ and $v(g)=\sum_jm_jv(T-a_j)$, some factor has
$u:=v(T-a)<0$.  For any $b\in k$ with $b\neq a$ we have
$T-b=(T-a)+(a-b)$ with $v(T-a)=u<0=v(a-b)$, so
\Cref{lem:ch32-semivaluation-calculus}\,\cref{itm:ch32-sc:strict} gives
$v(T-b)=u$.  Thus $v(T-b)=u$ for \emph{every} $b\in k$, and
multiplicativity yields $v(f)=u\deg f$ for every nonzero
$f\in k[T]$.  Writing $t:=-u>0$ gives \cref{itm:ch32-a1:infty}.

\emph{Uniqueness of the labels.}  The four families are distinguished by
their supports and by the sign of $v(T-a)$: type
\cref{itm:ch32-a1:endpoint} is the only one with nonzero support; type
\cref{itm:ch32-a1:triv} is the only one with $v\equiv0$ on
$k[T]\setminus\{0\}$; type \cref{itm:ch32-a1:ray} has $v\geq0$ with
$v(T-a)>0$ for exactly one $a$; type \cref{itm:ch32-a1:infty} has
$v<0$ on all $T-b$.  Within each family the parameters are recovered as
$a=$ the unique zero with positive value, and $t=v(T-a)$ or
$t=-v(T)$ respectively.

\emph{The multiplicative formulas.}  For $f=\sum_ic_i(T-a)^i$ with
$c_i\in k$, triviality of the ground norm gives $|c_i|\in\{0,1\}$, and
the seminorm $|f|_{a,r}=\max_{c_i\neq0}r^i$ is multiplicative because the
graded ring associated with the $r$-weighted filtration on $k[T-a]$ is
$k[\overline{T-a}]$, a domain --- this is the graded-reduction argument of
\Cref{ex:graded-gauss}, or a direct check on leading terms.  For $r<1$ the
maximum is attained at the smallest index with $c_i\neq0$, giving
$-\log|f|_{a,r}=t\ord_a(f)$; for $r=1$ it is $1$; for $r>1$ it is
attained at the largest index, giving $-\log|f|_{a,r}=-t\deg f$,
independently of $a$.
\end{proof}

\begin{remark}
\label{rem:ch32-no-type-four}
Over a nontrivially valued complete field the classification of points of
the analytic affine line involves four types
(\Cref{thm:berk-a1-classification}).  Type~\textup{III} records a disc of
radius outside the divisible closure of the ground value group, whereas
type~\textup{IV} is represented by a decreasing family of discs with empty
intersection.  Over a trivially valued algebraically closed field,
\Cref{thm:ch32-a1-classification} gives the complete list directly: the
trivial point, one finite-radius ray for each $a\in k$, the corresponding
proper-support endpoints, and the ray centered at infinity.  There are no
additional type~\textup{IV} ends.  Thus the projective line is the compact
tree described in \Cref{thm:ch32-p1-classification}, with branches indexed
by $\mathbf P^1(k)$.
\end{remark}

\subsection{The projective line, its tree, and its topology}
\label{sec:ch32-p1}

\begin{theorem}
\label{thm:ch32-p1-classification}
\index{projective line!trivially valued analytification}
Let $X=\mathbf P^1_k$.  Then
\[
 X^{\An}=\{v_{\triv}\}\ \sqcup\
 \coprod_{p\in\mathbf P^1(k)}\iota_p\bigl((0,+\infty]\bigr),
\]
where for each closed point $p$ the map
$\iota_p:[0,+\infty]\to X^{\An}$ of \Cref{prop:ch32-scaling-limits}
attached to $v=\ord_p$ is a homeomorphism onto a compactified ray with
$\iota_p(0)=v_{\triv}$, $\iota_p(t)=t\ord_p$, and
$\iota_p(+\infty)=v_{p,\triv}$.  Moreover:
\begin{enumerate}
\item\label{itm:ch32-p1:partition}
the half-open rays $\iota_p((0,+\infty])=c^{-1}(p)$ form a partition of
$X^{\An}\setminus\{v_\triv\}$ into open sets;
\item\label{itm:ch32-p1:strata}
$X^{\val}=\{v_\triv\}\sqcup\coprod_p\iota_p((0,+\infty))$ and
$X^{\triv}=\{v_\triv\}\sqcup\{v_{p,\triv}:p\in\mathbf P^1(k)\}$;
\item\label{itm:ch32-p1:affine}
$(\mathbf A^1_k)^{\An}=X^{\An}\setminus\{v_{\infty,\triv}\}$, which is
\Cref{thm:ch32-a1-classification} with $\iota_\infty(t)=v_{\infty,t}$;
\item\label{itm:ch32-p1:basis}
the sets
\[
 N(p_1,t_1;\ldots;p_r,t_r)
 :=X^{\An}\setminus\bigcup_{j=1}^r\iota_{p_j}\bigl([t_j,+\infty]\bigr),
 \qquad p_j\in\mathbf P^1(k),\ t_j>0,
\]
form a neighbourhood basis of $v_\triv$.
\end{enumerate}
\end{theorem}

\begin{proof}
Every valuation of $k(T)$ trivial on $k$ and different from
$v_\triv$ is a positive multiple of $\ord_p$ for a unique closed point
$p\in\mathbf P^1(k)$: on $\mathbf A^1$ this is
\Cref{thm:ch32-a1-classification}, cases \cref{itm:ch32-a1:ray} and
\cref{itm:ch32-a1:infty}, and $-t\deg=t\ord_\infty$ in the coordinate
$1/T$.  Every semivaluation with nonzero support is $v_{p,\triv}$ for a
unique closed point, by
\Cref{prop:ch32-support-partition}\,\cref{itm:ch32-sp:triv}.  This is the
displayed decomposition, and $\iota_p$ is a homeomorphism onto a
compactified ray by \Cref{prop:ch32-scaling-limits}, whose endpoints are
$v_{\mathbf P^1,\triv}=v_\triv$ at $t=0$ and $v_{p,\triv}$ at
$t=+\infty$ since $c(\ord_p)=p$.

\cref{itm:ch32-p1:partition}  For $t>0$ the center of $t\ord_p$ is $p$
and the center of $v_{p,\triv}$ is $p$, while $c(v_\triv)$ is the
generic point.  Hence $\iota_p((0,+\infty])=c^{-1}(p)$, and these sets are
open by \Cref{thm:ch32-support-center-maps}\,\cref{itm:ch32-scm:c}, $\{p\}$
being closed.

\cref{itm:ch32-p1:strata} and \cref{itm:ch32-p1:affine} are immediate from the
classification, using $U^{\An}=\{v:s(v)\cap U\neq\varnothing\}$ with
$U=\mathbf A^1$: the only point whose support misses $\mathbf A^1$ is
$v_{\infty,\triv}$.

\cref{itm:ch32-p1:basis}  A nonzero coherent ideal on $\mathbf P^1$ is
$\cO(-D)$ for an effective divisor $D=\sum_jm_jp_j$ --- being an ideal
sheaf on a smooth curve --- so
$v(\mathfrak a)=\sum_j m_j\,v(\mathfrak m_{p_j})$, and
$v(\mathfrak m_{p})=t$ if $v=\iota_{p}(t)$ and $=0$ if
$c(v)\neq p$.  Hence
\[
 \{v:v(\mathfrak a)<\varepsilon\}
 =X^{\An}\setminus\bigcup_j\iota_{p_j}\bigl([\varepsilon/m_j,+\infty]\bigr),
\]
which is of the stated form; conversely
$N(p_1,t_1;\ldots;p_r,t_r)$ is the finite intersection of the sets
$\{v:v(\mathfrak m_{p_j})<t_j\}$, each of which is one of the basic
conditions in \Cref{lem:ch32-neighbourhood-basis-triv}.  Now apply
\Cref{lem:ch32-neighbourhood-basis-triv}.
\end{proof}

\begin{example}
\label{ex:triv-p1-tree}
Let $X=\mathbf P^1_k$.  For each closed point $p$, the valuations
$t\ord_p$, $t\geq0$, form a ray issuing from $v_{X,\triv}$.  Its other
endpoint is $v_{p,\triv}$, reached as $t\to+\infty$.  The union of these
rays gives a compact star-shaped tree.  If $X$ is nodal, the
analytification is obtained from the tree of the normalization by identifying
the two proper-support endpoints over each node.  This foreshadows why
unibranchness is forced by the envelope property.
\end{example}

The assertion about a nodal curve deserves a proof, since it is the source
of the node/cusp comparison of \Cref{ex:bj-node-cusp}.

\begin{proposition}
\label{prop:ch32-nodal-curve}
\index{nodal curve!analytification}
Let $X$ be a connected projective curve over $k$ and
$\nu:X^\nu\to X$ its normalization.  Then
$\nu^{\An}:(X^{\nu})^{\An}\to X^{\An}$ is a continuous surjection which is
injective except over the non-normal points, and for $p\in X(k)$ with
$\nu^{-1}(p)=\{q_1,\ldots,q_m\}$ it identifies the $m$ endpoints
$v_{q_i,\triv}$ to the single endpoint $v_{p,\triv}$.  In particular, if
$X$ is nodal, $X^{\An}$ is obtained from the disjoint union of the trees
of the components of $X^\nu$ by gluing, at each node, the two
proper-support endpoints lying over it.
\end{proposition}

\begin{proof}
The morphism $\nu$ is finite and surjective, hence
$\nu^{\An}$ is proper by \Cref{thm:berk-proper-compact}.  It is
surjective as well: for $x\in X^{\An}$, base change to the completed residue
field $\mathcal H(x)$ gives a nonempty finite fibre; a maximal ideal of its
coordinate algebra yields a finite extension of $\mathcal H(x)$, whose
extended absolute value defines a point over $x$.  Thus $\nu^{\An}$ is a
surjective map of compact spaces.  Since $\nu$ is an isomorphism over
$X\setminus\Sing(X)$, the map $\nu^{\An}$ is bijective over
$\{v:c(v)\notin\Sing(X)\}=c^{-1}(X\setminus\Sing X)$, which is
\emph{closed}, indeed compact: the centre map is anticontinuous
(\Cref{thm:ch32-center-compact}\,\cref{itm:ch32-cc:general}), so the preimage
of the open set $X\setminus\Sing X$ is closed.  It is the complement of the
finitely many \emph{open} half-open rays $c^{-1}(p)$, $p\in\Sing X$.

Fix $p\in\Sing(X)$ with $\nu^{-1}(p)=\{q_1,\ldots,q_m\}$, and consider
the points of $(X^\nu)^{\An}$ with center some $q_i$.  Such a point is
either $t\ord_{q_i}$ with $t\in(0,+\infty)$, or the endpoint
$v_{q_i,\triv}$.  Indeed, each component of $X^\nu$ is a smooth
projective curve, and every nontrivial valuation of its function field
centered at $q_i$ is a positive multiple of $\ord_{q_i}$; this is
\Cref{prop:ch08-trivial-curve-scaling}\,(1), whose one-variable proof uses
that the local ring at $q_i$ is a discrete valuation ring.

If $w=t\ord_{q_i}$, its image $\nu^{\An}(w)$ is the valuation
$f\mapsto t\ord_{q_i}(f\circ\nu)$ of $k(X_{(i)})=k(X^\nu_{(i)})$, the
function fields agreeing because $\nu$ is birational.  Two such images
coincide only if the underlying valuations of the same function field
coincide; but a valuation determines its center on the proper scheme
$X^\nu$, so distinct $q_i$ on one component give distinct valuations, and
distinct components give valuations of distinct function fields.  Hence
$\nu^{\An}$ is injective on the union of the open rays over the $q_i$.

If $w=v_{q_i,\triv}$, then $\nu^{\An}(w)$ sends $f$ to
$v_{q_i,\triv}(f\circ\nu)$, which is $+\infty$ exactly when $f$
vanishes at $p$ and $0$ otherwise; that is, $\nu^{\An}(w)=v_{p,\triv}$,
independently of $i$.  Hence exactly the $m$ endpoints over $p$ are
identified.

For $X$ nodal, $m=2$ at each node and $X^\nu$ is smooth, so its
analytification is the disjoint union of the star-shaped trees of
\Cref{prop:ch08-trivial-curve-scaling}; the description follows.
\end{proof}

\subsection{Why the valuation space alone is not enough}
\label{sec:ch32-not-locally-compact}

\begin{proposition}
\label{prop:ch32-xval-not-locally-compact}
\index{valuation space!not locally compact}
Let $X$ be a smooth connected projective curve over an algebraically
closed $k$.  Then $X^{\val}=X^{\An}\setminus\{v_{p,\triv}:p\in X(k)\}$
is dense in $X^{\An}$ but is not locally compact at $v_\triv$.  In
particular $X^{\val}$ is not closed and not locally closed in
$X^{\An}$.
\end{proposition}

\begin{proof}
Density: by \Cref{prop:ch08-trivial-curve-scaling}, each
$v_{p,\triv}=\lim_{t\to+\infty}t\ord_p$ is a limit of points of
$X^{\val}$.

Suppose $N\subset X^{\val}$ were a compact neighbourhood of $v_\triv$ in
$X^{\val}$.  Then $N\supset N'\cap X^{\val}$ for some open
$N'\subset X^{\An}$ containing $v_\triv$.  The proof of
\Cref{thm:ch32-p1-classification}\,\cref{itm:ch32-p1:basis} uses only that a
coherent ideal on a smooth curve is the ideal of an effective divisor, so it
gives the same neighbourhood basis on $X$; we may therefore take
$N'=N(p_1,t_1;\ldots;p_r,t_r)$.  Choose $p\in X(k)$ distinct from
$p_1,\ldots,p_r$, which is possible because $X(k)$ is infinite.  Then
$N'$ contains the entire half-open ray $\iota_p((0,+\infty])$, so
$N\supset\iota_p((0,+\infty))$.

Consider the net $\bigl(\iota_p(t)\bigr)_{t\in(0,+\infty)}$, directed by
increasing $t$; it lies in $N$.  If $N$ were compact, the net would
have a cluster point in $N\subset X^{\val}$.  But in the compact
Hausdorff space $X^{\An}$ the net converges to $v_{p,\triv}$ by
\Cref{prop:ch32-scaling-limits}, so its unique cluster point in $X^{\An}$
is $v_{p,\triv}$, which is not in $X^{\val}$.  Contradiction.  For the
final assertion, a locally closed subspace of a compact Hausdorff space is
locally compact.
\end{proof}

\begin{remark}
\label{rem:ch32-moral-of-the-line}
\Cref{prop:ch32-xval-not-locally-compact} is the precise reason the theory
is set up on $X^{\An}$ rather than on the space $X^{\val}$ of honest
valuations, even though every divisorial and quasimonomial point lies in
$X^{\val}$.  The compactification is achieved exactly by the
proper-support points, and it costs the admission of the value $-\infty$
for potentials.  \Cref{sec:ch32-proper-support} pays that cost explicitly.
The alternative --- working on $X^{\val}$ with its own topology --- is
sometimes convenient locally, and is the setting of
\Cref{chap:local-valuative-psh}; but it is not available globally, because
the space does not provide the compact global domain used for weak
compactness and energy in \Cref{chap:bj-energy}.
\end{remark}

\section{Divisorial points}
\label{sec:ch32-divisorial}

\subsection{Definition and first invariants}
\label{sec:ch32-divisorial-def}

Assume now that $X$ is integral.  Let $\pi:Y\to X$ be a normal
projective birational model and $E\subset Y$ a prime divisor.  For
$t\in\mathbf Q_{>0}$, the valuation $t\ord_E$ is called divisorial.  We
also include $v_{X,\triv}$ as the limiting point at $t=0$, and write
$X^{\divis}$ for the resulting set.

\begin{definition}
\label{def:ch32-divisorial}
\index{divisorial valuation!over a trivially valued field}
Let $X$ be an integral projective $k$-variety.  A valuation
$v\in X^{\val}$ is \emph{divisorial} if either $v=v_{X,\triv}$ or
$v=t\ord_E$ with $t\in\mathbf Q_{>0}$ and $E$ a prime divisor on some
normal projective birational model $\pi:Y\to X$, where
$\ord_E:k(X)^\times=k(Y)^\times\to\mathbf Z$ is the order of vanishing
along $E$.  We write $X^{\divis}\subset X^{\val}$ for the set of
divisorial valuations.  For a possibly reducible $X$ we set
$X^{\divis}:=\coprod_\alpha X_\alpha^{\divis}$.
\end{definition}

Three comments on the definition are in order, and each is used later.

\emph{Why the trivial valuation is included.}  With $t=0$ the formula
$t\ord_E$ gives $v_{X,\triv}$, and by
\Cref{prop:ch32-scaling-limits} it is the limit of $t\ord_E$ as
$t\to0^+$.  Excluding it would make $X^{\divis}$ fail to be closed under
the limits that occur in \Cref{thm:triv-divisorial-density}, and would break
the statement $X^{\divis}=\{v_E\}$ for components $E$ of central fibres
of test configurations, since the trivial test configuration contributes the
trivial valuation.  This convention is that of
\cite[\S1.1.1]{BJ22} and is used without further comment.

\emph{Why the normalization is not fixed.}  Over a discretely valued field
one normalizes a divisorial point by $v(\varpi)=1$, that is,
$x_E=b_E^{-1}\ord_E$ with $b_E=\ord_E(\varpi)$; see
\Cref{rem:divisorial-normalization}.  Over a trivially valued field no such
normalization exists and the whole ray $\mathbf Q_{>0}\ord_E$ consists of
divisorial points.  This is the source of the ``factor of $b_E$''
discrepancies warned against in \Cref{app:normalizations}, and it is the
reason the Gauss extension of \Cref{sec:ch32-gauss} converts a ray into a
family of normalizations rather than into a single point.

\emph{Why only rational multiples.}  For $t$ irrational, $t\ord_E$ is a
perfectly good point of $X^{\An}$, of rational rank one but with value
group $t\mathbf Z\not\subset\mathbf Q$.  It is not divisorial: no prime
divisor on a birational model has $\ord$ equal to it, because the value
group of $\ord_F$ is $\mathbf Z$ and a rational multiple of it is
$\mathbf Q$-proportional to $\mathbf Z$.  Such points are quasimonomial
of rational rank one, and they occur; see
\Cref{ex:ch32-blowup-surface}\,\cref{itm:ch32-bs:irrational}.

\begin{proposition}
\label{prop:ch32-divisorial-invariants}
\index{divisorial valuation!invariants}
Let $X$ be an integral projective $k$-variety of dimension $n$, let
$Y\to X$ be a normal projective birational model, $E\subset Y$ a prime
divisor, and $t\in\mathbf R_{>0}$.  Put $v=t\ord_E$.  Then:
\begin{enumerate}
\item\label{itm:ch32-di:group}
$v$ is a valuation of $k(X)^\times$ trivial on $k^\times$, with value
group $t\mathbf Z$, hence of rational rank one and rank one;
\item\label{itm:ch32-di:residue}
its residue field is $k(E)$, of transcendence degree $n-1$ over $k$;
\item\label{itm:ch32-di:abhyankar}
$v$ is Abhyankar: equality holds in \Cref{thm:abhyankar-inequality},
\[
 \operatorname{rat.rk}(v)+\operatorname{trdeg}(v)=1+(n-1)=n=\dim X;
\]
\item\label{itm:ch32-di:center}
the center of $v$ on $X$ is the generic point of $\pi(E)$, and
$s(v)=X$; in particular $v\in X^{\val}$;
\item\label{itm:ch32-di:conversely}
conversely, a nontrivial $v\in X^{\val}$ with
$\operatorname{rat.rk}(v)=1$ and $\operatorname{trdeg}(v)=n-1$ is a
positive real multiple of some $\ord_E$, and is divisorial precisely when
that multiple is rational.
\end{enumerate}
\end{proposition}

\begin{proof}
\cref{itm:ch32-di:group} and \cref{itm:ch32-di:residue}  The local ring
$\cO_{Y,\eta_E}$ at the generic point of $E$ is a discrete valuation
ring, $Y$ being normal and $E$ of codimension one; its valuation is
$\ord_E$, with value group $\mathbf Z$ and residue field $k(E)$.
Multiplying by $t$ rescales the value group to $t\mathbf Z$ without
changing the valuation ring, hence without changing the residue field.  The
rational rank of $t\mathbf Z$ is one, and its rank as an ordered group is
one since it embeds in $\mathbf R$
(\Cref{prop:rank-one-real-embedding}).  Since $E$ is a prime divisor on
the $n$-dimensional $Y$, $\operatorname{trdeg}_kk(E)=n-1$.

\cref{itm:ch32-di:abhyankar}  Add the two numbers and compare with
\Cref{thm:abhyankar-inequality}, which asserts
$\operatorname{rat.rk}+\operatorname{trdeg}\leq n$; here equality holds.
See also \Cref{cor:rank-plus-trdeg}.

\cref{itm:ch32-di:center}  The support is $X$ because $v$ is a valuation
of the function field, hence has zero kernel on $\cO_X$.  The center is
the unique point $\xi\in X$ with $\cO_{X,\xi}$ dominated by the valuation
ring; since $\pi$ is proper and $\ord_E$ is centered at $\eta_E$ on
$Y$, the center on $X$ is $\pi(\eta_E)$, the generic point of
$\pi(E)$.

\cref{itm:ch32-di:conversely}  This is Zariski's classification of Abhyankar
valuations of rational rank one.  The conclusion that the value group is
discrete is part of the Zariski--Abhyankar theorem, not a consequence of the
numerical inequality alone.  In the present rank-one equality case the
theorem gives a generator $t>0$ and a prime divisor $E$ on a suitable
normal projective birational model such that $v=t\ord_E$; see
\cite[Ch.~VI, \S14]{ZariskiSamuelII}.  We use only
the direct implications \cref{itm:ch32-di:group}--\cref{itm:ch32-di:center}
below; \cref{itm:ch32-di:conversely} is quoted for orientation and is not
used in any proof of this book.
\end{proof}

\begin{proposition}
\label{prop:ch32-divisorial-linear-growth}
\index{linear growth!of divisorial valuations}
Let $X$ be a projective $k$-variety and $L$ an ample line bundle on
$X$.  For every $v\in X^{\divis}$ there is a constant $C=C(v,L)$ such
that
\[
 v(s)\leq Cm
 \qquad\text{for every }m\geq1\text{ and every regular }
 s\in H^0(X,L^{\otimes m}).
\]
In the notation of \Cref{def:bj-linear-growth}, $T_L(v)<+\infty$; thus
$X^{\divis}\subset X^{\lin}$.
\end{proposition}

\begin{proof}
We may assume $X$ integral, the general case following componentwise by
\Cref{def:ch32-divisorial}.  The case $v=v_{X,\triv}$ is trivial, with
$C=0$.  So let $v=t\ord_E$ with $t\in\mathbf Q_{>0}$, $E$ a prime
divisor on a normal projective birational model $\mu:Y\to X$ of dimension
$n$, and fix an ample line bundle $H$ on $Y$.

Let $s\in H^0(X,L^{\otimes m})$ be a regular section, that is, a nonzero
one, $X$ being integral, and put $a:=\ord_E(s)=\ord_E(\mu^*s)$.  The
divisor $\operatorname{div}(\mu^*s)$ is an effective Weil divisor on $Y$
in the class of $\mu^*(mL)$, and it contains $aE$; hence
$\operatorname{div}(\mu^*s)-aE$ is effective, and intersecting with the nef
class $H^{n-1}$ gives
\[
 a\,\bigl(E\cdot H^{n-1}\bigr)
 \leq\bigl(\operatorname{div}(\mu^*s)\cdot H^{n-1}\bigr)
 =m\,\bigl(\mu^*L\cdot H^{n-1}\bigr).
\]
Here $(E\cdot H^{n-1})>0$ because $H$ is ample and $E$ is an effective
nonzero divisor, and the middle inequality is the nonnegativity of the
intersection of an effective cycle with an ample class.  Therefore
\[
 \ord_E(s)=a\leq m\cdot
 \frac{(\mu^*L\cdot H^{n-1})}{(E\cdot H^{n-1})}=:mC_0 ,
\]
and $v(s)=t\,\ord_E(s)\leq tC_0m$.  Take $C=tC_0$.
\end{proof}

\begin{remark}
\label{rem:ch32-linear-growth-converse}
Linear growth is weaker than divisoriality.  Indeed, the same intersection
estimate proves linear growth for every real multiple $t\ord_E$, whereas
our divisorial convention retains only $t\in\mathbf Q_{>0}$.  Thus the
inclusion $X^{\divis}\subset X^{\lin}$ is strict as soon as $X$ has a
prime divisor; irrational quasimonomial points give further examples on
higher-dimensional cones.  The inclusion $X^{\lin}\subset X^{\val}$ is
strict in general: \Cref{prop:ch32-valuation-criterion} characterizes linear
growth by finite maximal vanishing order, while valuations with infinite
maximal vanishing order do occur.  The later theory identifies
$X^{\lin}$ with the nonpluripolar points of $X^{\An}$; see
\Cref{chap:bj-linear-growth} and \Cref{thm:bj-dirac-finite-energy}.
\end{remark}

\subsection{Density of divisorial points}
\label{sec:ch32-density}

\begin{theorem}\label{thm:triv-divisorial-density}
For an integral projective variety, $X^{\divis}$ is dense in $X^{\An}$.
Consequently, two continuous functions agreeing on $X^{\divis}$ agree on
all of $X^{\An}$.  Later, the BJ comparison theorem will show that two psh
functions agreeing on $X^{\divis}$ agree everywhere, even when they can be
$-\infty$ at proper-support points.
\end{theorem}

\begin{proof}
We reproduce the argument of \cite[Theorem~2.14]{BJ22}.  First recall two
facts proved there in Theorems~2.2 and~2.7 and Lemma~2.8.  The rational PL
functions form a vector lattice in $C^0(X^{\An})$, contain the rational
constants, and separate points.  Indeed, the topology of $X^{\An}$ is
generated by the ideal evaluations $\log|\mathfrak a|$, and the bounded
truncation
\[
             \max\{\log|\mathfrak a|,-m\}
\]
is the PL function of the flag ideal
$\mathfrak a+(\varpi^m)$.  The lattice form of the
Stone--Weierstrass theorem therefore makes $\PL(X)$ uniformly dense in
$C^0(X^{\An})$.

Next let $\phi\in\PL(X)$ vanish on $X^{\divis}$.  After passing to a
single integrally closed test configuration $\cX$, the test-configuration
description of PL functions writes $\phi=\phi_D$ for a vertical
$\mathbf Q$-Cartier divisor $D$ on $\cX$.  If $E$ ranges over the
irreducible components of $\cX_0$, the associated normalized valuations
$v_E$ belong to $X^{\divis}$; hence
$\phi_D(v_E)=0$ for every $E$.  The valuative effectivity criterion says
that $D\geq0$ exactly when all $\phi_D(v_E)\geq0$.  Applied to both
$D$ and $-D$, it gives $D=0$, and therefore $\phi=0$.

It remains to pass from PL functions to the topology.  If the closure
$F:=\overline{X^{\divis}}$ were proper, compact Hausdorff normality would
give a nonzero continuous function $h$ vanishing on $F$.  Choose PL
functions $\phi_j$ converging uniformly to $h$.  If
$\phi_j=\phi_D$ on a determination $\cX$, put
$c^+=\max_E\phi_D(v_E)$ and
$c^-=\min_E\phi_D(v_E)$.  Since $\phi_{\cX_0}=1$, the effectivity
criterion applied to $c^+\cX_0-D$ and
$D-c^-\cX_0$ gives
$c^-\leq\phi_D\leq c^+$ on all of $X^{\An}$.  Thus
\[
       \|\phi_j\|_{X^{\An}}
       =\sup_{v\in X^{\divis}}|\phi_j(v)|.
\]
Since $h$ vanishes on $X^{\divis}$, the right-hand side tends to zero,
whereas $\phi_j\to h$ uniformly; thus $h=0$, a contradiction.  Hence
$F=X^{\An}$.  The assertion about continuous functions follows
immediately from density.
\end{proof}

Here every nontrivial member of $X^{\divis}$ has full support.  Density
nevertheless holds in the \emph{whole} analytification: a net of full-support
divisorial valuations may converge to a semivaluation with a nonzero kernel.
Thus the theorem does not silently enlarge the term divisorial to include
divisors over proper subvarieties.

We isolate the two ingredients imported in the proof, so that the reader can
check them against \Cref{chap:bj-fs-pl-psh} and
\Cref{chap:bj-test-configurations} when those chapters are read.

\begin{remark}
\label{rem:ch32-density-ingredients}
The proof of \Cref{thm:triv-divisorial-density} uses exactly two statements
from the next two chapters, and no others.
\begin{enumerate}
\item\emph{PL functions are uniformly dense in $C^0(X^{\An})$.}
This is \eqref{eq:pl-stone-weierstrass} in \Cref{chap:bj-fs-pl-psh}, quoted
from \cite[Theorem~2.2]{BJ22}.  Its proof is the lattice Stone--Weierstrass
theorem applied to the algebra generated by the truncated ideal evaluations
$\max\{-v(\mathfrak a),-m\}$, which are PL by
\Cref{thm:tc-pl-correspondence}; the separation hypothesis is
\Cref{thm:ch32-ideals-generate}\,\cref{itm:ch32-ig:separate}, proved in the
present chapter.
\item\emph{A PL function is $\phi_D$ for a vertical
$\mathbf Q$-Cartier divisor $D$ on an integrally closed test
configuration, and $D\geq0$ if and only if $\phi_D(v_E)\geq0$ for every
component $E$ of the central fibre.}
This is \Cref{thm:tc-pl-correspondence} together with the valuative
effectivity criterion, \cite[Theorem~2.7 and Lemma~2.8]{BJ22}; the latter is
the vertical Serre criterion of \Cref{def:integrally-closed-tc}, in the form
``a vertical fractional ideal is contained in $\cO_{\cX}$ if and only if
$\ord_E\geq0$ on it for every component $E$ of $\cX_0$''.
\end{enumerate}
No circularity arises: neither statement uses
\Cref{thm:triv-divisorial-density}.
\end{remark}

\begin{corollary}
\label{cor:ch32-determination-divisorial}
Let $X$ be an integral projective $k$-variety.
\begin{enumerate}
\item\label{itm:ch32-dd:cont}
Two continuous functions on $X^{\An}$ agreeing on $X^{\divis}$ are
equal.
\item\label{itm:ch32-dd:usc}
If $\varphi,\psi:X^{\An}\to[-\infty,+\infty)$ are upper semicontinuous
and $\varphi\leq\psi$ on $X^{\divis}$, it does \emph{not} follow that
$\varphi\leq\psi$ on $X^{\An}$; an additional structural hypothesis is
required, and it is a hypothesis on $\varphi$, not on $\psi$: it suffices
that $\varphi$ be the usc regularization of its restriction to
$X^{\divis}$, as a $\theta$-psh function is by
\Cref{thm:bj-basic-closure}\,\cref{itm:bj-closure:determination}, for then
$\varphi(v)=\limsup_{X^{\divis}\ni w\to v}\varphi(w)
\leq\limsup_{w\to v}\psi(w)\leq\psi(v)$.  Asking it of $\psi$ would not
help: in the counterexample below $\psi\equiv0$ already equals the usc
regularization of its own restriction to $X^{\divis}$.
\item\label{itm:ch32-dd:sup}
If $\varphi$ is \emph{lower} semicontinuous and $\psi$ is \emph{upper}
semicontinuous, and $\varphi\leq\psi$ on $X^{\divis}$, then
$\varphi\leq\psi$ on $X^{\An}$.  The pairing is forced: with $\varphi$
usc and $\psi$ lsc the assertion is false, and the preceding clause is the
counterexample --- take $\psi\equiv0$ and
$\varphi=\mathbf 1_{\{v_0\}}$ for $v_0\notin X^{\divis}$.
\end{enumerate}
\end{corollary}

\begin{proof}
\cref{itm:ch32-dd:cont} is \Cref{thm:triv-divisorial-density}.

\cref{itm:ch32-dd:usc}  Take $\psi\equiv0$ and let $\varphi$ be the
indicator-type function $\varphi(v_0)=1$, $\varphi(v)=0$ for
$v\neq v_0$, where $v_0$ is any point not in $X^{\divis}$ --- for
instance a proper-support point, which exists as soon as $\dim X\geq1$.
This $\varphi$ is upper semicontinuous, since each superlevel set
$\{\varphi\geq c\}$ is $\varnothing$, $\{v_0\}$, or $X^{\An}$, and
$\{v_0\}$ is closed; it agrees with $\psi$ on $X^{\divis}$ and not
everywhere.  The example makes the point that upper semicontinuity controls
a function from above near a point but permits it to jump up \emph{at} the
point relative to the values on a prescribed dense subset.  This is why
\Cref{part:bj} always couples density statements with a regularization
theorem.

\cref{itm:ch32-dd:sup}  Suppose $\varphi(v)>\psi(v)$ for some
$v\in X^{\An}$ and choose $c\in\mathbf R$ with
$\psi(v)<c<\varphi(v)$.  The set
\[
 U:=\{\varphi>c\}\cap\{\psi<c\}
\]
is open --- the first factor because $\{\varphi\leq c\}$ is closed by lower
semicontinuity, the second because $\{\psi\geq c\}$ is closed by upper
semicontinuity --- and it contains $v$, hence is nonempty.  By
\Cref{thm:triv-divisorial-density} it contains a point
$w\in X^{\divis}$, and there $\varphi(w)>c>\psi(w)$, contradicting the
hypothesis $\varphi\leq\psi$ on $X^{\divis}$.
\end{proof}

\begin{remark}
\label{rem:ch32-density-not-enough}
\Cref{cor:ch32-determination-divisorial}\,\cref{itm:ch32-dd:usc} is the
elementary form of a recurring trap.  A $\theta$-psh function is upper
semicontinuous and may be $-\infty$ somewhere, so density of
$X^{\divis}$ alone does not force two psh functions agreeing on
$X^{\divis}$ to be equal.  The correct statement --- they \emph{are}
equal --- is the unconditional usc-extension and divisorial comparison
theorem of \cite[Theorem~4.22 and Corollary~4.23]{BJ22}, imported in
\Cref{thm:bj-basic-closure}\,\cref{itm:bj-closure:determination}.  It does
not require the envelope property.  The present chapter supplies only the
topological input.
\end{remark}

\begin{example}
\label{ex:ch32-density-line}
For $X=\mathbf P^1_k$, \Cref{thm:ch32-p1-classification} gives
\[
 X^{\divis}=\{v_\triv\}\ \sqcup\
 \coprod_{p\in\mathbf P^1(k)}\iota_p(\mathbf Q_{>0}),
\]
since the only divisorial valuations of $k(T)$ are the multiples of the
$\ord_p$.  Density is visible: each ray is a copy of $[0,+\infty]$ in
which $\mathbf Q_{>0}$ is dense, together with the endpoint
$v_{p,\triv}=\lim_{t\to+\infty}\iota_p(t)$.  Note that the endpoints
$v_{p,\triv}$ are \emph{not} divisorial --- they are proper-support
points --- and yet they lie in the closure of $X^{\divis}$; this is
precisely the phenomenon warned about after
\Cref{thm:triv-divisorial-density}.  Note also that $X^{\divis}$ is dense
but far from all of $X^{\An}$: its complement contains the uncountable set
$\coprod_p\iota_p\bigl((0,+\infty)\setminus\mathbf Q\bigr)$.
\end{example}

\subsection{Divisorial points and test configurations}
\label{sec:ch32-divisorial-tc}

The link between \Cref{def:ch32-divisorial} and the test configurations of
\Cref{chap:bj-test-configurations} is the following statement, which we
quote; it is the trivially valued counterpart of the model-component
description $x_E=b_E^{-1}\ord_E$ used throughout \Cref{part:bfj}.

\refstepcounter{theorem}\label{thm:ch32-divisorial-tc}
\begin{quote}
\textbf{Foundational input \thetheorem\ (divisorial points from central
fibres; source boundary).}
\index{test configuration!divisorial points from}
Let $X$ be a projective $k$-variety and $\cX$ an integrally closed test
configuration for $X$ in the sense of \Cref{def:integrally-closed-tc}.  For
each irreducible component $E$ of the central fibre $\cX_0$, the local
ring of $\cX$ at the generic point of $E$ is a discrete valuation ring;
setting $b_E:=\ord_E(\varpi)$ and
\[
 v_E:=\bigl(b_E^{-1}\ord_E\bigr)\big|_{k(X)},
\]
one obtains a divisorial valuation $v_E\in X^{\divis}$.  Conversely every
$v\in X^{\divis}$ arises in this way from some integrally closed test
configuration.  This is \cite[Lemma~1.20]{BJ22}, resting on
\cite[Theorem~4.6]{BHJ17}; the underlying algebraic input is Zariski's
theorem that a divisorial valuation is realized on a blowup, as recorded in
\Cref{thm:appB-rees-valuations}.
The two directions are imported from the cited results; no proof is claimed
here.
\end{quote}

\begin{remark}
\label{rem:ch32-bE}
The normalization by $b_E$ in \Cref{thm:ch32-divisorial-tc} is not a
choice: it is forced by the requirement $w(\varpi)=1$ on the
test-configuration side, and it disappears after restriction to $k(X)$,
where no element plays the role of $\varpi$.  The identity
$\sigma(v_E)=b_E^{-1}\ord_E$ of \eqref{eq:ch30-gauss-divisorial} says that
the Gauss extension puts the factor back.  A formula in $\PL(X)$ that
looks as though it has lost a factor $b_E$ relative to the corresponding
formula in \Cref{part:bfj} has almost always lost it here; the unit tests of
\Cref{sec:appJ-unit-tests} are designed to catch this.
\end{remark}

\begin{example}
\label{ex:ch32-normal-cone-valuation}
Let $X$ be an integral projective variety, $Z\subsetneq X$ a nonempty
closed subscheme with nonzero ideal $\mathfrak b$, and let
$\cX\to X\times\mathbf A^1$ be the normalized
blowup of $Z\times\{0\}$, the normalized deformation to the normal cone.
Its central fibre contains the strict transform $\widetilde X$ of
$X\times\{0\}$ and may contain several exceptional prime components
$P_1,\ldots,P_r$.  The associated divisorial valuations are
\[
 v_{\widetilde X}=v_{X,\triv},
 \qquad
 v_{P_i}=b_i^{-1}\ord_{P_i}|_{k(X)},
 \qquad b_i=\ord_{P_i}(\varpi).
\]
The restrictions $\ord_{P_i}|_{k(X)}$ are the Rees valuations of
$\mathfrak b$, and $b_i=\ord_{P_i}(\mathfrak b)$; see
\Cref{thm:appB-rees-valuations}.  Thus each nontrivial point is the
corresponding normalized Rees valuation.  When $\mathfrak b$ has a single
Rees valuation, there is a single exceptional prime relevant to this
formula.
Two cautions.  The set-theoretic recipe
$f\mapsto\max\{m:f\in\overline{\mathfrak b^{\,m}}\}$ computes
$\lfloor\min_i v_{P_i}(f)\rfloor$ (and $\lfloor v_P(f)\rfloor$ in the
single-Rees-valuation case), not a valuation:
for $\mathfrak b=(x^2,y^3)\subset k[x,y]$ it returns $0$ at $x$ and
$1$ at $x^2$, whereas the normalized Rees valuation has
$v_P(x)=\tfrac12$ and $v_P(y)=\tfrac13$
(\Cref{ex:appB-monomial-rees}).  A general
$\mathfrak b$ has several Rees valuations --- take
$\mathfrak b=\mathfrak m_p\cap\mathfrak m_q$ for distinct closed points
--- so there is no single exceptional prime in general.  In particular,
for $Z=\{p\}$ a smooth closed point of a smooth $X$, the exceptional
divisor is irreducible and $v_P=\ord_p$ is the order of vanishing at $p$;
the associated PL function is
$\phi(v)=-\min\{v(\mathfrak m_p),1\}$.  The homogeneous potential
$-v(\mathfrak m_p)$ that appears when the truncation is removed is
\eqref{eq:ch08-ideal-potential}; its truncation mechanism is
\Cref{ex:ch08-truncation}.
\end{example}

\section{Quasimonomial points}
\label{sec:ch32-quasimonomial}

Quasimonomial valuations enlarge this computable locus.  Suppose $Y$ is
regular at the generic point of a stratum of an SNC divisor
$E_1+\cdots+E_r$, with local equations $z_1,\ldots,z_r$.  For
$\alpha\in\mathbf R_{\geq0}^r$, define
\begin{equation}
\label{eq:quasimonomial-valuation}
 v_\alpha\!\left(\sum_{\beta\in\mathbf N^r}c_\beta z^\beta\right)
 =\min_{c_\beta\ne0}\langle\alpha,\beta\rangle .
\end{equation}
Completion of the regular local ring makes the expansion meaningful, and
faithful flatness shows that the restriction is independent of the chosen
expansion.  Coordinate changes preserving the SNC stratum induce the same
valuation.  Varying $\alpha$ produces a cone; the slice
$\sum_i\alpha_i=1$ is the corresponding simplex.

\subsection{The construction, in detail}
\label{sec:ch32-qm-construction}

We set up the data precisely, because two things differ from the discretely
valued construction of \Cref{sec:skeleton-monomial}, and both differences
matter.

\begin{definition}
\label{def:ch32-snc-pair}
\index{SNC pair!over a trivially valued field}
Let $X$ be an integral projective $k$-variety of dimension $n$.  An
\emph{SNC pair over $X$} is a pair $(Y,E)$ consisting of a projective
birational morphism $\pi:Y\to X$ with $Y$ regular, together with a
reduced divisor $E=E_1+\cdots+E_N$ on $Y$ with simple normal crossings
whose support contains the exceptional locus of $\pi$.  For
$J\subset\{1,\ldots,N\}$ we write $E_J:=\bigcap_{j\in J}E_j$ and, for a
connected component $W$ of $E_J$, $\eta_W$ for its generic point; by
\Cref{lem:snc-normal-form} the local ring $\cO_{Y,\eta_W}$ is regular of
dimension $\#J$ with regular system of parameters
$(z_j)_{j\in J}$ given by local equations of the $E_j$.  The
\emph{quasimonomial cone} of $(Y,E)$ is
\[
 \sigma_{Y,E}:=\left(\coprod_{(J,W)}\mathbf R_{\geq0}^{J}\right)\big/\!\sim,
\]
where the face obtained by setting the coordinates in $J\setminus I$ to
zero is identified with the cone indexed by the component of $E_I$
containing $W$.  The \emph{dual complex} $\Delta_{Y,E}$ is the base
$\sum_j\alpha_j=1$ of this cone complex; equivalently, it has a
$(\#J-1)$-simplex for each pair $(J,W)$, with the same face identifications.
\end{definition}

\emph{Difference one: the cone has a vertex, not a base.}  Over a discretely
valued field the constraint $v(\varpi)=1$ reads
$\sum_jb_j\alpha_j=1$ and cuts the cone $\mathbf R_{\geq0}^J$ down to a
simplex; the dual complex is that simplex, and the skeleton is a compact
polyhedron sitting inside $X^{\An}$
(\Cref{thm:skeleton-embedding}).  Over a trivially valued field there is no
such constraint, the whole cone $\mathbf R_{\geq0}^J$ is realized inside
$X^{\An}$, and the origin is the trivial valuation $v_{X,\triv}$.  The
\emph{dual complex} of \Cref{def:ch32-snc-pair} is a choice of base for the
cone, not an intrinsic slice; the choice $\sum_j\alpha_j=1$ is the one used
in the drawings, and the choice $A_X=1$ is the one used in K-stability.

\emph{Difference two: zero weights encode faces.}
Allowing $\alpha_j=0$ produces the quasimonomial
valuation of the smaller stratum $J_+=\{j:\alpha_j>0\}$; this makes the
cones fit together along faces.  The value $\alpha=0$ gives
$v_{X,\triv}$.

\begin{proposition}
\label{prop:ch32-qm-well-defined}
\index{quasimonomial valuation!construction}
With the notation of \Cref{def:ch32-snc-pair}, fix $(J,W)$ and
$\alpha\in\mathbf R_{>0}^J$.  Then:
\begin{enumerate}
\item\label{itm:ch32-qm:min}
for $f\in\widehat{\cO}_{Y,\eta_W}\simeq\kappa(\eta_W)[\![z_j:j\in J]\!]$,
the minimum in \eqref{eq:quasimonomial-valuation} is attained, and
\eqref{eq:quasimonomial-valuation} defines a semivaluation on the completion;
\item\label{itm:ch32-qm:mult}
this semivaluation is multiplicative, and its restriction to
$\cO_{Y,\eta_W}$ has zero kernel, hence defines a valuation
$v_\alpha$ of $k(Y)^\times=k(X)^\times$ trivial on $k^\times$;
\item\label{itm:ch32-qm:indep}
$v_\alpha$ does not depend on the choice of coefficient field in
Cohen's structure theorem, nor on the choice of local equations $z_j$;
\item\label{itm:ch32-qm:face}
if $\alpha_j\to0$ for $j\notin J_+$ with $\alpha|_{J_+}$ fixed, then
$v_\alpha\to v_{\alpha|_{J_+}}$, the quasimonomial valuation of the
stratum $J_+$; in particular the cones of \Cref{def:ch32-snc-pair} glue
continuously along faces.
\end{enumerate}
\end{proposition}

\begin{proof}
\cref{itm:ch32-qm:min}  Since all $\alpha_j>0$, one has
$\langle\alpha,\beta\rangle\geq(\min_j\alpha_j)\,|\beta|$, and for a given
bound $C$ only finitely many $\beta\in\mathbf N^J$ satisfy
$|\beta|\leq C/\min_j\alpha_j$; so the set of values
$\{\langle\alpha,\beta\rangle:c_\beta\neq0\}$ is discrete near $0$ and
the infimum is attained.  The ultrametric inequality holds because the
coefficient of $z^\beta$ in $f+g$ is $c_\beta+d_\beta$, which vanishes
unless $c_\beta\neq0$ or $d_\beta\neq0$; and $v_\alpha(1)=0$,
$v_\alpha(0)=+\infty$ by convention.

\cref{itm:ch32-qm:mult}  Multiplicativity is the assertion that the
associated graded ring of the $\alpha$-weighted filtration
\[
 \mathfrak a_{\geq t}:=\Bigl\{\textstyle\sum_\beta c_\beta z^\beta:
 c_\beta=0\ \text{whenever}\ \langle\alpha,\beta\rangle<t\Bigr\}
\]
is a domain.  That graded ring is
$\bigoplus_t\mathfrak a_{\geq t}/\mathfrak a_{>t}
\simeq\kappa(\eta_W)[Z_j:j\in J]$, a polynomial ring in which $Z_j$ is the
class of $z_j$ in degree $\alpha_j$; it is an integral domain, and
therefore the initial form of a product is the product of the initial forms,
which is exactly $v_\alpha(fg)=v_\alpha(f)+v_\alpha(g)$.  This is the
argument of \Cref{lem:monomial-valuation}, which is written there over a
discretely valued base but uses only the structure of the completed local
ring at the stratum and is valid verbatim here; the only change is that the
ground field $k$ replaces the residue field of the discretely valued base,
so no term $v(\varpi)$ occurs.

For the kernel, let $f\in\cO_{Y,\eta_W}$ be nonzero; its image in the
completion is nonzero by Krull's intersection theorem, so some
$c_\beta\neq0$ and $v_\alpha(f)<+\infty$.  Hence $v_\alpha$ restricted
to $\cO_{Y,\eta_W}$ has zero kernel, and extends to the fraction field
$k(Y)=k(X)$ by $v_\alpha(f/g)=v_\alpha(f)-v_\alpha(g)$.  Triviality on
$k^\times$ holds because a nonzero constant has nonzero coefficient
$c_0$.

\cref{itm:ch32-qm:indep}  The graded ring $\operatorname{gr}_\alpha$ and the
filtration $(\mathfrak a_{\geq t})$ are intrinsic to the boundary ideals:
they are the weighted monomial filtration generated by the products of the
principal ideals $(z_j)$.  Replacing $z_j$ by another local equation
multiplies it by a unit and leaves these ideals unchanged.  The associated
graded construction therefore does not depend on a coefficient field or on
the chosen generators.  Since
$v_\alpha(f)=\sup\{t:f\in\mathfrak a_{\geq t}\}$, neither does the
valuation.

\cref{itm:ch32-qm:face}  Fix $f\in\cO_{Y,\eta_W}$ nonzero and write it in
the completion, and let
$S=\{\beta\in\mathbf N^J:c_\beta\neq0\}$.  Dickson's lemma gives a finite
set $M\subset S$ of componentwise minimal elements such that every
$\beta\in S$ dominates some $\mu\in M$.  For every nonnegative weight
vector $\gamma$ one therefore has
\[
 \min_{\beta\in S}\langle\gamma,\beta\rangle
 =\min_{\mu\in M}\langle\gamma,\mu\rangle.
\]
The right-hand side is the minimum of finitely many continuous functions of
$\gamma$.  Letting the coordinates outside $J_+$ tend to zero thus gives
the monomial value attached to the component of $E_{J_+}$ containing $W$.
Convergence in $X^{\An}$ is pointwise convergence on regular functions, so
the cones glue continuously.
\end{proof}

\begin{theorem}
\label{thm:ch32-quasimonomial-invariants}
\index{quasimonomial valuation!invariants}
Let $(Y,E)$ be an SNC pair over the integral projective $k$-variety
$X$ of dimension $n$, let $(J,W)$ be a stratum with generic point
$\eta_W$, and let $\alpha\in\mathbf R_{\geq0}^J$, $\alpha\neq0$.  Put
$J_+:=\{j\in J:\alpha_j>0\}$, $q:=\#J_+$, and
$d:=\dim_{\mathbf Q}\sum_{j\in J_+}\mathbf Q\alpha_j$.  Then:
\begin{enumerate}
\item\label{itm:ch32-qi:center}
the center of $v_\alpha$ on $Y$ is the generic point of the stratum
$E_{J_+}$ containing $W$, and its center on $X$ is the generic point of
the image of that stratum;
\item\label{itm:ch32-qi:group}
the value group of $v_\alpha$ is $\sum_{j\in J_+}\mathbf Z\alpha_j$, of
rational rank $d$;
\item\label{itm:ch32-qi:residue}
the residue field of $v_\alpha$ is a purely transcendental extension of
$\kappa(\eta_{W'})$ of transcendence degree $q-d$, where $W'$ is the
component of $E_{J_+}$ containing $W$; hence
$\operatorname{trdeg}(v_\alpha)=(n-q)+(q-d)=n-d$;
\item\label{itm:ch32-qi:abhyankar}
$v_\alpha$ is Abhyankar:
$\operatorname{rat.rk}(v_\alpha)+\operatorname{trdeg}(v_\alpha)=d+(n-d)=n$;
\item\label{itm:ch32-qi:vertex}
if $J_+=\{j\}$, then $v_\alpha=\alpha_j\ord_{E_j}$; in particular the
rays of the cone through the vertices of $\Delta_{Y,E}$ are the divisorial
rays $\mathbf R_{>0}\ord_{E_j}$;
\item\label{itm:ch32-qi:divisorial}
$v_\alpha$ is divisorial if and only if
$\alpha=t(p_j)_{j\in J_+}$ with $(p_j)$ a primitive vector of positive
integers and $t\in\mathbf Q_{>0}$; otherwise
$v_\alpha\notin X^{\divis}$.
\end{enumerate}
\end{theorem}

\begin{proof}
\cref{itm:ch32-qi:center}  By construction $v_\alpha\geq0$ on
$\cO_{Y,\eta_{W'}}$, and $v_\alpha>0$ on its maximal ideal, since a
nonunit there has $c_0=0$ and every monomial $z^\beta$, $\beta\neq0$,
supported on $J_+$ has weight $>0$; while $z_j$ for
$j\in J\setminus J_+$ has weight $0$ and is therefore a unit for
$v_\alpha$, which is why the center is the generic point of $E_{J_+}$
rather than of $E_J$.  Uniqueness of centers
(\Cref{prop:centers-proper-separated}) identifies it.  The center on $X$ is
the image, by the argument of
\Cref{prop:ch32-divisorial-invariants}\,\cref{itm:ch32-di:center}.

\cref{itm:ch32-qi:group}  By multiplicativity the value group is generated by
the values on $\cO_{Y,\eta_{W'}}$, and every such value is of the form
$\langle\alpha,\beta\rangle$ with $\beta\in\mathbf N^{J_+}$; conversely
each $\alpha_j=v_\alpha(z_j)$ occurs.  Hence the group is
$\sum_{j\in J_+}\mathbf Z\alpha_j$, whose rational rank is $d$ by
definition of $d$.

\cref{itm:ch32-qi:residue}  The residue field of $v_\alpha$ is the
degree-zero part of the graded fraction field of
$\operatorname{gr}_\alpha\simeq\kappa(\eta_{W'})[Z_j:j\in J_+]$, where
$\deg Z_j=\alpha_j$.  A ratio of monomials $Z^\beta/Z^\gamma$ has degree
zero exactly when $\delta:=\beta-\gamma$ lies in
$L:=\ker\bigl(\mathbf Z^{J_+}\to\mathbf R,\ \delta\mapsto
\langle\alpha,\delta\rangle\bigr)$, a subgroup of rank $q-d$.  The
degree-zero part is therefore the fraction field of the group algebra
$\kappa(\eta_{W'})[L]$, a Laurent polynomial ring in $q-d$ variables,
hence purely transcendental of transcendence degree $q-d$.  Since
$E_{J_+}$ has codimension $q$ in $Y$
(\Cref{lem:snc-normal-form}\,(3)), $\kappa(\eta_{W'})$ has transcendence
degree $n-q$ over $k$.  Adding gives $n-d$.

\cref{itm:ch32-qi:abhyankar}  Add \cref{itm:ch32-qi:group} and
\cref{itm:ch32-qi:residue} and compare with
\Cref{thm:abhyankar-inequality}.

\cref{itm:ch32-qi:vertex}  If $J_+=\{j\}$, then
$\widehat{\cO}_{Y,\eta_{E_j}}$ is a complete discrete valuation ring with
uniformizer $z_j$ and \eqref{eq:quasimonomial-valuation} reads
$v_\alpha(f)=\alpha_j\ord_{z_j}(f)=\alpha_j\ord_{E_j}(f)$.

\cref{itm:ch32-qi:divisorial}  If $v_\alpha$ is divisorial, then its value
group is $t\mathbf Z$ with $t\in\mathbf Q_{>0}$ by
\Cref{prop:ch32-divisorial-invariants}\,\cref{itm:ch32-di:group}, so
$d=1$ and all $\alpha_j$ lie in $t\mathbf Z$, hence are rational
multiples of one another.  Dividing their common integer vector by its
greatest common divisor gives a unique primitive positive vector $(p_j)$
and a scalar $t>0$.  Conversely, if
$\alpha=t(p_j)_{j\in J_+}$ with $(p_j)$ primitive and positive, then
$t^{-1}v_\alpha$ is $\mathbf Z$-valued with residue field of
transcendence degree $n-1$, and it is realized as $\ord_F$ on the
toroidal weighted blowup of $Y$ along $E_{J_+}$ with weights $(p_j)$;
this is the toroidal subdivision statement
\Cref{thm:skeleton-subdivision}, whose construction does not use the ground
valuation.  Then $v_\alpha=t\ord_F$, divisorial when $t\in\mathbf Q$ and
not otherwise, by the value-group argument just given.
\end{proof}

\begin{proposition}
\label{prop:ch32-log-discrepancy-affine}
\index{log discrepancy!affine on quasimonomial cones}
Let $X$ be a normal $\mathbf Q$-Gorenstein integral projective
$k$-variety and $(Y,E)$ an SNC pair over $X$.  Define, for a stratum
$(J,W)$ and $\alpha\in\mathbf R_{\geq0}^J$,
\begin{equation}
\label{eq:ch32-log-discrepancy-affine}
 A_X(v_\alpha):=\sum_{j\in J}\alpha_j\,A_X(\ord_{E_j}),
 \qquad
 A_X(\ord_{E_j}):=1+\ord_{E_j}\bigl(K_{Y}-\pi^*K_X\bigr).
\end{equation}
Then:
\begin{enumerate}
\item\label{itm:ch32-lda:consistent}
the definition is consistent, in the sense that whenever $v_\alpha$ is
itself divisorial, say $v_\alpha=t\ord_F$, the value
\eqref{eq:ch32-log-discrepancy-affine} agrees with
$t\bigl(1+\ord_F(K_{Y'}-\pi'^*K_X)\bigr)$ computed on any model $Y'$
carrying $F$, that is, with \eqref{eq:appJ-div-discrepancy};
\item\label{itm:ch32-lda:positive}
$A_X$ is finite, continuous and homogeneous on $\sigma_{Y,E}$, and
compatible with the face maps of
\Cref{prop:ch32-qm-well-defined}\,\cref{itm:ch32-qm:face}; if $X$ is klt,
then it is positive away from the zero valuation;
\item\label{itm:ch32-lda:link}
if $X$ is klt, then
$\operatorname{Link}_{A_X}(\sigma_{Y,E}\setminus\{0\})$ is a compact
polyhedral complex, namely the copy of $\Delta_{Y,E}$ normalized by
$A_X=1$.
\end{enumerate}
\end{proposition}

\begin{proof}
\cref{itm:ch32-lda:consistent}  The second formula in
\eqref{eq:ch32-log-discrepancy-affine} is \eqref{eq:appJ-div-discrepancy}
at $c=1$, so the case $\#J_+=1$ is a tautology.  In general, by
\Cref{thm:ch32-quasimonomial-invariants}\,\cref{itm:ch32-qi:divisorial}, a
divisorial $v_\alpha$ has $\alpha=t(p_j)_{j\in J_+}$ with
$t\in\mathbf Q_{>0}$ and $(p_j)$ a primitive positive integer vector, and
$v_\alpha=t\ord_F$
with $F$ the exceptional divisor of the toroidal weighted blowup with
weights $(p_j)$; since log discrepancy of a divisorial valuation does not
depend on the model carrying the divisor, it suffices to compute on that
blowup.  Let
$\rho:Y'\to Y$ be the weighted blowup along $E_{J_+}$ with weights
$(p_j)$, with exceptional prime divisor $F$, so that
$v_\alpha=t\ord_F$ by
\Cref{thm:ch32-quasimonomial-invariants}\,\cref{itm:ch32-qi:divisorial}.
The toroidal computation of the relative canonical divisor gives
\[
 K_{Y'}-\rho^*K_Y=\Bigl(\sum_{j\in J_+}p_j-1\Bigr)F,
 \qquad
 \rho^*E_j=E_j'+p_jF\quad(j\in J_+),
\]
whence
\[
 \ord_F\bigl(K_{Y'}-\rho^*\pi^*K_X\bigr)
 =\Bigl(\sum_{j\in J_+}p_j\Bigr)-1
  +\sum_{j\in J_+}p_j\,\ord_{E_j}(K_Y-\pi^*K_X),
\]
so that
\[
\begin{split}
 A_X(\ord_F)&=1+\ord_F(K_{Y'}-\rho^*\pi^*K_X)\\
 &=\sum_{j\in J_+}p_j\Bigl(1+\ord_{E_j}(K_Y-\pi^*K_X)\Bigr)
 =\sum_{j\in J_+}p_jA_X(\ord_{E_j}).
\end{split}
\]
Multiplying by $t$ gives exactly the right-hand side of
\eqref{eq:ch32-log-discrepancy-affine}, which proves consistency.

\cref{itm:ch32-lda:positive}  The right-hand side of
\eqref{eq:ch32-log-discrepancy-affine} is an affine function of $\alpha$
with finitely many terms, hence finite, continuous and homogeneous of degree
one; compatibility with the face maps is the observation that the terms with
$\alpha_j=0$ drop out, which matches
\Cref{prop:ch32-qm-well-defined}\,\cref{itm:ch32-qm:face}.  If $X$ is klt,
then $A_X(\ord_{E_j})>0$ for every $j$, which gives positivity away from
the zero valuation.  Without the klt hypothesis no positivity is asserted.

\cref{itm:ch32-lda:link}  Assume $X$ is klt and apply
\Cref{thm:ch32-link-cone} with
$\Omega=\sigma_{Y,E}\setminus\{0\}$ and $A=A_X$.  Compactness of
$\{A_X=1\}\cap\sigma_{Y,E}$ holds because it is a finite union of compact
simplices: the piece lying in the stratum cone $\mathbf R_{\geq0}^J$ is cut
out by the single affine equation $\sum_j\alpha_jA_X(\ord_{E_j})=1$ with
positive coefficients, hence is a compact simplex.
\end{proof}

\begin{remark}
\label{rem:ch32-link-single-model}
When $X$ is klt, \Cref{prop:ch32-log-discrepancy-affine} produces a compact
link inside a fixed model's quasimonomial cone.  This is compatible with
\Cref{thm:ch32-no-compact-link}: the latter only records the fixed-point
obstruction to a gauge on all of $X^{\An}$.  No compactness claim is made
here for a union or an inverse limit of normalized links over all models.
\end{remark}

\begin{theorem}
\label{thm:ch32-qm-density}
\index{skeleton!inverse limit}
Let $X$ be an integral projective $k$-variety.
\begin{enumerate}
\item\label{itm:ch32-qmd:dense}
Write $X^{\mathrm{qm}}$ for the set of quasimonomial points, that is, the
union of the cones $\sigma_{Y,E}$ over all SNC pairs over $X$.  Every
vertex ray of such a cone is a divisorial ray.  Conversely, if
$\operatorname{char}k=0$, then $X^{\divis}\subset X^{\mathrm{qm}}$, and
consequently $X^{\mathrm{qm}}$ is dense in $X^{\An}$.
\item\label{itm:ch32-qmd:abhyankar}
Every quasimonomial point is Abhyankar.  Conversely, if
$\operatorname{char}k=0$ and $X$ is smooth, every Abhyankar valuation of
$k(X)/k$ is quasimonomial; this converse is \cite[\S3]{JM12}, and is not
used below.
\item\label{itm:ch32-qmd:limit}
If $\operatorname{char}k=0$ and $X$ is smooth, then $X^{\An}$ is the
inverse limit of the dual complexes of SNC test configurations, with the
retraction maps as transition morphisms; this is
\cite[Appendix~A]{BJ22}, whose discretely valued counterpart is
\Cref{thm:skeleton-inverse-limit}.
\end{enumerate}
\end{theorem}

\begin{proof}
\cref{itm:ch32-qmd:dense}  That a vertex ray is divisorial is
\Cref{thm:ch32-quasimonomial-invariants}\,\cref{itm:ch32-qi:vertex}.
Conversely, let $v=t\ord_E$ with $t\in\mathbf Q_{>0}$ and $E$ a prime
divisor on a normal projective birational model $Y_0\to X$.  In
characteristic zero, \Cref{thm:appB-resolution} provides a log resolution
$Y\to Y_0$ of the pair consisting of $Y_0$, its singular locus, and the
exceptional locus, on which the strict transform of $E$ appears as a
component of an SNC divisor whose support contains the exceptional locus over
$X$; taking $\alpha=t$ at the corresponding vertex gives
$v=v_\alpha$.  Density then follows from
\Cref{thm:triv-divisorial-density}.  In positive characteristic the
existence of such a log resolution is not known in dimension $\geq4$, and
we make no claim.

\cref{itm:ch32-qmd:abhyankar}  The direct statement is
\Cref{thm:ch32-quasimonomial-invariants}\,\cref{itm:ch32-qi:abhyankar}.  The
converse is quoted.

\cref{itm:ch32-qmd:limit}  Quoted.
\end{proof}

\subsection{The affine plane: monomial valuations and their invariants}
\label{sec:ch32-affine-plane}

\begin{example}
\label{ex:ch32-affine-plane}
\index{monomial valuation!invariants, tabulated}
Let $X=\mathbf A^2_k$ with coordinates $x,y$, and let
$E_1=\{x=0\}$, $E_2=\{y=0\}$, an SNC divisor with stratum the origin.
For $w=(w_1,w_2)\in\mathbf R^2_{\geq0}\setminus\{0\}$ let $v_w$ be the
monomial valuation \eqref{eq:quasimonomial-valuation}, that is,
\[
 v_w\Bigl(\sum_{\beta\in\mathbf N^2}c_\beta x^{\beta_1}y^{\beta_2}\Bigr)
 =\min_{c_\beta\neq0}\bigl(\beta_1w_1+\beta_2w_2\bigr).
\]
\Cref{thm:ch32-quasimonomial-invariants} and
\Cref{lem:monomial-invariants} give the following table, in which
$d=\dim_{\mathbf Q}(\mathbf Qw_1+\mathbf Qw_2)$ counted only over the
nonzero coordinates.

\begin{center}
\small
\begin{tabular}{l|c|c|c|c|c}
\toprule
$w$&center on $\mathbf A^2$&$\Gamma_{v_w}$&
 $\operatorname{rat.rk}$&$\operatorname{trdeg}$&$A_{\mathbf A^2}$\\
\midrule
$(s,0)$, $s>0$&generic pt of $E_1$&$s\mathbf Z$&$1$&$1$&$s$\\
$(0,s)$, $s>0$&generic pt of $E_2$&$s\mathbf Z$&$1$&$1$&$s$\\
$(p,q)\in\mathbf Z_{>0}^2$ coprime&origin&$\mathbf Z$&$1$&$1$&$p+q$\\
$c(p,q)$, $c>0$, $(p,q)$ as above&origin&$c\mathbf Z$&$1$&$1$&$c(p+q)$\\
$(1,\theta)$, $\theta\notin\mathbf Q$&origin&
 $\mathbf Z+\theta\mathbf Z$&$2$&$0$&$1+\theta$\\
\bottomrule
\end{tabular}
\end{center}

Every row satisfies Abhyankar's equality
$\operatorname{rat.rk}+\operatorname{trdeg}=2$.  The support of every
$v_w$ is all of $\mathbf A^2$, so all these points lie in
$X^{\val}$.  The third row is divisorial: $v_{(p,q)}=\ord_F$ for $F$
the exceptional divisor of the weighted blowup of weights $(p,q)$, and
$A(\ord_F)=1+(p+q-1)=p+q$ by \eqref{eq:appJ-div-discrepancy}, in agreement
with \eqref{eq:ch32-log-discrepancy-affine} and with
\Cref{prop:ch08-monomial-invariants}\,(2).  The fifth row is not divisorial,
by
\Cref{thm:ch32-quasimonomial-invariants}\,\cref{itm:ch32-qi:divisorial}; it is
the irrational point analysed in \Cref{ex:gallery-irrational}.

Three explicit numbers, to fix the picture.
\begin{enumerate}
\item\emph{The origin's own valuation.}  $w=(1,1)$ gives
$v_{(1,1)}=\ord_0$, the order of vanishing at the origin, with
$A(\ord_0)=2$ and $v(\mathfrak m_0)=1$.
\item\emph{A weighted point.}  $w=(2,3)$ gives the divisorial valuation of
the $(2,3)$-weighted blowup, with $A=5$,
$v(\mathfrak m_0)=\min\{2,3\}=2$, and
$v\bigl((x^2,y^3)\bigr)=\min\{4,9\}=4$; the associated homogeneous
potential is
$\varphi_{(x^2,y^3)}(v_w)=\max\{-2w_1,-3w_2\}$, computed in
\Cref{ex:ch08-monomial-multiplicity}.
\item\emph{The $A=1$ link.}  By
\eqref{eq:ch32-log-discrepancy-affine}, $A_{\mathbf A^2}(v_w)=w_1+w_2$,
so the normalized link is the segment
$\{w_1+w_2=1,\ w\geq0\}$, joining $\ord_{E_1}$ to $\ord_{E_2}$, with
$\tfrac12\ord_0$ at its midpoint $(1/2,1/2)$.
\end{enumerate}
The factor $\tfrac12$ is the normalization warning of
\Cref{rem:divisorial-normalization} in its trivially valued form.
\end{example}

\subsection{A quasimonomial point on a blown-up surface}
\label{sec:ch32-blowup-surface}

\begin{example}
\label{ex:ch32-blowup-surface}
\index{quasimonomial valuation!on a blown-up surface}
Let $X=\mathbf A^2_k$ with coordinates $(x,y)$, let $p$ be the origin,
and construct a tower
\[
 Y_2\xrightarrow{\ \pi_2\ }Y_1\xrightarrow{\ \pi_1\ }X
\]
as follows: $\pi_1$ is the blowup of $p$, with exceptional prime divisor
$E_1$; and $\pi_2$ is the blowup of the point $q\in E_1$ corresponding
to the tangent direction $y=0$, with exceptional prime divisor $E_2$.
Write $\pi=\pi_1\circ\pi_2$ and $E_1'$ for the strict transform of
$E_1$ on $Y_2$.  The pair $(Y_2,E_1'+E_2)$ is an SNC pair over
$X$, and $E_1'\cap E_2$ is a single closed point $\eta$, a
zero-dimensional stratum.

\emph{Coordinates.}  On $Y_1$, work in the chart $(x,t)$ with
$(x,y)=(x,xt)$, so $E_1=\{x=0\}$ and $q=(0,0)$.  Blow up $q$ and use
the chart $(z_1,z_2)$ with
\[
 x=z_1z_2,\qquad t=z_2 ,
 \qquad\text{so that}\qquad
 E_1'=\{z_1=0\},\quad E_2=\{z_2=0\}.
\]
Then
\begin{equation}
\label{eq:ch32-blowup-coordinates}
 x=z_1z_2,\qquad y=xt=z_1z_2^2 .
\end{equation}
Let $v_\alpha$, $\alpha=(\alpha_1,\alpha_2)\in\mathbf R^2_{\geq0}$, be
the quasimonomial valuation at $\eta$ with $v_\alpha(z_i)=\alpha_i$.
\begin{enumerate}
\item\label{itm:ch32-bs:values}
\emph{Values on the original coordinates.}  By
\eqref{eq:ch32-blowup-coordinates},
\[
\begin{gathered}
 v_\alpha(x)=\alpha_1+\alpha_2,
 \qquad
 v_\alpha(y)=\alpha_1+2\alpha_2,\\
 v_\alpha(\mathfrak m_p)=\min\{v_\alpha(x),v_\alpha(y)\}
 =\alpha_1+\alpha_2 .
\end{gathered}
\]
At the vertices: $\alpha=(1,0)$ gives $\ord_{E_1}$, with
$\ord_{E_1}(x)=\ord_{E_1}(y)=1$; and $\alpha=(0,1)$ gives
$\ord_{E_2}$, with $\ord_{E_2}(x)=1$, $\ord_{E_2}(y)=2$.
\item\label{itm:ch32-bs:discrepancy}
\emph{Log discrepancies.}  On $Y_1$ one has
$K_{Y_1}-\pi_1^*K_X=E_1$, so $A_X(\ord_{E_1})=1+1=2$.  On $Y_2$,
$K_{Y_2}-\pi_2^*K_{Y_1}=E_2$ and $\pi_2^*E_1=E_1'+E_2$, so
\[
 K_{Y_2}-\pi^*K_X=\pi_2^*E_1+E_2=E_1'+2E_2,
 \qquad
 A_X(\ord_{E_2})=1+2=3 .
\]
By \eqref{eq:ch32-log-discrepancy-affine},
\begin{equation}
\label{eq:ch32-blowup-discrepancy}
 A_X(v_\alpha)=2\alpha_1+3\alpha_2 .
\end{equation}
The normalized link $\{A_X=1\}$ is the segment joining
$\tfrac12\ord_{E_1}$ to $\tfrac13\ord_{E_2}$.
\item\label{itm:ch32-bs:consistency}
\emph{A consistency check with \Cref{ex:ch32-affine-plane}.}  Take
$\alpha=(1,1)$.  Then $v_\alpha(x)=2$ and $v_\alpha(y)=3$, and in fact
$v_{(1,1)}=v_{(2,3)}$, the monomial valuation of weights $(2,3)$ in the
coordinates $(x,y)$.  Indeed a monomial $x^ay^b$ becomes
$z_1^{a+b}z_2^{a+2b}$, and the linear map
$(a,b)\mapsto(a+b,a+2b)$ is injective, so distinct monomials of
$k[x,y]$ go to distinct monomials of $k[z_1,z_2]$ and no cancellation can
occur; therefore
$v_{(1,1)}\bigl(\sum c_{ab}x^ay^b\bigr)=\min_{c_{ab}\neq0}(2a+3b)$, which
is $v_{(2,3)}$.  The log discrepancies agree:
\eqref{eq:ch32-blowup-discrepancy} gives $2+3=5$, and the table in
\Cref{ex:ch32-affine-plane} gives $p+q=2+3=5$.
\item\label{itm:ch32-bs:irrational}
\emph{An irrational quasimonomial point.}  Take
$\alpha=(1,\sqrt2)$.  Then the value group is
$\mathbf Z+\sqrt2\,\mathbf Z$, of rational rank $2$; the residue field is
$k$, of transcendence degree $0$; Abhyankar's equality reads
$2+0=2$.  By
\Cref{thm:ch32-quasimonomial-invariants}\,\cref{itm:ch32-qi:divisorial} this
point is not divisorial, and no rescaling makes it so.  Its log discrepancy
is $A_X=2+3\sqrt2$, and
$v_\alpha(\mathfrak m_p)=1+\sqrt2$.  Note the contrast with the point
$t\ord_{E_1}$ for $t$ irrational, which is also non-divisorial but has
rational rank one: non-divisoriality has two distinct causes, an irrational
scaling and an irrational \emph{ratio} of weights, and only the second is
intrinsic.
\item\label{itm:ch32-bs:center}
\emph{Centers.}  For $\alpha_1,\alpha_2>0$ the center of $v_\alpha$ on
$Y_2$ is $\eta=E_1'\cap E_2$ and its center on $X$ is $p$; for
$\alpha=(s,0)$ the center on $Y_2$ is the generic point of $E_1'$ and
on $X$ is still $p$.  The center map is constant equal to $p$ on the
whole open cone, in accordance with
\Cref{prop:ch32-scaling-action}\,\cref{itm:ch32-sa:point}, while the center on
$Y_2$ jumps at the boundary: this is the anticontinuity of
\Cref{thm:ch32-support-center-maps}\,\cref{itm:ch32-scm:c} in action.
\end{enumerate}
\end{example}

\begin{remark}
\label{rem:ch32-blowup-lessons}
\Cref{ex:ch32-blowup-surface} contains, in two dimensions, every phenomenon
of \Cref{sec:ch32-quasimonomial}: a two-dimensional cone of valuations with
a one-dimensional link; log discrepancy affine on the cone with the two
vertex values $2$ and $3$; a rational interior point that is divisorial
after a toroidal blowup, and equals a monomial valuation in the original
coordinates; an irrational interior point of rational rank two that is not
divisorial and never will be; and a center map on $X$ that is constant on
the cone while the center on the model jumps.  The reader who wants the
same computation with the Monge--Amp\`ere measure attached should compare
\Cref{ex:ch30-surface-blowup}, where the corresponding discretely valued
surface computation is carried out with its energy.
\end{remark}

\section{Proper support and the value \texorpdfstring{$-\infty$}{-infinity}}
\label{sec:ch32-proper-support}

\subsection{Three families of proper-support points}
\label{sec:ch32-proper-support-examples}

\begin{definition}
\label{def:ch32-proper-support}
\index{proper support}
A point $v\in X^{\An}$ has \emph{proper support} if
$v\notin X^{\val}$, that is, if $s(v)$ is not an irreducible component of
$X$.  We write $X^{\An}\setminus X^{\val}$ for the set of such points.
\end{definition}

Three families are worth having in mind, and all three are used later.

\begin{example}
\label{ex:ch32-proper-support-trivial}
Let $Z\subsetneq X$ be an integral closed subscheme which is not an
irreducible component.  Then $v_{Z,\triv}$ has proper support: it is the
point of $X^{\An}$ with $v(f)=+\infty$ if $f|_Z=0$ and $0$ otherwise.
It is simultaneously a fixed point of the scaling action
(\Cref{thm:ch32-trivial-locus}\,\cref{itm:ch32-tl:fixed}) and a
proper-support point.  For $Z=\{p\}$ a closed point, $v_{p,\triv}$ is the
endpoint of every compactified ray whose valuations are centered at $p$:
by \Cref{prop:ch32-scaling-limits}, $\lim_{t\to+\infty}tv=v_{p,\triv}$
whenever $c(v)=p$.  On the projective line these are exactly the points
$\iota_p(+\infty)$ of \Cref{thm:ch32-p1-classification}.
\end{example}

\begin{example}
\label{ex:ch32-curve-germ}
\index{semivaluation!curve germ}
Let $X$ be a variety over $k$ of dimension $n\geq2$, let
$C\subset X$ be an integral curve, and let $p\in C(k)$ be a closed point
at which $C$ is regular.  Define
\begin{equation}
\label{eq:ch32-curve-germ}
 \ord_{(C,p)}(f):=\ord_p\bigl(f|_C\bigr)\in[0,+\infty],
 \qquad f\in\cO_{X,p},
\end{equation}
where $\ord_p$ is the discrete valuation of the regular local ring
$\cO_{C,p}$, with $\ord_p(0)=+\infty$.  Then $\ord_{(C,p)}$ is a
semivaluation on $\cO_{X,p}$, and it extends to a point of $X^{\An}$
with
\[
 s\bigl(\ord_{(C,p)}\bigr)=C,
 \qquad
 c\bigl(\ord_{(C,p)}\bigr)=p .
\]
Since $\dim C=1<n$, the point has proper support.
Concretely, for $X=\mathbf A^2$, $C=\{y=0\}$ and $p$ the origin,
\[
 \ord_{(C,p)}\Bigl(\sum_{a,b}c_{ab}x^ay^b\Bigr)
 =\min\{a:c_{a0}\neq0\},
\]
with the convention that the minimum of the empty set is $+\infty$; thus
$\ord_{(C,p)}(x)=1$ and $\ord_{(C,p)}(y)=+\infty$.  This point is the
limit of full-support monomial valuations:
\begin{equation}
\label{eq:ch32-curve-germ-limit}
 \lim_{t\to+\infty}v_{(1,t)}=\ord_{(C,p)} ,
\end{equation}
because $v_{(1,t)}\bigl(\sum c_{ab}x^ay^b\bigr)
=\min_{c_{ab}\neq0}(a+tb)$, which for large $t$ equals
$\min\{a:c_{a0}\neq0\}$ if some $c_{a0}\neq0$, and tends to $+\infty$
otherwise.  So a proper-support point is genuinely a limit of points of
$X^{\val}$, and no amount of restricting attention to valuations avoids
it.
\end{example}

\begin{example}
\label{ex:ch32-proper-support-subvariety}
Let $Y\subset X$ be integral closed, not a component, and let
$w\in Y^{\val}$ be any valuation of $k(Y)$.  Its image under the closed
embedding $Y^{\An}\hookrightarrow X^{\An}$ of
\Cref{prop:ch32-support-partition}\,\cref{itm:ch32-sp:closed} is a
proper-support point of $X^{\An}$.  In particular, whenever
$\dim Y\geq1$, the proper-support locus of $X^{\An}$ contains an entire
copy of $Y^{\An}$, with all of its own divisorial and quasimonomial
points.  \Cref{ex:ch32-curve-germ} is this construction with $Y=C$ and
$w=\ord_p$.
\end{example}

\subsection{Where potentials take the value \texorpdfstring{$-\infty$}{-infinity}}
\label{sec:ch32-minus-infinity}

Fix a nonzero coherent ideal $\mathfrak a$ and recall the homogeneous
potential
\[
 \varphi_{\mathfrak a}(v)=-v(\mathfrak a)\in[-\infty,0]
\]
of \eqref{eq:ch08-ideal-potential}.  By
\Cref{prop:ch32-tropical-identities}\,\cref{itm:ch32-ti:identities},
\begin{equation}
\label{eq:ch32-potential-polar-locus}
 \{\varphi_{\mathfrak a}=-\infty\}
 =\{v\in X^{\An}:s(v)\subset Z(\mathfrak a)\}
 =Z(\mathfrak a)^{\An}_{\mathrm{red}},
\end{equation}
the analytification of the reduced zero locus, a closed subset of
$X^{\An}$ contained in the proper-support locus whenever $\mathfrak a$
is nonzero at every generic point of $X$ --- automatic for integral
$X$, but not for the reducible $X$ allowed by
\Cref{sec:ch32-standing}: for $X=V(x_0x_1)\subset\mathbf P^2$ and
$\mathfrak a=\mathcal I_{X_1}$, which is a nonzero coherent ideal because
its stalk at the generic point of $X_2$ is the whole local ring, the
point $v_{X_1,\triv}$ lies in $\{\varphi_{\mathfrak a}=-\infty\}$ and
has full support.  Three consequences deserve emphasis.

\emph{The polar locus is never empty when $\mathfrak a$ is a proper
nonzero ideal.}  Indeed $Z(\mathfrak a)$ then contains a closed point
$p$, and $v_{p,\triv}\in\{\varphi_{\mathfrak a}=-\infty\}$.

\emph{The polar locus meets no scaling orbit in a proper subset.}  It is
scaling stable, because $s(tv)=s(v)$; and it consists of points at which
every homogeneous potential built from $\mathfrak a$ is $-\infty$.

\emph{The polar locus is where a Fubini--Study function degenerates.}  If
\[
 \phi=\max_j\{-v(s_j)+\lambda_j\}
\]
is a Fubini--Study function in the sense of \Cref{def:gf-fs-functions}, built
from finitely many sections
$s_j$ of a line bundle, then $\phi(v)=-\infty$ exactly when every
$s_j$ vanishes along $s(v)$, that is, when $s(v)$ is contained in the
common zero locus.  This is \Cref{ex:gf-proper-support-pole}; the present
chapter supplies the geometric description of the locus, and
\Cref{chap:bj-fs-pl-psh} the analytic consequences.

\begin{example}
\label{ex:ch32-minus-infinity-p2}
\index{potential!polar locus of}
Let $X=\mathbf P^2_k$ with homogeneous coordinates $[x_0:x_1:x_2]$, let
$L=\{x_0=0\}$ be a line with ideal sheaf $\mathfrak a=\cO(-L)$, and let
$p=[0:1:0]\in L$.  Then for $v\in X^{\An}$:
\[
 \varphi_{\mathfrak a}(v)=-v(\mathfrak a)=
 \begin{cases}
 0,&c(v)\notin L,\\
 -v(x_0/x_1),&c(v)\in L\setminus\{x_1=0\},\ s(v)\not\subset L,\\
 -\infty,&s(v)\subset L .
 \end{cases}
\]
Four sample values.
\begin{enumerate}
\item At $v=v_{X,\triv}$: $\varphi_{\mathfrak a}=0$.
\item At $v=t\ord_L$ with $t>0$: $\varphi_{\mathfrak a}(v)=-t$, so
$\varphi_{\mathfrak a}$ is unbounded below already on $X^{\val}$ --- it
is homogeneous and nonzero, hence unbounded on every nontrivial orbit.
\item At $v=\ord_{(C,p)}$ for $C$ a line through $p$ different from
$L$: $C\not\subset L$, so $v(\mathfrak a)=\ord_p\bigl((x_0/x_1)|_C\bigr)
=1$ and $\varphi_{\mathfrak a}(v)=-1$, a finite value at a
proper-support point.
\item At $v=\ord_{(L,p)}$, or at $v=v_{L,\triv}$, or at
$v=v_{p,\triv}$: $s(v)\subset L$, so
$\varphi_{\mathfrak a}(v)=-\infty$.
\end{enumerate}
Items (3) and (4) together show that the value $-\infty$ is not a property
of proper support as such: it is a property of the \emph{position} of the
support relative to the zero locus.  Both possibilities occur at points of
the same kind, and the same potential distinguishes them.  This is the reason
the theory must keep track of supports and cannot simply discard the
non-valuation points.
\end{example}

\subsection{The maximum principle for decreasing functions}
\label{sec:ch32-maximum-principle}

\begin{definition}
\label{def:ch32-decreasing}
\index{function!decreasing}\index{function!generically finite}
A function $\varphi:X^{\An}\to\mathbf R\cup\{-\infty\}$ is
\emph{decreasing} if it reverses the order \eqref{eq:semivaluation-order}:
$v\geq w$ implies $\varphi(v)\leq\varphi(w)$.  It is
\emph{generically finite} if $\varphi|_{X_\alpha^{\An}}\not\equiv-\infty$
for every irreducible component $X_\alpha$ of $X$.
\end{definition}

The homogeneous potentials $\varphi_{\mathfrak a}$ are decreasing, by
\Cref{prop:ch32-tropical-identities}\,\cref{itm:ch32-ti:monotone} and
\eqref{eq:semivaluation-order}; so are all the Fubini--Study and psh
functions of \Cref{chap:bj-fs-pl-psh}, a fact proved there.  The following
statement is \cite[Lemma~1.4]{BJ22}; it is the reason ``supremum'' is so
often replaced by ``value at a trivial semivaluation'' in \Cref{part:bj}.

\begin{proposition}
\label{prop:ch32-maximum-principle}
\index{maximum principle!for decreasing functions}
Let $X$ be a projective $k$-variety and let
$\varphi:X^{\An}\to\mathbf R\cup\{-\infty\}$ be decreasing.
\begin{enumerate}
\item\label{itm:ch32-mp:sup}
For every integral closed $Y\subset X$,
\begin{equation}
\label{eq:ch32-maximum-principle}
 \sup_{Y^{\An}}\varphi=\varphi\bigl(v_{Y,\triv}\bigr).
\end{equation}
\item\label{itm:ch32-mp:gf}
$\varphi$ is generically finite if and only if
$\varphi(v_{X_\alpha,\triv})>-\infty$ for every irreducible component
$X_\alpha$.
\item\label{itm:ch32-mp:usc}
If moreover $\varphi$ is upper semicontinuous, and if $v\in X^{\An}$ has
$Z:=\overline{\{c(v)\}}$, then $t\mapsto\varphi(tv)$ is
non-increasing on $(0,+\infty)$ and
\[
 \lim_{t\to+\infty}\varphi(tv)=\varphi\bigl(v_{Z,\triv}\bigr)
 =\sup_{Z^{\An}}\varphi .
\]
\end{enumerate}
\end{proposition}

\begin{proof}
\cref{itm:ch32-mp:sup}  Let $v\in Y^{\An}$, so $s(v)\subset Y$.  Then
$v\geq v_{s(v),\triv}\geq v_{Y,\triv}$: the first inequality is the left
half of \eqref{eq:trivial-endpoints}, and the second holds because, for a
coherent ideal $\mathfrak b$ with
$v_{Y,\triv}(\mathfrak b)=+\infty$, one has $Y\subset Z(\mathfrak b)$,
hence $s(v)\subset Z(\mathfrak b)$ and
$v_{s(v),\triv}(\mathfrak b)=+\infty$; while
$v_{Y,\triv}(\mathfrak b)=0$ is dominated by anything.  Since $\varphi$
is decreasing, $\varphi(v)\leq\varphi(v_{Y,\triv})$.  As
$v_{Y,\triv}\in Y^{\An}$, the supremum is attained there.

\cref{itm:ch32-mp:gf}  Apply \cref{itm:ch32-mp:sup} with $Y=X_\alpha$:
$\varphi|_{X_\alpha^{\An}}\equiv-\infty$ if and only if its supremum is
$-\infty$, if and only if $\varphi(v_{X_\alpha,\triv})=-\infty$.

\cref{itm:ch32-mp:usc}  For $0<s\leq t$ we have $tv\geq sv$, since
$(tv)(\mathfrak b)=t\,v(\mathfrak b)\geq s\,v(\mathfrak b)=(sv)(\mathfrak b)$
for every $\mathfrak b$; hence $\varphi(tv)\leq\varphi(sv)$ and
$t\mapsto\varphi(tv)$ is non-increasing.  By
\eqref{eq:trivial-endpoints} we have $tv\leq v_{Z,\triv}$ for every
$t$, so $\varphi(tv)\geq\varphi(v_{Z,\triv})$ and
$\lim_{t\to+\infty}\varphi(tv)\geq\varphi(v_{Z,\triv})$.  Conversely
$tv\to v_{Z,\triv}$ by \Cref{prop:ch32-scaling-limits}, so upper
semicontinuity gives
$\limsup_{t\to+\infty}\varphi(tv)\leq\varphi(v_{Z,\triv})$.  The two
inequalities give the limit, and
$\varphi(v_{Z,\triv})=\sup_{Z^{\An}}\varphi$ is
\cref{itm:ch32-mp:sup} applied to $Y=Z$.
\end{proof}

\begin{remark}
\label{rem:ch32-maximum-principle-reading}
\Cref{prop:ch32-maximum-principle} has two faces.  Read forwards, it says
that a decreasing function on $X^{\An}$ is controlled from above by
finitely many numbers, its values at the trivial valuations of the
components; this is what makes the normalization
$\sup_{X^{\An}}\varphi=0$ of \eqref{eq:normalized-psh-space} a normalization
at a point rather than a genuine supremum, and it is used constantly in
compactness arguments.  Read backwards, it says that if
$\varphi(v_{X_\alpha,\triv})=-\infty$ for some component then $\varphi$
is identically $-\infty$ on that whole component's analytification: the
value $-\infty$ at a single trivial valuation is fatal, whereas the value
$-\infty$ at a proper-support point is not.  This asymmetry is exactly the
content of the generic finiteness condition
\Cref{def:ch32-decreasing}, and it is why \Cref{part:bj} keeps the two
notions --- ``$-\infty$ somewhere'' and ``$\equiv-\infty$ on a
component'' --- rigorously apart.
\end{remark}

\subsection{Proper-support points are dense}
\label{sec:ch32-proper-support-dense}

\begin{proposition}
\label{prop:ch32-proper-support-dense}
\index{proper support!density}
Let $X$ be a projective $k$-variety.
\begin{enumerate}
\item\label{itm:ch32-psd:one}
The set of $v\in X^{\An}$ with $\dim s(v)=1$ is dense in $X^{\An}$ if
and only if every irreducible component of $X$ has dimension at least
$1$.
\item\label{itm:ch32-psd:two}
The set $X^{\An}\setminus X^{\val}$ of proper-support points is dense in
$X^{\An}$ if and only if every irreducible component of $X$ has
dimension at least $2$.
\end{enumerate}
\end{proposition}

\begin{proof}
A subset of $X^{\An}$ is dense if and only if its intersection with
$X_\alpha^{\An}$ is dense in $X_\alpha^{\An}$ for every irreducible
component, because $X^{\An}=\bigcup_\alpha X_\alpha^{\An}$ is a finite
closed cover.  So we may assume $X$ integral of dimension $n$.

If $n=0$, then $X^{\An}$ is a point and both displayed subsets are empty,
so neither is dense, as the two equivalences assert.  If $n=1$, the set of
points with one-dimensional
support is $X^{\val}$, which is dense by
\Cref{prop:ch32-xval-not-locally-compact} --- or directly, since each
$v_{p,\triv}$ is the limit of $t\ord_p$ --- while
$X^{\An}\setminus X^{\val}=\{v_{p,\triv}:p\in X(k)\}$ is not dense: its
closure is contained in the proper closed subset $X^{\triv}$ of
\Cref{thm:ch32-trivial-locus}\,\cref{itm:ch32-tl:char}, which omits every
$t\ord_p$ with $0<t<\infty$.

Assume now $n\geq2$.  It suffices to prove that the set $\mathcal C$ of
semivaluations with one-dimensional support is dense, since every such point
lies in $X^{\An}\setminus X^{\val}$; this gives both
\cref{itm:ch32-psd:one} and \cref{itm:ch32-psd:two}.  By
\Cref{thm:triv-divisorial-density} it is enough to show that every
$\ord_E\in X^{\divis}$ lies in the closure of $\mathcal C$; here
$\pi:Y\to X$ is a normal projective birational model and $E\subset Y$ a
prime divisor.  Replacing $Y$ by a resolution of an open neighbourhood of
the generic point of $E$ --- or simply working at a point where $Y$ is
regular, which is possible since $Y$ is normal, hence regular in
codimension one, so regular at the generic point of $E$ and therefore on a
dense open subset of $E$ --- we may assume $Y$ and $E$ both regular
near the point chosen below.

By \Cref{def:trivial-analytification} a basic neighbourhood of $\ord_E$ in
$X^{\An}$ is determined by an affine open $U\subset X$ meeting
$\pi(E)$, finitely many $f_1,\ldots,f_r\in\cO_X(U)$ and
$\varepsilon>0$.  We construct $v\in\mathcal C$ with the exact equalities
$v(f_i)=\ord_E(f_i)$ for all $i$, which is stronger.

Let $Z\subset Y$ be the union of the irreducible components of
$\sum_i\pi^*\operatorname{div}(f_i)$ different from $E$, together with
the singular locus of $Y$ and the singular locus of $E$.  Each of these
meets $E$ in a proper closed subset: the first because a component
different from $E$ cannot contain $E$; the second because $Y$ is
normal, hence regular in codimension one, so regular at the generic point of
$E$; the third because $E$ is integral, hence regular at its generic
point.  Since $\pi^{-1}(U)\cap E$ is a nonempty open subset of $E$, we
may choose a closed point $p\in E\cap\pi^{-1}(U)$ outside $Z$; at $p$
both $Y$ and $E$ are regular.

Embed $Y$ in a projective space by a very ample line bundle and let
$\Lambda$ be a general linear subspace of codimension $n-1$ through
$p$.  By Bertini's theorems on irreducibility and on smooth points ---
valid over any algebraically closed field, see \cite[\S6]{Jou80} --- the
scheme $C:=Y\cap\Lambda$ is, for general $\Lambda$, an integral curve,
regular at $p$, and meeting $E$ transversely at $p$, the last because
a general $\Lambda$ through $p$ is transverse to the tangent space
$T_pE$.  Moreover, for any fixed closed subset $W\subsetneq Y$, cutting
$W$ by $n-1$ general hyperplanes through $p$ leaves a set of dimension
at most $\dim W-(n-1)\leq0$; hence $C\not\subset W$.  Applying this to
$W=$ the exceptional locus of $\pi$, which is a proper closed subset,
gives $C\not\subset\operatorname{Exc}(\pi)$, so $\pi(C)$ is an integral
curve in $X$ meeting $U$.

Let $w:=\ord_{(C,p)}\in Y^{\An}$ be the curve-germ semivaluation of
\Cref{ex:ch32-curve-germ}, and let $v:=\pi^{\An}(w)\in X^{\An}$; then
$s(v)=\pi(C)$ is one-dimensional, so $v\in\mathcal C$, and $s(v)$ meets
$U$, so $v\in U^{\An}$.  Finally, near $p$ the divisor
$\pi^*\operatorname{div}(f_i)$ equals $m_iE$ with
$m_i=\ord_E(f_i)$, because all other components have been excluded by the
choice of $p\notin Z$; writing $z$ for a local equation of $E$ at
$p$, we get $\pi^*f_i=u_iz^{m_i}$ with $u_i$ a unit in
$\cO_{Y,p}$.  Restricting to $C$ and using transversality, which gives
$\ord_p(z|_C)=1$, and $\ord_p(u_i|_C)=0$, we obtain
\[
 v(f_i)=w(\pi^*f_i)=\ord_p\bigl(u_i|_C\bigr)+m_i\ord_p\bigl(z|_C\bigr)
 =m_i=\ord_E(f_i).
\]
Thus $v$ lies in the prescribed neighbourhood of $\ord_E$, and
$\ord_E\in\overline{\mathcal C}$.

Finally, the ``only if'' directions: if some component has dimension $0$,
its analytification is a point of $X^{\triv}$, which has zero-dimensional
support, so the set in \cref{itm:ch32-psd:one} misses a nonempty open subset
of $X^{\An}$; if some component has dimension $1$, the argument of the
case $n=1$ above shows that the proper-support points are not dense in that
component's analytification.
\end{proof}

\begin{warning}{The value $-\infty$ is not an edge case}
Restricting the theory to $X^{\val}$ would lose the compact global domain
used later.  Indeed, $X^{\val}$ is not compact and is not even locally
compact at the trivial valuation
(\Cref{prop:ch32-xval-not-locally-compact}).  Moreover, proper-support points
are dense as soon as $\dim X\geq2$
(\Cref{prop:ch32-proper-support-dense}), so they cannot be discarded
topologically.  Finally, the value of a psh function at a
proper-support point is not determined by its restriction to $X^{\val}$ in
the manner of a continuous extension.  What is true --- that a
$\theta$-psh function is the
smallest usc extension of its restriction to $X^{\divis}$ --- is the
unconditional theorem \cite[Theorem~4.22]{BJ22}, imported in
\Cref{thm:bj-basic-closure}\,\cref{itm:bj-closure:determination}.
\emph{Fourth}, the statement of the
envelope theorem itself, \Cref{thm:bj-final-envelope}, is about functions
that may equal $-\infty$ at proper-support points, and its proof passes
through birational models on which the polar loci move; discarding them would
make the induction impossible.  Accordingly, in this book every psh function
is defined on all of $X^{\An}$, with values in
$\mathbf R\cup\{-\infty\}$, and the locus where it is $-\infty$ is an
object of study, not an exception.
\end{warning}

\section{Gauss extension and the nontrivially valued companion}
\label{sec:ch32-gauss}

Put $K=k((\varpi))$ and normalize $v_K(\varpi)=1$.  If
$v\in X^{\An}$ has full support, its Gauss extension is
\begin{equation}
\label{eq:gauss-extension-triv}
 \sigma(v)\!\left(\sum_i f_i\varpi^i\right)
       :=\min_i\{v(f_i)+i\}.
\end{equation}
The same construction on residue fields treats arbitrary supports.  Formula
\eqref{eq:gauss-extension-triv} gives a continuous section
\begin{equation}
\label{eq:triv-gauss-section}
                   \sigma:X^{\An}\longrightarrow X_K^{\An}
\end{equation}
of the restriction map.  Its image consists of the canonical Gauss points in
the fibres, namely the unique Shilov-boundary points of the fibres of
$X_K^{\An}\to X^{\An}$.  Multiplicativity follows from the ordinary Gauss
lemma, or from the associated graded residue ring.

\subsection{What is already proved, and what is added here}
\label{sec:ch32-gauss-imported}

The construction \eqref{eq:gauss-extension-triv}, its well-definedness, and
its basic properties are established in full in \Cref{sec:ch30-gauss}, in the
notation \eqref{eq:gauss-extension} with the uniformizer called $t$.  We
do not repeat any of it.  For the reader's convenience the five statements
used below are:

\begin{enumerate}
\item\label{itm:ch32-gi:welldef}
\Cref{thm:ch30-gauss-well-defined}: \eqref{eq:gauss-extension-triv} defines a
semivaluation on $A\otimes_kK$ for every finitely generated $k$-algebra
$A$ and every semivaluation $v$ on $A$ trivial on $k^\times$; the
construction is compatible with localization, glues to
\eqref{eq:triv-gauss-section}, satisfies $r\circ\sigma=\mathrm{id}$ for the
restriction map $r:X_K^{\An}\to X^{\An}$, and is a homeomorphism onto a
compact image.
\item\label{itm:ch32-gi:centers}
\Cref{thm:ch30-gauss-centers-supports}: centers are preserved under the
identification $(X_R)_s=X$, see \eqref{eq:ch30-gauss-center}; supports are
base changed, see \eqref{eq:ch30-gauss-support}; divisorial points go to
divisorial points and quasimonomial to quasimonomial.
\item\label{itm:ch32-gi:scaling}
\Cref{prop:ch30-gauss-scaling}: the scaling action corresponds to changing
the normalization of $v_K$, via \eqref{eq:ch30-gauss-scaling}, and to no
action on $X_K^{\An}$ over $K$.
\item\label{itm:ch32-gi:image}
\Cref{thm:ch30-gauss-image}: the image of $\sigma$ is exactly the set of
$\mathbf G_m(k)$-invariant points $y$ with $y(\varpi)=1$, formula
\eqref{eq:ch30-gauss-image}; the fibres of $r$ are large, by
\Cref{ex:ch30-gauss-image-strict}.
\item\label{itm:ch32-gi:dictionary}
\Cref{thm:ch30-gauss-model-dictionary}: the model dictionary, identifying
test configurations with models over $R$, PL functions with model
functions \eqref{eq:ch30-gauss-pl}, vertical intersection numbers
\eqref{eq:ch30-gauss-intersection}, and total masses
\eqref{eq:ch30-gauss-volume}.
\end{enumerate}

What the present section adds is the transfer of the specific structures of
\emph{this} chapter --- ideal evaluations, the order, the trivial locus, the
proper-support locus, and links --- together with the complexification
discussion.

\subsection{Transfer of the structures of this chapter}
\label{sec:ch32-gauss-transfer}

\begin{theorem}
\label{thm:ch32-gauss-transfer}
\index{Gauss extension!transfer of structures}
Let $k$ be trivially valued and algebraically closed, $X$ an integral
projective $k$-variety, $K=k(\!(\varpi)\!)$, $R=k[\![\varpi]\!]$, and
let $\sigma:X^{\An}\to X_K^{\An}$ be the Gauss section
\eqref{eq:triv-gauss-section}, with retraction $r$.
\begin{enumerate}
\item\label{itm:ch32-gt:ideals}
For every coherent ideal $\mathfrak a\subset\cO_X$, writing
$\mathfrak a_K:=\mathfrak a\cdot\cO_{X_K}$,
\begin{equation}
\label{eq:ch32-gauss-ideal}
 \sigma(v)(\mathfrak a_K)=v(\mathfrak a)
 \qquad\text{for all }v\in X^{\An}.
\end{equation}
\item\label{itm:ch32-gt:order}
$\sigma$ is order preserving for the orders
\eqref{eq:semivaluation-order} on $X^{\An}$ and on $X_K^{\An}$ restricted
to base-changed ideals; more precisely, $v\geq w$ in $X^{\An}$ if and
only if $\sigma(v)(\mathfrak a_K)\geq\sigma(w)(\mathfrak a_K)$ for all
coherent $\mathfrak a$ on $X$.
\item\label{itm:ch32-gt:trivial}
$\sigma$ maps $X^{\triv}$ bijectively onto the set of points of
$X_K^{\An}$ of the form $\sigma(v_{Y,\triv})$; these are \emph{not}
trivial semivaluations of $X_K$, since
$\sigma(v_{Y,\triv})(\varpi)=1$.  In particular $\sigma$ does not
commute with the two scaling actions, there being none on the target.
\item\label{itm:ch32-gt:support}
$\sigma$ maps $X^{\val}$ into $X_K^{\val}$ and maps the
proper-support locus into the proper-support locus, with
$\dim s(\sigma(v))=\dim s(v)$; the induced map on supports is base
change, \eqref{eq:ch30-gauss-support}.
\item\label{itm:ch32-gt:link}
If $A:\Omega\to\mathbf R_{>0}$ is a continuous homogeneous gauge on a
scaling-stable $\Omega\subset X^{\An}$, then $\sigma$ carries
$\operatorname{Link}_A(\Omega)$ homeomorphically onto a subset of
$X_K^{\An}$; but the images of the different links
$\operatorname{Link}_A(\Omega)$ and
$\operatorname{Link}_{cA}(\Omega)$, $c>0$, are disjoint unless
$c=1$, and no link is distinguished on the target.
\item\label{itm:ch32-gt:divisorial}
$\sigma(X^{\divis})$ is the set of $\mathbf G_m(k)$-invariant divisorial
points $y$ of $X_K^{\An}$ with $y(\varpi)=1$, and
$\sigma(t\ord_E)$ for $t\in\mathbf Q_{>0}$ is the divisorial point
$b_F^{-1}\ord_F$ of an explicit prime divisor $F$ on a model over $R$,
namely the one produced by \Cref{thm:ch30-gauss-model-dictionary}.
\end{enumerate}
\end{theorem}

\begin{proof}
\cref{itm:ch32-gt:ideals}  Fix an affine chart $U=\Spec A$ with
$c(v)\in U$, and generators $f_1,\ldots,f_s$ of $\mathfrak a(U)$.  Then
$\mathfrak a_K(U_K)$ is generated by the same $f_i$, viewed in
$A\otimes_kK$.  By \eqref{eq:ch30-gauss-center} the center of
$\sigma(v)$ on the trivial model $X_R$ is $c(v)$, so the same elements
generate the stalk of $\mathfrak a_K$ at the center of $\sigma(v)$.
Since $f_i\in A$ has $\sigma(v)(f_i)=v(f_i)$ by
\Cref{thm:ch30-gauss-well-defined}\,\cref{itm:ch30-gw:semival}, taking the
minimum gives \eqref{eq:ch32-gauss-ideal}.

\cref{itm:ch32-gt:order}  Immediate from \cref{itm:ch32-gt:ideals} and
\eqref{eq:semivaluation-order}: both sides of the asserted equivalence are
the same family of inequalities.

\cref{itm:ch32-gt:trivial}  By \eqref{eq:ch30-gauss-support} the support of
$\sigma(v_{Y,\triv})$ is the generic point of $Y_K$; but
$\sigma(v_{Y,\triv})(\varpi)=1\neq0$ by
\Cref{thm:ch30-gauss-well-defined}\,\cref{itm:ch30-gw:semival}, so
$\sigma(v_{Y,\triv})$ is not a trivial semivaluation of $X_K$.  It is in
fact the Gauss point of the trivial model of $Y_K$, as recorded in
\Cref{rem:ch30-gauss-ray-endpoints}.  Injectivity is
\eqref{eq:gauss-section}.

\cref{itm:ch32-gt:support}  By \eqref{eq:ch30-gauss-support} the support of
$\sigma(v)$ is the base change of the support of $v$; base change along
the field extension $K/k$ preserves dimension and irreducibility, the
latter because $K/k$ is a separable, indeed regular, extension.  Hence
$\dim s(\sigma(v))=\dim s(v)$, and $s(\sigma(v))$ is an irreducible
component of $X_K$ exactly when $s(v)$ is one of $X$.

\cref{itm:ch32-gt:link}  $\sigma$ is a homeomorphism onto its image by
\Cref{thm:ch30-gauss-well-defined}\,\cref{itm:ch30-gw:cont}, so it carries
any subspace homeomorphically onto its image.  The disjointness assertion
follows because $\operatorname{Link}_{cA}(\Omega)
=c^{-1}\operatorname{Link}_A(\Omega)$ and $\sigma$ is injective while
$\sigma(c^{-1}v)\neq\sigma(v)$ for $c\neq1$ and
$v\notin X^{\triv}$.  That no link on the target is distinguished is the
content of \Cref{prop:ch30-gauss-scaling}: rescaling on the source
corresponds to changing $v_K(\varpi)$ and then rescaling the resulting
valuation.  Hence a constant rescaling of a link does not arise from an
action on the fixed normalized target $X_K^{\An}$; the source gauge is
additional data.

\cref{itm:ch32-gt:divisorial}  Combine
\Cref{thm:ch30-gauss-image} with
\Cref{thm:ch30-gauss-centers-supports}\,\cref{itm:ch30-gcs:divisorial} and
\Cref{thm:ch30-gauss-model-dictionary}\,\cref{itm:ch30-gmd:points}: the
latter exhibits, for a divisorial $v\in X^{\divis}$, an integrally closed
test configuration $\cX$ with a component $E$ of $\cX_0$ such that
$v=v_E$, and identifies $\sigma(v_E)$ with $b_E^{-1}\ord_E$ on the
model $\cX_R$.
\end{proof}

\begin{remark}
\label{rem:ch32-gauss-not-transfer}
Three structures of this chapter have no counterpart on the target.
\emph{The scaling action}: $X_K^{\An}$ carries no $\mathbf R_{>0}$-action
over $K$, because every point satisfies $y(\varpi)=1$; see
\Cref{prop:ch30-gauss-scaling} and the warning after it.
\emph{The trivial locus}: there is no analogue of $X^{\triv}$ in
$X_K^{\An}$, since every analytic point restricts to the prescribed
nontrivial valuation of $K$, and hence has $y(\varpi)=1$.
\emph{Compactness of the fibres of $c$}: the center map of $X_K^{\An}$
refers to a model over $R$, not to $X_K$, and the sets $c^{-1}(U)$ of
\Cref{thm:ch32-center-compact} are not the preimages under $r$ of anything.
The correct bookkeeping is the four-data ledger of
\Cref{cor:ch30-gauss-four-data}: centers, supports, classes, and total mass.
\end{remark}

\subsection{The test-configuration realization}
\label{sec:ch32-gauss-tc}

If $\cX\to\mathbf A^1_k$ is a test configuration, then
$\cX\setminus\cX_0\simeq X\times\mathbf G_m$.  The Gauss extension embeds
$X^{\An}$ into $\cX^{\An}$, centered on the central fibre and normalized
by $v(\varpi)=1$.  In this test-configuration realization, its image is
exactly the set of semivaluations invariant under the
$\mathbf G_m(k)$-action on $\cX$ and satisfying $v(\varpi)=1$;
invariance here refers to the action on the total space, not merely to
triviality on scalar constants.  Vertical divisors on $\cX$ can therefore
be evaluated on $X^{\An}$.  This is the bridge used in the next chapter.

Three points deserve to be made precise, since they are the ones on which
\Cref{chap:bj-test-configurations} and \Cref{chap:bj-fs-pl-psh} rest.

\emph{The realization is independent of $\cX$.}  A test configuration
$\cX$ is a compactification of $X\times\mathbf G_m$ over
$\mathbf A^1$, and the embedding
$\sigma:X^{\An}\hookrightarrow(\cX\setminus\cX_0)^{\An}
\subset\overline{\cX}^{\An}$ depends only on the identification
\eqref{eq:tc-punctured-trivialization}, not on $\cX$ itself.  This is
\cite[Remark~1.9]{BJ22}, and in the present book it is the statement that
$\sigma$ is defined on $X^{\An}$ once and for all by
\eqref{eq:triv-gauss-section}, the test configuration entering only through
the model $\cX_R$ of \Cref{thm:ch30-gauss-model-dictionary}.

\emph{Vertical divisors become PL functions.}  A vertical
$\mathbf Q$-Cartier divisor $D$ on $\cX$, that is, a
$\mathbf G_m$-invariant $\mathbf Q$-Cartier divisor supported on
$\cX_0$, defines $\phi_D\in\PL(X)$ by $\phi_D(v)=\sigma(v)(D)$, with
the sign convention of \eqref{eq:appJ-model-divisor}; this is
\Cref{thm:tc-pl-correspondence} and, on the model side,
\eqref{eq:ch30-gauss-pl}.  Effective $D$ give nonnegative
$\phi_D$, by the sign test of \eqref{eq:appJ-special-fibre-test}.

\emph{Scaling changes the ground normalization.}  For $t>0$, let
$\sigma_\lambda$ denote Gauss extension formed with
$v_K(\varpi)=\lambda$.  Then the exact relation is
\[
                    \sigma(tv)=t\,\sigma_{1/t}(v),
\]
which is \eqref{eq:ch30-gauss-scaling}.  Thus source scaling is represented
by changing the ground normalization and then rescaling; it is not an action
on the fixed slice $v_K(\varpi)=1$.  For rational $t$, the corresponding
model operation is obtained by finite base change followed by normalization,
as in \Cref{prop:tc-valuative-invariance}.

\begin{example}
\label{ex:ch32-gauss-ray-numbers}
Let $X=\mathbf P^1_k$ with coordinate $T$, $p=0$, and consider the
divisorial ray $v_t:=t\ord_0$, $t\in(0,+\infty)$.  Its Gauss extension
is the monomial point of $(\mathbf P^1_K)^{\An}$ with
\[
 \sigma(v_t)(T)=t,\qquad \sigma(v_t)(\varpi)=1,
\]
that is, the Gauss point of the disc of radius $|\varpi|^t$; see
\Cref{ex:ch08-gauss-ray}, where the same computation is carried out in the
tropical coordinate of \Cref{sec:ch08-p1-tree}.  The two endpoints behave as
follows.  At $t\to0^+$ the source point tends to $v_{X,\triv}$ and the
image tends to $\sigma(v_{X,\triv})$, the Gauss point of the unit disc,
which is the divisorial point of the special fibre of the trivial model
$\mathbf P^1_R$.  At $t\to+\infty$ the source point tends to
$v_{0,\triv}$, of proper support, and the image tends to
$\sigma(v_{0,\triv})$, the rigid point $T=0$ of
$(\mathbf P^1_K)^{\An}$, which carries no Monge--Amp\`ere mass of a bounded
potential, the rigid point being pluripolar
(\Cref{ex:ch23-sup-not-usc}, \Cref{prop:bounded-ma-no-polar-mass}).  Two numerical checks:
$\sigma(v_t)(T^m)=mt$ for every $m$, so the image ray is parametrized
linearly; and $\sigma(2v_t)=\sigma(v_{2t})$ has
$\sigma(2v_t)(T)=2t$ while $\sigma(2v_t)(\varpi)=1$, confirming that
scaling on the source does not act on the target but moves along the
apartment.
\end{example}

\subsection{Complexification}
\label{sec:ch32-complexification}

When $k=\mathbf C$, one may also regard $\cX\to\mathbf A^1_\mathbf C$ as
a complex degeneration over a punctured disc.  Formal completion and base
change to $\mathbf C((\varpi))$ produce its nontrivially valued analytic
generic fibre, while \eqref{eq:triv-gauss-section} selects the invariant slice
corresponding to the trivially valued analytification.  This is the precise
sense in which a complex degeneration gives a complexification of the
radial valuation picture.  There is no canonical map
$X(\mathbf C)\to X^{\An}$ carrying the Euclidean topology to the Berkovich
topology; the comparison runs through the degeneration.

We make the negative statement precise, since it is often assumed away.

\begin{proposition}
\label{prop:ch32-no-complexification-map}
\index{complexification!absence of a canonical map}
Let $X$ be an integral projective $\mathbf C$-variety with
$\dim X\geq1$, and give $\mathbf C$ the trivial absolute value.
\begin{enumerate}
\item\label{itm:ch32-nc:points}
The assignment $p\mapsto v_{p,\triv}$ is an injection
$X(\mathbf C)\hookrightarrow X^{\An}$, and its image is a set of isolated
points of $X^{\triv}$ in the sense of
\Cref{ex:ch32-trivial-locus-curve}; the induced topology on
$X(\mathbf C)$ is the discrete topology, not the Euclidean one.
\item\label{itm:ch32-nc:no-cont}
The assignment $p\mapsto v_{p,\triv}$ is not continuous on any connected
positive-dimensional component of $X(\mathbf C)^{\mathrm{eucl}}$.  More
generally, every continuous map from $X(\mathbf C)^{\mathrm{eucl}}$ whose
image is contained in $\{v_{p,\triv}:p\in X(\mathbf C)\}$ is locally
constant.
\end{enumerate}
\end{proposition}

\begin{proof}
\cref{itm:ch32-nc:points}  Injectivity is
\Cref{thm:ch32-trivial-locus}\,\cref{itm:ch32-tl:bijection}.  For each
$p\in X(\mathbf C)$ the set $\{v:v(\mathfrak m_p)>1\}$ is open in
$X^{\An}$ and meets $\{v_{q,\triv}:q\in X(\mathbf C)\}$ only in
$v_{p,\triv}$, because $v_{q,\triv}(\mathfrak m_p)=0$ for $q\neq p$
and $=+\infty$ for $q=p$.  Hence the subspace topology on the image is
discrete.

\cref{itm:ch32-nc:no-cont}  The displayed target subset is discrete by
\cref{itm:ch32-nc:points}; therefore any continuous map into it is locally
constant.  On a connected positive-dimensional component the assignment is
injective and nonconstant, so it cannot be continuous.
\end{proof}

\begin{remark}
\label{rem:ch32-complexification-content}
\Cref{prop:ch32-no-complexification-map} does not say that there is no
relation between the complex and the trivially valued pictures; it says the
relation is not a map of spaces defined pointwise.  Three genuine
comparisons exist, and each is developed elsewhere in this book.
\begin{enumerate}
\item\emph{Through a degeneration.}  A test configuration $(\cX,\cL)$ for
$(X,L)$ over $\mathbf C$ has, on the one hand, a central fibre and a
$\mathbf G_m$-action, which produce a PL function on $X^{\An}$ by
\Cref{thm:tc-pl-correspondence}; and, on the other hand, a family of
complex manifolds over the punctured disc, which produces a subgeodesic ray
of metrics on $L$.  The dictionary between these two is the subject of
\Cref{chap:rays-filtrations}.
\item\emph{Through filtrations.}  A filtration of the section ring of
$(X,L)$ gives, on the non-Archimedean side, a Fubini--Study function via
\eqref{eq:norm-fs-function}, and on the complex side a family of
Bergman-type metrics; the two are compared in
\Cref{chap:dxz-bj-comparison}.
\item\emph{Through hybrid spaces.}  A one-parameter family of complex
varieties degenerating at $\varpi=0$ can be compactified by attaching
$X^{\An}$ as a boundary, so that Euclidean and Berkovich topologies live on
one space; this is the hybrid construction of
\Cref{chap:hybrid-spaces}.
\end{enumerate}
None of these is a pointwise map $X(\mathbf C)\to X^{\An}$, and each carries
its own hypotheses.  A statement of the form ``the trivially valued
analytification is the complexification of the valuation picture'' is
meaningful only relative to one of these three constructions, and the
present chapter fixes only the target.
\end{remark}

\begin{remark}
\label{rem:ch32-larger-trivial-field}
There is a fourth, quite different, change of base worth separating from the
Gauss extension: if $k'/k$ is an extension of algebraically closed fields,
both trivially valued, then \Cref{lem:ch30-base-change-map} provides a
continuous surjection $X_{k'}^{\An}\to X^{\An}$ of compact spaces, and the
structures of this chapter --- supports, centers, ideal evaluations, the
order and the scaling action --- are compatible with it.  The supply of
divisorial points may increase, since a prime divisor on a model of
$X_{k'}$ need not descend to $k$; no strictness is asserted for every
extension or every $X$.  We do not use this base change in
\Cref{part:bj}, and record it only to prevent it from being confused with
\eqref{eq:triv-gauss-section}, whose target is a \emph{nontrivially} valued
analytification and whose source and target are analytifications of schemes
over incomparable valued fields, as emphasized in
\Cref{sec:ch30-gauss-setting}.
\end{remark}

\section{Dictionary and provenance}
\label{sec:ch32-dictionary}

\subsection{The algebra--analysis dictionary}
\label{sec:ch32-dict-forward}

The table records each correspondence together with the place where it is
proved.

\begin{center}
\small
\begin{tabular}{p{.30\textwidth}|p{.34\textwidth}|p{.26\textwidth}}
\toprule
\textbf{algebraic side}&\textbf{analytic side}&\textbf{proved in}\\
\midrule
coherent ideal $\mathfrak a$
 &continuous $v\mapsto v(\mathfrak a)\in[0,+\infty]$
 &\Cref{thm:ch32-ideals-generate}\\
product, sum of ideals
 &sum, minimum of values
 &\Cref{prop:ch32-tropical-identities}\\
scheme point $\xi\in X$
 &trivial semivaluation $v_{\overline{\{\xi\}},\triv}$
 &\Cref{thm:ch32-trivial-locus}\\
specialization $\xi\rightsquigarrow\xi'$
 &order $v\geq w$ on centers
 &\Cref{thm:ch32-order-centers}\\
affine open $U=\Spec A$
 &compact $c^{-1}(U)=\mathcal M(A)$
 &\Cref{thm:ch32-center-compact}\\
prime divisor $E$ on a model
 &divisorial ray $\mathbf R_{>0}\ord_E$
 &\Cref{prop:ch32-divisorial-invariants}\\
component of an integrally closed $\cX_0$
 &divisorial point $v_E$
 &\Cref{thm:ch32-divisorial-tc}\\
SNC stratum with weights $\alpha$
 &quasimonomial point $v_\alpha$
 &\Cref{prop:ch32-qm-well-defined}\\
relative canonical divisor
 &log discrepancy, affine on cones
 &\Cref{prop:ch32-log-discrepancy-affine}\\
integral curve germ $(C,p)$
 &proper-support point $\ord_{(C,p)}$
 &\Cref{ex:ch32-curve-germ}\\
zero locus $Z(\mathfrak a)$
 &polar locus $\{\varphi_{\mathfrak a}=-\infty\}$
 &\eqref{eq:ch32-potential-polar-locus}\\
test configuration $\cX$
 &model $\cX_R$ over $R$ via $\sigma$
 &\Cref{thm:ch32-gauss-transfer}\\
\bottomrule
\end{tabular}
\end{center}

\subsection{Numerical checkpoints}
\label{sec:ch32-dict-numbers}

The following identities recur and are the fastest way to detect a lost
factor.  All are proved above.
\begin{enumerate}
\item $A_{\mathbf A^2}(v_w)=w_1+w_2$ for the monomial valuation of weights
$w$; in particular $A(\ord_0)=2$, not $1$
(\Cref{ex:ch32-affine-plane}).
\item On the tower of two blowups of \Cref{ex:ch32-blowup-surface},
$A(\ord_{E_1})=2$, $A(\ord_{E_2})=3$, and
$A(v_\alpha)=2\alpha_1+3\alpha_2$.
\item The quasimonomial point $\alpha=(1,1)$ of that example is the
monomial valuation of weights $(2,3)$, and both computations of its log
discrepancy give $5$.
\item $\sigma(t\ord_p)(T)=t$ and $\sigma(t\ord_p)(\varpi)=1$ on
$\mathbf P^1$; a divisorial point of the trivially valued side with
parameter $t$ becomes a point of the apartment at distance $t$ from the
Gauss point (\Cref{ex:ch32-gauss-ray-numbers}).
\item A component $E$ of a central fibre gives $v_E=b_E^{-1}\ord_E|_{k(X)}$
with $b_E=\ord_E(\varpi)$: the factor $b_E$ lives on the model side and
vanishes on the trivially valued side (\Cref{rem:ch32-bE}).
\end{enumerate}

\subsection{Source note}
\label{sec:ch32-source-note}

The semivaluation, topology, scaling, curve, and Gauss-extension statements
are from \cite[\S1]{BJ22}, building on Berkovich geometry and
\cite{BHJ17}.  Two safeguards will be used throughout
\Crefrange{chap:trivial-analytification}{chap:bj-stability}: values at
proper-support points are retained, and compactness is formulated with nets.
The density theorem, including the PL Stone--Weierstrass and vertical-divisor
steps reproduced above, is \cite[Theorems~2.2, 2.7, and~2.14 and
Lemma~2.8]{BJ22}.

We record the provenance of the individual items in the order in which they
appear.  The identification of the points of a trivially valued
analytification with semivaluations, and the partition by support, is
Berkovich's construction \cite{Berk12} specialized to the trivial absolute
value; the presentation as a disjoint union of spaces of valuations of
subvarieties is \cite[\S1.1.2]{BJ22}.  The evaluation of a semivaluation on
an ideal, the tropical identities \eqref{eq:ideal-tropical-identities}, and
the tropical-spectrum description are \cite[\S1.1.3]{BJ22}, whose local
antecedent is \cite[\S1.2]{JM12}; the affine form proved here as
\Cref{thm:ch32-tropical-spectrum} is a complete statement of what that
description asserts, and its proof is elementary.  The proof of compactness
by centers, \Cref{thm:ch32-center-compact}, is the standard reduction of the
properness theorem to the compactness of a Banach spectrum
\cite[Theorem~1.2.1]{Berk12}; we have written it out because it also produces
the compact sets $c^{-1}(U)$ and the anticontinuity of the center map,
which in Berkovich's language is the assertion that $c^{-1}(U)$ is a
$k$-analytic domain.  Anticontinuity of the center map and continuity of
the support map are \cite[\S1.1.3]{BJ22}.

The countability statements of \Cref{sec:ch32-countability} are the two
halves of \cite[Remark~1.1]{BJ22}.  We have proved metrizability over a
countable field by exhibiting a countable generating family of continuous
functions, and the failure of first countability over an uncountable field by
the Baire-type argument of \Cref{lem:ch32-uncountable-field}, which is
classical.  The sequential compactness statement quoted in the warning after
\Cref{thm:ch32-not-first-countable} rests on a theorem of Poineau on
Berkovich spaces; we neither prove nor use it, and we have therefore given no
citation key for it beyond the reference in \cite[Remark~1.1]{BJ22}.

The partial order \eqref{eq:semivaluation-order} and its description by
centers, \Cref{thm:ch32-order-centers}, is \cite[Lemma~4.4]{JM12}
transported to the global setting, as in \cite[\S1.1.4]{BJ22}; the
sandwiching \eqref{eq:trivial-endpoints} is \cite[Example~1.2]{BJ22}.  The
identification of the fixed-point set of the scaling action with the Zariski
space, \Cref{thm:ch32-trivial-locus}, is \cite[\S1.1.5]{BJ22}.  The maximum
principle \Cref{prop:ch32-maximum-principle} is \cite[Lemma~1.4]{BJ22}.  The
factor $t$ in the induced action on functions
\eqref{eq:ch32-scaling-on-functions} is the convention of
\cite[(1.3)]{BJ22}, whose reason is the geometric interpretation via base
change; compare \Cref{prop:ch30-gauss-scaling} here.

The description of the analytification of a curve, \Cref{sec:ch32-line}, is
\cite[\S1.1.6]{BJ22}; the classification of the points of the trivially
valued affine line given as \Cref{thm:ch32-a1-classification} is folklore and
is proved here in full.  The failure of local compactness of $X^{\val}$,
\Cref{prop:ch32-xval-not-locally-compact}, is stated in the same place; the
proof by a net along a ray is the natural one and is written out here.

Divisorial valuations over a trivially valued field, including the convention
that $v_{X,\triv}$ counts as one, are \cite[\S1.1.1]{BJ22}; their relation
to integrally closed test configurations, \Cref{thm:ch32-divisorial-tc}, is
\cite[Lemma~1.20]{BJ22}, resting on \cite[Theorem~4.6]{BHJ17} and ultimately
on Zariski's realization theorem, for which our reference is
\cite[Ch.~VI]{ZariskiSamuelII}.  Linear growth of divisorial valuations,
\Cref{prop:ch32-divisorial-linear-growth}, is \cite[Proposition~1.28]{BJ22};
the intersection-theoretic proof is the one given there.  The density theorem
\Cref{thm:triv-divisorial-density} is \cite[Theorem~2.14]{BJ22}, and the
density of proper-support points \Cref{prop:ch32-proper-support-dense} is
\cite[Lemma~2.15]{BJ22}; the curve-germ construction in its proof is due to
the same source, and the Bertini input is \cite[\S6]{Jou80}.

Quasimonomial valuations over a trivially valued base are classical; the
construction by graded reduction that we use is the one of
\Cref{lem:monomial-valuation}, and its invariants
\Cref{thm:ch32-quasimonomial-invariants} are the trivially valued form of
\Cref{lem:skeleton-monomial-invariants}.  The equivalence between
quasimonomial and Abhyankar in characteristic zero is
\cite[\S3]{JM12}, and the realization of $X^{\An}$ as an inverse limit of
dual complexes for smooth $X$ in characteristic zero is
\cite[Appendix~A]{BJ22}; neither is used in any proof of this book.  The
affine-plane computations of \Cref{ex:ch32-affine-plane} reproduce, with the
invariants displayed, the monomial gallery of
\Cref{lem:monomial-invariants} and \Cref{ex:gallery-divisorial}; the
surface computation of \Cref{ex:ch32-blowup-surface} is elementary and is
included because the same tower recurs in \Cref{chap:bj-envelope-statement}
and in \Cref{ex:ch30-surface-blowup}.

Everything about the Gauss extension is imported from
\Cref{sec:ch30-gauss}, where it is proved in full; the underlying statements
are \cite[\S1.3]{BJ22} and \cite[Lemma~4.2]{BHJ17}, and the identification of
the image with the $\mathbf G_m(k)$-invariant slice is
\cite[Remark~1.9]{BJ22}.  The transfer statement
\Cref{thm:ch32-gauss-transfer} and the negative complexification statement
\Cref{prop:ch32-no-complexification-map} are formulated here for the purposes
of \Cref{part:comparisons} and are proved in full.  The three genuine
complex--non-Archimedean comparisons listed in
\Cref{rem:ch32-complexification-content} are developed in
\Cref{chap:rays-filtrations}, \Cref{chap:dxz-bj-comparison} and
\Cref{chap:hybrid-spaces} respectively.

Finally, a remark on what this chapter does \emph{not} contain.  There is no
Monge--Amp\`ere operator, no energy, and no psh function here; the classes
$\PL(X)$, $\FS(X)$ and $\PSH(\theta)$ are introduced in
\Cref{chap:bj-fs-pl-psh}, after the test configurations of
\Cref{chap:bj-test-configurations}.  The only forward references made above
are to \Cref{thm:tc-pl-correspondence} and
\eqref{eq:pl-stone-weierstrass}, both used in the proof of
\Cref{thm:triv-divisorial-density} and both independent of it, as
\Cref{rem:ch32-density-ingredients} records.
